\documentclass[11pt,a4paper,oneside]{book}
\usepackage[T1]{fontenc}
\usepackage[utf8]{inputenc}
\usepackage[indonesian]{babel}
\usepackage{lmodern}
\usepackage{amsmath,amsthm,amssymb,amscd}
\usepackage{graphicx}
\usepackage[a4paper,margin=22mm,headheight=15pt]{geometry}
\usepackage{microtype}
\usepackage[hidelinks,pdfusetitle]{hyperref}
\usepackage{bookmark}
% language package supplied by unit wrapper

% UTF-8 input supplied by unit wrapper

%\usepackage[utf8]{inputenc}

\usepackage{ifthen}

%\usepackage{fancyhdr}

%Aus Grundvorspann

%Packages

\usepackage{amscd}

\usepackage{amssymb}

\usepackage{epsfig} %Graphik

\usepackage{verbatim,moreverb} %

\usepackage{enumitem}

\usepackage{eurosym}

\usepackage{epigraph}

\setcounter{MaxMatrixCols}{20}

%Mathematische Einzelbefehle

%[[Projekt:Semantische Vorlagen/Mathematischer Modus/Latex]]

\newcommand{\mathdisp}[2]{\[ #1 \, #2 \]}

\newcommand{\mathdisplayteile}[3]{\[ #1 \] \[ #2 \, #3 \]}

\newcommand{\mathind}[3]{$ #1$, $#2 $#3}

\newcommand{\mathbed}[8]{$ #1 ${\ifthenelse{\equal{#2}{}}{, }{ #2 }}$ #3 ${\ifthenelse{\equal{#5}{}}{}{{\ifthenelse{\equal{#4}{}}{, }{ #4 }}}$ #5 $}{\ifthenelse{\equal{#7}{}}{}{{\ifthenelse{\equal{#6}{}}{, }{ #6 }}}$ #7 $}#8}

\newcommand{\mathbeddisp}[8]{\[ #1 {\ifthenelse{\equal{#2}{}}{,\, }{ \text{ #2 } }} #3 {\ifthenelse{\equal{#5}{}}{}{{\ifthenelse{\equal{#4}{}}{,\, }{ \text{ #4 } }}} #5 }{\ifthenelse{\equal{#7}{}}{}{{\ifthenelse{\equal{#6}{}}{,\, }{ \text{ #6 } }}} #7 }#8 \]}

\newcommand{\mathbedtermdisp}[8]{\[ #1\] #2 $ #3 $#4 $ #5 $#6$ #7 $#8}

\newcommand{\mathkon}[4]{$ #1 $ #2 $ #3 $#4}

\newcommand{\mathkor}[5]{{\ifthenelse{\equal{#1}{}}{}{#1 }}$ #2 $ #3 $ #4 $#5}

\newcommand{\mathkork}[5]{{\ifthenelse{\equal{#1}{}}{}{#1 }}$ #2 $ #3 $ #4 $#5}

\newcommand{\mathlist}[5]{$ #1 $\ifthenelse{\equal{#2}{}}{,}{ #2} $ #3 $\ifthenelse{\equal{#4}{}}{,}{ #4} $ #5 $}

\newcommand{\mathlistdisp}[5]{\[  #1  \ifthenelse{\equal{#2}{}}{,}{ \text{ #2 }}  #3  \ifthenelse{\equal{#4}{}}{,}{\text{ #4 } }  #5 \] }

\newcommand{\alignier}[2]{\begin{align*} #1 \, #2 \end{align*}}

\newcommand{\mathlistvierdisp}[7]{\[ #1,\, #3, \, #5 \text{ #6 } #7  \]}

%[[Projekt:Semantische Vorlagen/Mathematische Vergleichsvorlagen/Latex]]

\newcommand{\mavergleichskette}[4]{$ #1 #2 #3#4$}

\newcommand{\mavergleichskettedisp}[4]{\[  #1 #2 #3 #4 \]}

\newcommand{\mavergleichskettedisplang}[4]{\[  #1 #2 #3 #4 \]}

\newcommand{\mavergleichskettealign}[4]{ \begin{eqnarray*}  #1 #2 #3 #4 \end{eqnarray*}  }

\newcommand{\vergleichskette}[9]{ #1 \, #2 \, #3 \ifthenelse
{\equal {#4}{ }} {} {\, #4 \, #5 }  \ifthenelse {\equal {#6}{ } }{}
{\, #6 \, #7 } \ifthenelse {\equal {#8}{ }}{} {\, #8 \, #9 }  }

\newcommand{\vergleichskettelang}[9]{ #1 \, #2 \, #3 \ifthenelse
{\equal {#4}{ }} {} {\, #4 \, #5 }  \ifthenelse {\equal {#6}{ } }{}
{\, #6 \, #7 } \ifthenelse {\equal {#8}{ }}{} {\, #8 \, #9 }  }

\newcommand{\vergleichskettefortsetzung}[8]{ \, #1 \, #2 \ifthenelse
{\equal {#3}{ }} {} {\, #3 \, #4 }  \ifthenelse {\equal {#5}{ } }{}
{\, #5 \, #6 } \ifthenelse {\equal {#7}{ }}{} {\, #7 \, #8 } }

\newcommand{\vergleichskettealign}[9]{ #1 & #2 & #3 \ifthenelse
{\equal {#4}{ }} {} {\cr & #4 & #5 }  \ifthenelse {\equal {#6}{ }
}{} {\cr & #6 & #7 } \ifthenelse {\equal {#8}{ }}{} {\cr & #8 & #9 }
}

\newcommand{\vergleichskettefortsetzungalign}[8]{ \cr & #1 & #2 \ifthenelse
{\equal {#3}{ }} {} {\cr & #3 & #4 }  \ifthenelse {\equal {#5}{ }
}{} {\cr & #5 & #6 } \ifthenelse {\equal {#7}{ }}{} {\cr & #7 & #8 }
}

\newcommand{\mavergleichskettealigndrucklinks}[4]{ \begin{eqnarray*}  #1 #2 #3 #4 \end{eqnarray*}  }

\newcommand{\vergleichskettealigndrucklinks}[9]{  &  & #1 \cr & #2 & #3 \ifthenelse
{\equal {#4}{ }} {} {\cr & #4 & #5 }  \ifthenelse {\equal {#6}{ }
}{} {\cr & #6 & #7 } \ifthenelse {\equal {#8}{ }}{} {\cr & #8 & #9 }
}

\newcommand{\mavergleichskettealignhandlinks}[4]{ \begin{eqnarray*}  #1 #2 #3 #4 \end{eqnarray*}  }

\newcommand{\vergleichskettealignhandlinks}[9]{ #1 & #2 & #3 \ifthenelse
{\equal {#4}{ }} {} {\cr & #4 & #5 }  \ifthenelse {\equal {#6}{ }
}{} {\cr & #6 & #7 } \ifthenelse {\equal {#8}{ }}{} {\cr & #8 & #9 }
}

\newcommand{\mavergleichskettedisphandlinks}[4]{ \[ #1 #2 #3 #4 \] }

\newcommand{\vergleichskettedisphandlinks}[9]{  #1 \, #2 \, #3 \ifthenelse
{\equal {#4}{ }} {} {\, #4 \, #5 }  \ifthenelse {\equal {#6}{ } }{}
{\, #6 \, #7 } \ifthenelse {\equal {#8}{ }}{} {\, #8 \, #9 }  }

\newcommand{\zeilemitteil}[2]{ \ifthenelse {\equal {#2} {} } {#1} { #1 \cr & & #2}   }

\newcommand{\mavergleichskettek}[4]{$ #1 #2 #3#4$}

\newcommand{\vergleichskettek}[9]{ #1 \, #2 \, #3 \ifthenelse
{\equal {#4}{ }} {} {\, #4 \, #5 }  \ifthenelse {\equal {#6}{ } }{}
{\, #6 \, #7 } \ifthenelse {\equal {#8}{ }}{} {\, #8 \, #9 }  }

\newcommand{\vergleichskettefortsetzungk}[8]{ \, #1 \, #2 \ifthenelse
{\equal {#3}{ }} {} {\, #3 \, #4 }  \ifthenelse {\equal {#5}{ } }{}
{\, #5 \, #6 } \ifthenelse {\equal {#7}{ }}{} {\, #7 \, #8 } }

%[[Projekt:Semantische Vorlagen/Mathematische Strukturvorlagen/Latex]]

\newcommand{\N}   {{\mathbb N}}
\newcommand{\Z}   {{\mathbb Z}}
\newcommand{\Q}   {{\mathbb Q}}
\newcommand{\R}   {{\mathbb R}}
\newcommand{\C}   {{\mathbb C}}
\newcommand{\Complex}  {{\mathbb C}}

\newcommand{\defeq}{:=}

\newcommand{\defeqr}{=:}

\newcommand{\maabb}[4]{${\ifthenelse {\equal {#1}{} } {} {#1 \colon }}\,  #2 \rightarrow #3 #4 $}

\newcommand{\maabbele}[6]{$ {\ifthenelse {\equal {#1}{} } {} {#1 \colon }}\,   #2 \rightarrow  #3 , \, #4 \mapsto #5 #6 $}

\newcommand{\maabbdisp}[4]{\[ {\ifthenelse {\equal {#1}{} } {} {#1 \colon }}\,  #2 \longrightarrow #3 #4 \]}

\newcommand{\maabbeledisp}[6]{\[ {\ifthenelse {\equal {#1}{} } {} {#1 \colon }}\,   #2 \longrightarrow  #3 , \, #4 \longmapsto #5 #6 \]}

%Variante bei ifequal-Problem

\newcommand{\maabbnamedisp}[4]{\[ #1 \colon \, #2 \longrightarrow #3 #4 \]}

\newcommand{\maabbeledispvar}[6]{\[ {\ifthenelse {\equal {#1}{} } {} {#1 \colon }}\,   #2 \longrightarrow  #3 , \, #4 \longmapsto #5 #6  \]}

\newcommand{\maabbelementzeiledisplay}[6]{\[ {\ifthenelse {\equal {#1}{} } {} {#1 \colon }}\, #2 \longrightarrow #3 , \]  	\[ #4 \longmapsto #5 #6 \]}

\newcommand{\maabbelementdoppelzeiledisplay}[6]{\[ {\ifthenelse {\equal {#1}{} } {} {#1 \colon }}\, #2 \longrightarrow #3 , \] \[ #4 \longmapsto \] \[  #5 #6 \]}

\newcommand{ \arccot  } {\operatorname{arccot } }

\newcommand{\betrag}[1]{\left| #1 \right|}

\newcommand{\strukturgarbe}{ {\mathcal O} }

\newcommand{\schnittring}[1]{\Gamma (#1 , {\mathcal O})}

\newcommand{\tensor}{\otimes}

%[[Projekt:Semantische Vorlagen/Textgestaltung/Latex]]

\newcommand{\einrueckung}[1]{\par \hspace{1cm}#1 \par}

\newcommand{\zusatz}[2]{(#1)#2}

\newcommand{\zusatzklammer}[3]{(#1#2)#3}

\newcommand{\zusatzgs}[3]{\-- #1 \-- #3}

\newcommand{\zusatzfussnote}[3]{#3\footnote{#1#2}}

%Feinabstimmung

\newcommand{\latexdruckeinskurz}{\!}

\newcommand{\latexdruckzweikurz}{\! \!}

\newcommand{\latexdruckdreikurz}{\! \! \!}

%[[Projekt:Semantische Vorlagen/Satzumgebungen/Latex]]

\newtheorem{fakt}{Fakta}[section] %so wird abschnittsweise %durchnummeriert. Bei [] wird % einfach nummeriert.

\newtheorem{Fakt}{fakt}[Fakt]

\newtheorem{Proposition}[fakt]{Proposisi}

\newtheorem{Korollar}[fakt]{Korolari}

\newtheorem{Lemma}[fakt]{Lema}

\newtheorem{Satz}[fakt]{Teorema}

\newtheorem{Theorem}[fakt]{Teorema}

\theoremstyle{definition}

\newtheorem{Axiom}[fakt]{Aksioma}

\newtheorem{aufgabe}[fakt]{Soal}

\newtheorem{Aufgabe}[fakt]{Soal}

\newtheorem{klausuraufgabe}{Soal ujian}

\newtheorem{Klausuraufgabe}{Soal ujian}

\newtheorem{beispiel}[fakt]{Contoh}

\newtheorem{Beispiel}[fakt]{Contoh}

\newtheorem{Konstruktion}[fakt]{Konstruksi}

\newtheorem{konstruktion}[fakt]{Konstruksi}

\newtheorem{Verfahren}[fakt]{Verfahren}

\newtheorem{verfahren}[fakt]{Verfahren}

\newtheorem{bemerkung}[fakt]{Catatan}

\newtheorem{Bemerkung}[fakt]{Catatan}

\newtheorem{definition}[fakt]{Definisi}

\newtheorem{Definition}[fakt]{Definisi}

\newtheorem{notation}[fakt]{Notation}

\newtheorem{Notation}[fakt]{Notation}

\newtheorem{Frage}[fakt]{Frage}

\newtheorem{frage}[fakt]{Frage}

\newtheorem{problem}[fakt]{Problem}

\newtheorem{Problem}[fakt]{Problem}

\newtheorem{situation}[fakt]{Situation}

\newtheorem{Situation}[fakt]{Situation}

%Einlesegestaltung

%[[Projekt:Semantische Vorlagen/Umgebungsdesign/Latex]]

%Aufgabengestaltung

\newcommand{\inputaufgabe}[4]{\aufgabevorskip \begin{aufgabe} \punkte {#1} \aufgabepunktskip #2 #3 #4\end{aufgabe} \aufgabenachskip}

\newcommand{\inputaufgabegibtloesung}[4]{\aufgabevorskip \begin{aufgabe}\hspace{-1.8 mm}* \punkte {#1} \aufgabepunktskip #2 #3 #4\end{aufgabe} \aufgabenachskip}

\newcommand{\inputaufgabepunktenummervonhand}[3]{\aufgabevorskip {\bf Aufgabe #1} (#2 Punkte) \aufgabepunktskip #3  \aufgabenachskip}

\newcommand{\inputaufgabeloesung}[2]{ \begin{aufgabe} #1 \end{aufgabe} \aufgabenachskip

\bigskip

Lösung

\bigskip #2

}

\newcommand{\inputaufgabeklausurloesung}[3]{\par \newpage \begin{aufgabe} {\ifthenelse {\equal {#1}{}}{} {\ifthenelse {\equal {#1}{1}} {(1 Punkt)} {(#1 Punkte)}}} \par \bigskip #2 \end{aufgabe} \underline{L\"osung:} \par \bigskip #3 }

\newcommand{\inputaufgabeloesungvar}[2]{\aufgabevorskip Aufgabe: #1 \aufgabenachskip
Lösung: #2}

\newcommand{\inputaufgabepunkteloesung}[3]{\aufgabevorskip \begin{aufgabe} {\ifthenelse {\equal {#1}{}}{} {\ifthenelse {\equal {#1}{1}} {(1 Punkt)} {(#1 Punkte)}}}  \aufgabepunktskip #2 \end{aufgabe} \aufgabenachskip \loesung #3}

\newcommand{\loesung}[1]{\underline{Lösung:} #1}

\newcommand{\punkte}[1]{\ifthenelse {\equal {#1}{}}{} {\ifthenelse {\equal {#1}{1}} {(1 poin)} {(#1 poin)}}}

\newcommand{\inputexercise}[4]{\aufgabevorskip \begin{exercise} \points {#1} \aufgabepunktskip #2 #3 #4\end{exercise} \aufgabenachskip}

\newcommand{\points}[1]{\ifthenelse {\equal {#1}{}}{} {\ifthenelse {\equal {#1}{1}} {(1 point)} {(#1 points)}}}

%Beispielgestaltung

\newcommand{\inputbeispiel}[2]
{\beispielvorskip
\begin{beispiel}
\beispielbenennung{#1}
\beispielbenennungnachskip #2
\end{beispiel}
\beispielnachskip}

\newcommand{\beispielbenennung}[1]{ \ifthenelse {\equal {#1}{} \OR \equal {#1}{   }  }{}{(#1)}     }

%Bemerkunggestaltung

\newcommand{\inputbemerkung}[2]
{\bemerkungvorskip
\begin{bemerkung}
\bemerkungbenennung{#1}
\bemerkungbenennungnachskip #2
\end{bemerkung}
\bemerkungnachskip}

\newcommand{\bemerkungbenennung}[1]{  \ifthenelse {\equal {#1}{} \OR \equal {#1}{   }  }{}{(#1)}}

\newcommand{\inputverfahren}[2]
{\bemerkungvorskip
\begin{verfahren}
\bemerkungbenennung{#1}
\bemerkungbenennungnachskip #2
\end{verfahren}
\bemerkungnachskip}

\newcommand{\inputkonstruktion}[2]
{\bemerkungvorskip
\begin{konstruktion}
\bemerkungbenennung{#1}
\bemerkungbenennungnachskip #2
\end{konstruktion}
\bemerkungnachskip}

\newcommand{\inputfrage}[2]
{\bemerkungvorskip
\begin{frage}
\problembenennung{#1}
\bemerkungbenennungnachskip #2
\end{frage}
\bemerkungnachskip}

\newcommand{\fragebenennung}[1]{  \ifthenelse {\equal {#1}{} \OR \equal {#1}{   }  }{}{(#1)}}

\newcommand{\inputproblem}[2]
{\bemerkungvorskip
\begin{problem}
\problembenennung{#1}
\bemerkungbenennungnachskip #2
\end{problem}
\bemerkungnachskip}

\newcommand{\problembenennung}[1]{  \ifthenelse {\equal {#1}{} \OR \equal {#1}{   }  }{}{(#1)}}

%Situationsgestaltung

\newcommand{\inputsituation}[2]
{\begin{situation}
\situationbenennung{#1} #2
\end{situation}
}

\newcommand{\situationbenennung}[1]{\ifthenelse {\equal {#1}{} \OR \equal {#1}{  }  }{}{(#1)}}

%Definitiongestaltung

\newcommand{\inputdefinition}[2]
{\definitionvorskip
\begin{definition}
\definitionbenennung{#1}
\definitionbenennungnachskip #2
\end{definition}
\definitionnachskip}

\newcommand{\inputaxiom}[2]
{\definitionvorskip
\begin{Axiom}
\definitionbenennung{#1}
\definitionbenennungnachskip #2
\end{Axiom}
\definitionnachskip}

\newcommand{\definitionbenennung}[1]{\ifthenelse {\equal {#1}{} \OR \equal {#1}{  }  }{}{(#1)}}

\newcommand{\inputnotation}[2]
{\definitionvorskip \begin{notation} \definitionbenennung{#1} \definitionbenennungnachskip #2 \end{notation} \definitionnachskip}

%Fakt- und Beweisgestaltung

\renewcommand{\proofname}{\hspace{0cm}{\it Bukti}}

\newcommand{\beweis}[1]{\begin{proof} #1 \end{proof}}

\newcommand{\inputfakt}[4]
{\faktvorskip
\begin{#2}

%\label{#1}   (Absetzen, sonst  kann das folgende dahinter rutschen)

\faktbenennung{#3}
\faktbenennungnachskip #4
\end{#2}
\faktnachskip
}

\newcommand{\inputfaktbeweis}[5]
{\faktvorskip
\begin{#2}

%\label{#1}   (Absetzen, sonst  kann das folgende dahinter rutschen)

\faktbenennung{#3}
\faktbenennungnachskip #4
\end{#2}
\faktnachskip
\beweis{#5}}

\newcommand{\inputfaktbeweisnichtvorgefuehrt}[5]{\inputfaktbeweis {#1} {#2} {#3} {#4} {Dieser Beweis wurde in der Vorlesung nicht vorgef\"uhrt.} }

\newcommand{\inputfaktbeweistrivial}[4]{\inputfaktbeweis {#1} {#2} {#3} {#4} {Das ist trivial.} }

\newcommand{\inputfaktuebergangbeweis}[6]
{\faktvorskip
\begin{#2}

%\label{#1} (Absetzen, sonst kann das
% folgende dahinter rutschen)

\faktbenennung{#3}
\faktbenennungnachskip #4
\end{#2}
\faktnachskip #5 \faktnachskip
\beweis{#6}}

\newcommand{\inputbeweis}[1]{\beweis{#1} }

\newcommand{\faktbenennung}[1]{\ifthenelse {\equal {#1}{} \OR \equal {#1}{  }  }{}{(#1)}}

%Gestaltung der Faktstruktur

\newcommand{\faktsituation}[1]{\faktsituationskip #1}

\newcommand{\faktvoraussetzung}[1]{\faktvoraussetzungskip #1}

\newcommand{\faktvoraussetzungleer}[1]{ \hspace{-0,15cm} }

\newcommand{\faktvoraussetzungpos}[1]{\faktvoraussetzungskip #1}

\newcommand{\faktuebergang}[1]{\faktuebergangskip #1}

\newcommand{\faktuebergangleer}[1]{\faktuebergangskip \hspace{-0,15cm} }

\newcommand{\faktfolgerung}[1]{\faktfolgerungskip #1}

\newcommand{\faktzusatz}[1]{\faktzusatzskip #1}

%[[Projekt:Semantische Vorlagen/Beweisgestaltung/Latex]]

\newcommand{\teilbeweis}[5]{#1#2#3#4#5}

\newcommand{\fallunterscheidung}[5]{#1 #2\ifthenelse {\equal {#3}{}} {} { #3}\ifthenelse {\equal {#4}{}} {} { #4}\ifthenelse {\equal {#5}{}} {} { #5}}

\newcommand{\fallunterscheidungzwei}[2]{#1 #2}

\newcommand{\fallunterscheidungdrei}[3]{#1 #2 #3 }

\newcommand{\fallunterscheidungvier}[4]{#1 #2 #3 #4 }

\newcommand{\fallunterscheidungfuenf}[5]{#1 #2 #3 #4 #5 }

%[[Projekt:Semantische Vorlagen/Stichwortumsetzung/Latex]]

\newcommand{\betonung}[2]{\emph{#1}#2}

\newcommand{\stichwort}[2]{\emph{#1}#2}

\newcommand{\stichwortpraemath}[3]{$#1$-\emph{#2}#3}

\newcommand{\definitionswort}[2]{\emph{#1}#2}

\newcommand{\definitionswortpraemath}[3]{$#1$-\emph{#2}#3}

\newcommand{\definitionswortteil}[2]{\emph{#1}#2}

\newcommand{\definitionswortenp}[2]{\emph{#1}#2}

\newcommand{\definitionsverweis}[3]{#1#3}

\newcommand{\definitionsverweisanfuehrung}[3]{\anfuehrung {#1} {#3} }

%Anführungszeichen

\newcommand{\anfuehrung}[2]{"`{#1}"'#2}

\newcommand{\anfuehrungenglisch}[2]{``{#1}''#2}

%[[Projekt:Semantische Auflistungsvorlagen/Latex]]

%Aufzählungen

\newcommand{\aufzaehlungeins}[1]{\begin{enumerate} \item  #1 \end{enumerate}}

\newcommand{\aufzaehlungzwei}[2]{\begin{enumerate} \item  #1 \item  #2 \end{enumerate}}

\newcommand{\aufzaehlungdrei}[3]{\begin{enumerate} \item  #1 \item  #2 \item  #3  \end{enumerate}}

\newcommand{\aufzaehlungvier}[4]{\begin{enumerate} \item  #1 \item  #2 \item  #3 \item  #4 \end{enumerate}}

\newcommand{\aufzaehlungfuenf}[5]{\begin{enumerate} \item  #1 \item  #2 \item  #3 \item  #4 \item  #5  \end{enumerate}}

\newcommand{\aufzaehlungsechs}[6]{\begin{enumerate} \item  #1 \item  #2 \item  #3 \item  #4 \item  #5 \item  #6 \end{enumerate}}

\newcommand{\aufzaehlungsieben}[7]{\begin{enumerate} \item  #1 \item  #2 \item  #3 \item  #4 \item  #5 \item #6 \item  #7 \end{enumerate}}

\newcommand{\aufzaehlungacht}[8]{\begin{enumerate} \item  #1 \item  #2 \item  #3 \item  #4 \item  #5 \item #6 \item  #7 \item #8 \end{enumerate}}

\newcommand{\aufzaehlungneun}[9]{\begin{enumerate} \item  #1 \item  #2 \item  #3 \item  #4 \item  #5 \item #6 \item  #7  \item  #8  \item  #9 \end{enumerate}}

\newcommand{\itemfuenf}[5]{\item  #1 \item  #2 \item  #3 \item  #4 \item  #5}

\newcommand{\itemsechs}[6]{\item  #1 \item  #2 \item  #3 \item  #4 \item  #5 \item  #6}

\newcommand{\itemsieben}[7]{\item  #1 \item  #2 \item  #3 \item  #4 \item  #5 \item  #6 \item #7}

\newcommand{\aufzaehlungeinsabc}[1]{\begin{enumerate}[label=(\alph*)]  #1 \end{enumerate}}

\newcommand{\aufzaehlungzweiabc}[2]{\begin{enumerate}[label=(\alph*)] \item  #1 \item  #2 \end{enumerate}}

\newcommand{\aufzaehlungdreiabc}[3]{\begin{enumerate}[label=(\alph*)] \item  #1 \item  #2 \item  #3  \end{enumerate}}

\newcommand{\aufzaehlungvierabc}[4]{\begin{enumerate}[label=(\alph*)] \item  #1 \item  #2 \item  #3 \item  #4 \end{enumerate}}

\newcommand{\aufzaehlungfuenfabc}[5]{\begin{enumerate}[label=(\alph*)] \item  #1 \item  #2 \item  #3 \item  #4 \item  #5  \end{enumerate}}

\newcommand{\aufzaehlungsechsabc}[6]{\begin{enumerate}[label=(\alph*)] \item  #1 \item  #2 \item  #3 \item  #4 \item  #5 \item  #6 \end{enumerate}}

\newcommand{\aufzaehlungsiebenabc}[7]{\begin{enumerate}[label=(\alph*)] \item  #1 \item  #2 \item  #3 \item  #4 \item  #5 \item #6 \item  #7 \end{enumerate}}

\newcommand{\aufzaehlungachtabc}[8]{\begin{enumerate}[label=(\alph*)] \item  #1 \item  #2 \item  #3 \item  #4 \item  #5 \item #6 \item  #7 \item #8 \end{enumerate}}

\newcommand{\aufzaehlungneunabc}[9]{\begin{enumerate}[label=(\alph*)] \item  #1 \item  #2 \item  #3 \item  #4 \item  #5 \item #6 \item  #7  \item  #8  \item  #9 \end{enumerate}}

\newcommand{\aufzaehlungzweireihe}[2]{\begin{enumerate} #1 #2 \end{enumerate}}

\newcommand{\aufzaehlungdreibeweis}[3]{(1) #1 \par (2) #2 \par (3) #3 }

\newcommand{\aufzaehlungvierbeweis}[4]{(1) #1 \par (2) #2 \par (3) #3  \par (4) #4}

\newcommand{\aufzaehlungfuenfbeweis}[5]{(1) #1 \par (2) #2 \par (3) #3  \par (4) #4  \par (5) #5  }

%Auflistungen

\newcommand{\listitem}{\par \listskip \listdot}

\newcommand{\auflistungzwei}[2]{\par \listskip \listdot  #1 \par \listskip \listdot  #2 \par \listskip}

\newcommand{\auflistungdrei}[3]{\par \listskip \listdot  #1 \par \listskip \listdot  #2\par \listskip \listdot  #3\par\listskip }

\newcommand{\auflistungvier}[4]{\par \listskip \listdot  #1 \par \listskip \listdot  #2\par \listskip \listdot  #3\par \listskip \listdot  #4  \par \listskip }

\newcommand{\auflistungfuenf}[5]{\par \listskip \listdot  #1 \par \listskip \listdot  #2\par \listskip \listdot  #3\par \listskip \listdot  #4  \par \listskip  \listdot  #5 \par \listskip   }

\newcommand{\auflistungsechs}[6]{\par \listskip \listdot  #1 \par \listskip \listdot  #2\par \listskip \listdot  #3\par \listskip \listdot  #4  \par \listskip  \listdot  #5 \par \listskip \listdot  #6 \par \listskip}

\newcommand{\auflistungsieben}[7]{\par \listskip \listdot  #1 \par \listskip \listdot  #2\par \listskip \listdot  #3\par \listskip \listdot  #4  \par \listskip  \listdot  #5 \par \listskip \listdot  #6 \par \listskip \listdot  #7 \par \listskip }

\newcommand{\auflistungacht}[8]{\par \listskip \listdot  #1 \par \listskip \listdot  #2\par \listskip \listdot  #3\par \listskip \listdot  #4  \par \listskip  \listdot  #5 \par \listskip \listdot  #6 \par \listskip \listdot  #7 \par \listskip \listdot  #8  \par \listskip}

\newcommand{\auflistungneun}[9]{\par \listskip \listdot  #1 \par \listskip \listdot  #2\par \listskip \listdot  #3\par \listskip \listdot  #4  \par \listskip  \listdot  #5 \par \listskip \listdot  #6 \par \listskip \listdot  #7 \par \listskip \listdot  #8  \par \listskip  \listdot  #9 \par \listskip  }

\newcommand{\listitemfuenf}[5]{\listitem  #1 \listitem  #2 \listitem  #3 \listitem  #4 \listitem  #5}

\newcommand{\auflistungzweireihe}[2]{ #1 #2 }

%Gestaltung der Auflistungen

\newcommand{\listdot}{$\bullet \,$}

%Bildgestaltung

\newcommand{\inputbild}[1]{\includegraphics[scale]{#1}}

%Tabellengestaltung

%[[Projekt:Semantische Vorlagen/Zeilenkonfiguration/Latex]]

\newcommand{\zeileundzwei}[2]{#1 & #2}

\newcommand{\zeileunddrei}[3]{#1 & #2 & #3}

\newcommand{\zeileundvier}[4]{#1 & #2 & #3 & #4}

\newcommand{\zeileundfuenf}[5]{#1 & #2 & #3 & #4 & #5}

\newcommand{\zeileundsechs}[6]{#1 & #2 & #3 & #4 & #5 & #6}

\newcommand{\zeileundsieben}[7]{#1 & #2 & #3 & #4 & #5 & #6 & #7}

\newcommand{\zeileundacht}[8]{#1 & #2 & #3 & #4 & #5 & #6 & #7 & #8}

\newcommand{\zeileundneun}[9]{#1 & #2 & #3 & #4 & #5 & #6 & #7 & #8 & #9}

\newcommand{\mazeileundzwei}[2]{$#1$ & $#2$}

\newcommand{\mazeileunddrei}[3]{$#1$ & $ #2$ & $#3$}

\newcommand{\mazeileundvier}[4]{$#1$ & $#2$ & $#3$ & $#4$}

\newcommand{\mazeileundfuenf}[5]{$#1$ & $#2$ & $#3$ & $#4$ & $#5$}

\newcommand{\mazeileundsechs}[6]{$#1$ & $#2$ & $#3$ & $#4$ & $#5$ & $#6$}

\newcommand{\mazeileundsieben}[7]{$#1$ & $#2$ & $#3$ & $#4$ & $#5$ & $#6$ & $#7$}

\newcommand{\mazeileundacht}[8]{$#1$ & $#2$ & $#3$ & $#4$ & $#5$ & $#6$ & $#7$ & $#8$}

\newcommand{\mazeileundneun}[9]{$#1$ & $#2$ & $#3$ & $#4$ & $#5$ & $#6$ & $#7$ & $#8$ & $#9$}

%besser als Listen?
%Liste mit und (&) bzw. Bruch.

\newcommand {\listesechsmaund}[6]{$#1$ & $#2$& $#3$ & $#4$ & $#5$ & $#6$}

\newcommand {\listesiebenmaund}[7]{$#1$ & $#2$& $#3$ & $#4$ & $#5$ & $#6$ & $#7$}

\newcommand{\listezweiund}[2]{#1}

\newcommand{\listedreiund}[3]{#1}

\newcommand{\listesiebenund}[7]{#1}

\newcommand{\listesechsbruch}[6]{#1 \\ \hline #2 \\ \hline #3 \\ \hline #4 \\ \hline #5 \\ \hline #6}

\newcommand{\leitzeileunddrei}[3]{#1}

\newcommand{\madoppelzeileunddrei}[6]{$#1$&$#2$&$#3$\\ \hline $#4$&$#5$&$#6$}

%[[Projekt:Semantische Vorlagen/Tabellen/Latex]]

\newcommand{\tabelleviervier}[4]{\par \bigskip \begin{tabular}{|c|c|c|c|} \hline #1 \\ \hline #2 \\ \hline #3 \\ \hline #4 \\ \hline \end{tabular}}

\newcommand{\tabellefuenfvier}[5]{\par \bigskip \begin{tabular}{|c|c|c|c|c|} \hline #1 \\ \hline #2 \\ \hline #3 \\ \hline #4 \\ \hline #5  \\ \hline \end{tabular}}

%Mit math. Einträgen - ohne Leitzeile

\newcommand{\matabellezweivier}[5]{\begin{center} \begin{tabular}{|c|c|} \hline #2 \\ \hline #3 \\ \hline #4 \\ \hline #5  \\ \hline \end{tabular} \end{center} }

\newcommand{\matabellezweifuenf}[6]{\begin{center} \begin{tabular}{|c|c|} \hline     #2 \\ \hline #3 \\ \hline #4 \\ \hline #5 \\ \hline #6  \\ \hline \end{tabular} \end{center} }

\newcommand{\matabellezweisechs}[7]{\begin{center} \begin{tabular}{|c|c|} \hline     #2 \\ \hline #3 \\ \hline #4 \\ \hline #5 \\ \hline #6 \\ \hline #7   \\ \hline \end{tabular} \end{center} }

\newcommand{\matabellezweisieben}[8]{\begin{center} \begin{tabular}{|c|c|} \hline     #2 \\ \hline #3 \\ \hline #4 \\ \hline #5 \\ \hline #6 \\ \hline #7 \\ \hline #8  \\ \hline \end{tabular} \end{center} }

%Mit math. Einträgen - Mit Leitzeile

\newcommand{\matabellesiebenzwoelf}[3]{\begin{center} \begin{tabular}{|c||c|c|c|c|c|c|c|}  \hline  #1   \\ \hline \hline #2 \\ \hline #3 \\ \hline \end{tabular} \end{center} }

\newcommand{\matabellezwoelfzwoelf}[3]{\begin{center} \begin{tabular}{|c||c|c|c|c|c|c|c|c|c|c|c|c|}  \hline  #1   \\ \hline \hline #2 \\ \hline #3 \\ \hline \end{tabular} \end{center} }

\newcommand{\matabelledreineun}[5]{\begin{center} \begin{tabular}{|c|c|c|}\hline  #1 \\ \hline  #2 \\ \hline #3 \\ \hline   #4 \\ \hline  #5 \\ \hline \end{tabular} \end{center} }

\newcommand{\matabelledreizehn}[6]{\begin{center} \begin{tabular}{|c|c|c|}\hline  #1 \\ \hline  #2 \\ \hline #3 \\ \hline   #4 \\ \hline  #5 \\ \hline  #6 \\ \hline \end{tabular} \end{center} }

\newcommand{\matabelledreielf}[6]{\begin{center} \begin{tabular}{|c|c|c|}\hline  #1 \\ \hline  #2 \\ \hline #3 \\ \hline   #4 \\ \hline  #5 \\ \hline  #6 \\ \hline \end{tabular} \end{center} }

%[[Projekt:Semantische Vorlagen/Klausur/Punktetabellen/Latex]]

\newcommand{\punktetabellezehn}
{\vspace{3ex} \begin{center} \renewcommand{\arraystretch}{1.5} \tabcolsep=0.78ex \begin{tabular}{@{}|c|c|c|c|c|c|c|c|c|c|c|c|c|}
\hline Aufgabe & 1 & 2 & 3 & 4 & 5 & 6 & 7 & 8 & 9 & 10 & $\sum$
\\ \hline m\"ogliche Pkt. & \aeins & \azwei & \adrei & \avier & \afuenf & \asechs   & \asieben & \aacht  & \aneun & \azehn & \aelf
\\ \hline erhaltene Pkt.  & \phantom{11} & \phantom{11} & \phantom{11} & \phantom{11} & \phantom{11}& \phantom{11} & \phantom{11} & \phantom{11} & \phantom{11} & & \\ \hline \end{tabular} \end{center} \vspace{2ex} }

\newcommand{\punktetabelleelf}
{\vspace{3ex} \begin{center} \renewcommand{\arraystretch}{1.5} \tabcolsep=0.78ex \begin{tabular}{@{}|c|c|c|c|c|c|c|c|c|c|c|c|c|c|}
\hline Aufgabe & 1 & 2 & 3 & 4 & 5 & 6 & 7 & 8 & 9 & 10 & 11 & $\sum$
\\ \hline m\"ogliche Pkt. & \aeins & \azwei & \adrei & \avier & \afuenf & \asechs   & \asieben & \aacht  & \aneun & \azehn & \aelf & \azwoelf
\\ \hline erhaltene Pkt.  & \phantom{11} & \phantom{11} & \phantom{11} & \phantom{11} & \phantom{11}& \phantom{11} & \phantom{11} & \phantom{11} & \phantom{11} & \phantom{11} & & \\ \hline \end{tabular} \end{center} \vspace{2ex} }

\newcommand{\punktetabellezwoelf}
{\vspace{3ex} \begin{center} \renewcommand{\arraystretch}{1.5} \tabcolsep=0.78ex \begin{tabular}{@{}|c|c|c|c|c|c|c|c|c|c|c|c|c|c|c|}
\hline Aufgabe & 1 & 2 & 3 & 4 & 5 & 6 & 7 & 8 & 9 & 10 & 11 & 12    & $\sum$
\\ \hline m\"ogliche Pkt. & \aeins & \azwei & \adrei & \avier & \afuenf & \asechs   & \asieben & \aacht  & \aneun & \azehn & \aelf & \azwoelf & \adreizehn
\\ \hline erhaltene Pkt.  & \phantom{11} & \phantom{11} & \phantom{11} & \phantom{11} & \phantom{11}& \phantom{11} & \phantom{11} & \phantom{11}
& \phantom{11} & \phantom{11} & \phantom{11} & & \\ \hline \end{tabular} \end{center} \vspace{2ex} }

\newcommand{\punktetabelledreizehn}
{\vspace{3ex} \begin{center} \renewcommand{\arraystretch}{1.5} \tabcolsep=0.78ex \begin{tabular}{@{}|c|c|c|c|c|c|c|c|c|c|c|c|c|c|c|c|}
\hline Aufgabe & 1 & 2 & 3 & 4 & 5 & 6 & 7 & 8 & 9 & 10 & 11 & 12 & 13   & $\sum$
\\ \hline m\"ogliche Pkt. & \aeins & \azwei & \adrei & \avier & \afuenf & \asechs   & \asieben & \aacht  & \aneun & \azehn & \aelf & \azwoelf & \adreizehn & \avierzehn
\\ \hline erhaltene Pkt.  & \phantom{11} & \phantom{11} & \phantom{11} & \phantom{11} & \phantom{11}& \phantom{11} & \phantom{11} & \phantom{11} & \phantom{11}
& \phantom{11} & \phantom{11} & \phantom{11} & & \\ \hline \end{tabular} \end{center} \vspace{2ex} }

\newcommand{\punktetabellevierzehn}
{\vspace{3ex} \begin{center} \renewcommand{\arraystretch}{1.5} \tabcolsep=0.78ex \begin{tabular}{@{}|c|c|c|c|c|c|c|c|c|c|c|c|c|c|c|c|c|}
\hline Aufgabe & 1 & 2 & 3 & 4 & 5 & 6 & 7 & 8 & 9 & 10 & 11 & 12 & 13 & 14  & $\sum$
\\ \hline m\"ogliche Pkt. & \aeins & \azwei & \adrei & \avier & \afuenf & \asechs   & \asieben & \aacht  & \aneun & \azehn & \aelf & \azwoelf & \adreizehn & \avierzehn & \afuenfzehn
\\ \hline erhaltene Pkt. & \phantom{11} & \phantom{11} & \phantom{11} & \phantom{11} & \phantom{11} & \phantom{11}& \phantom{11} & \phantom{11} & \phantom{11} & \phantom{11}
& \phantom{11} & \phantom{11} & \phantom{11} & & \\ \hline \end{tabular} \end{center} \vspace{2ex} }

\newcommand{\punktetabellefuenfzehn}
{\vspace{3ex} \begin{center} \renewcommand{\arraystretch}{1.5} \tabcolsep=0.78ex \begin{tabular}{@{}|c|c|c|c|c|c|c|c|c|c|c|c|c|c|c|c|c|c|}
\hline Aufgabe & 1 & 2 & 3 & 4 & 5 & 6 & 7 & 8 & 9 & 10 & 11 & 12 & 13 & 14 & 15 & $\sum$
\\ \hline m\"ogliche Pkt. & \aeins & \azwei & \adrei & \avier & \afuenf & \asechs   & \asieben & \aacht  & \aneun & \azehn & \aelf & \azwoelf & \adreizehn & \avierzehn & \afuenfzehn
& \asechzehn
\\ \hline erhaltene Pkt. & \phantom{11} & \phantom{11} & \phantom{11}
& \phantom{11} & \phantom{11} & \phantom{11} & \phantom{11}& \phantom{11} & \phantom{11} & \phantom{11} & \phantom{11}
& \phantom{11} & \phantom{11} & \phantom{11} & & \\ \hline \end{tabular} \end{center} \vspace{2ex} }

\newcommand{\punktetabellesechzehn}
{\vspace{3ex} \begin{center} \renewcommand{\arraystretch}{1.5} \tabcolsep=0.78ex \begin{tabular}{@{}|c|c|c|c|c|c|c|c|c|c|c|c|c|c|c|c|c|c|}
\hline Aufgabe & 1 & 2 & 3 & 4 & 5 & 6 & 7 & 8 & 9 & 10 & 11 & 12 & 13 & 14 & 15 & 16 & $\sum$
\\ \hline m\"ogliche Pkt. & \aeins & \azwei & \adrei & \avier & \afuenf & \asechs   & \asieben & \aacht  & \aneun & \azehn & \aelf & \azwoelf & \adreizehn & \avierzehn & \afuenfzehn
& \asechzehn   & \asiebzehn
\\ \hline erhaltene Pkt. & \phantom{11} & \phantom{11} & \phantom{11} & \phantom{11}
& \phantom{11} & \phantom{11} & \phantom{11} & \phantom{11}& \phantom{11} & \phantom{11} & \phantom{11} & \phantom{11}
& \phantom{11} & \phantom{11} & \phantom{11} & & \\ \hline \end{tabular} \end{center} \vspace{2ex} }

\newcommand{\punktetabellesiebzehn}
{\vspace{3ex} \begin{center} \renewcommand{\arraystretch}{1.5} \tabcolsep=0.78ex \begin{tabular}{@{}|c|c|c|c|c|c|c|c|c|c|c|c|c|c|c|c|c|c|c|c|}
\hline Aufgabe & 1 & 2 & 3 & 4 & 5 & 6 & 7 & 8 & 9 & 10 & 11 & 12 & 13 & 14 & 15 & 16 & 17 & $\sum$
\\ \hline m\"ogliche Pkt. & \aeins & \azwei & \adrei & \avier & \afuenf & \asechs   & \asieben & \aacht  & \aneun & \azehn & \aelf & \azwoelf & \adreizehn & \avierzehn & \afuenfzehn
& \asechzehn   & \asiebzehn & \aachtzehn
\\ \hline erhaltene Pkt. & \phantom{11} & \phantom{11} & \phantom{11} & \phantom{11}
& \phantom{11} & \phantom{11} & \phantom{11} & \phantom{11} & \phantom{11}& \phantom{11} & \phantom{11} & \phantom{11} & \phantom{11}
& \phantom{11} & \phantom{11} & \phantom{11} & & \\ \hline \end{tabular} \end{center} \vspace{2ex} }

\newcommand{\punktetabelleachtzehn}
{\vspace{3ex} \begin{center} \renewcommand{\arraystretch}{1.5} \tabcolsep=0.78ex \begin{tabular}{@{}|c|c|c|c|c|c|c|c|c|c|c|c|c|c|c|c|c|c|c|c|c|}
\hline Aufgabe & 1 & 2 & 3 & 4 & 5 & 6 & 7 & 8 & 9 & 10 & 11 & 12 & 13 & 14 & 15 & 16 & 17 & 18 & $\sum$
\\ \hline m\"ogliche Pkt. & \aeins & \azwei & \adrei & \avier & \afuenf & \asechs   & \asieben & \aacht  & \aneun & \azehn & \aelf & \azwoelf & \adreizehn & \avierzehn & \afuenfzehn
& \asechzehn   & \asiebzehn & \aachtzehn & \aneunzehn
\\ \hline erhaltene Pkt. & \phantom{11} & \phantom{11} & \phantom{11} & \phantom{11} & \phantom{11}
& \phantom{11} & \phantom{11} & \phantom{11} & \phantom{11} & \phantom{11}& \phantom{11} & \phantom{11} & \phantom{11} & \phantom{11}
& \phantom{11} & \phantom{11} & \phantom{11} & & \\ \hline \end{tabular} \end{center} \vspace{2ex} }

\newcommand{\punktetabelleneunzehn}
{\vspace{3ex} \begin{center} \renewcommand{\arraystretch}{1.5} \tabcolsep=0.78ex \begin{tabular}{@{}|c|c|c|c|c|c|c|c|c|c|c|c|c|c|c|c|c|c|c|c|c|c|}
\hline Aufgabe & 1 & 2 & 3 & 4 & 5 & 6 & 7 & 8 & 9 & 10 & 11 & 12 & 13 & 14 & 15 & 16 & 17 & 18 & 19 & $\sum$
\\ \hline m\"ogliche Pkt. & \aeins & \azwei & \adrei & \avier & \afuenf & \asechs   & \asieben & \aacht  & \aneun & \azehn & \aelf & \azwoelf & \adreizehn & \avierzehn & \afuenfzehn
& \asechzehn   & \asiebzehn & \aachtzehn & \aneunzehn  & \azwanzig
\\ \hline erhaltene Pkt. & \phantom{11} & \phantom{11} & \phantom{11} & \phantom{11} & \phantom{11} & \phantom{11}
& \phantom{11} & \phantom{11} & \phantom{11} & \phantom{11} & \phantom{11}& \phantom{11} & \phantom{11} & \phantom{11} & \phantom{11}
& \phantom{11} & \phantom{11} & \phantom{11} & & \\ \hline \end{tabular} \end{center} \vspace{2ex} }

\newcommand{\punktetabellezwanzig}
{\vspace{3ex} \begin{center} \renewcommand{\arraystretch}{1.5} \tabcolsep=0.78ex \begin{tabular}{@{}|c|c|c|c|c|c|c|c|c|c|c|c|c|c|c|c|c|c|c|c|c|c|c|}
\hline Aufgabe & 1 & 2 & 3 & 4 & 5 & 6 & 7 & 8 & 9 & 10 & 11 & 12 & 13 & 14 & 15 & 16 & 17 & 18 & 19 & 20 & $\sum$
\\ \hline m\"ogliche Pkt. & \aeins & \azwei & \adrei & \avier & \afuenf & \asechs   & \asieben & \aacht  & \aneun & \azehn & \aelf & \azwoelf & \adreizehn & \avierzehn & \afuenfzehn
& \asechzehn   & \asiebzehn & \aachtzehn & \aneunzehn  & \azwanzig & \aeinundzwanzig
\\ \hline erhaltene Pkt. & \phantom{11} & \phantom{11} & \phantom{11} & \phantom{11} & \phantom{11} & \phantom{11} & \phantom{11}
& \phantom{11} & \phantom{11} & \phantom{11} & \phantom{11} & \phantom{11}& \phantom{11} & \phantom{11} & \phantom{11} & \phantom{11}
& \phantom{11} & \phantom{11} & \phantom{11} & & \\ \hline \end{tabular} \end{center} \vspace{2ex} }

\newcommand{\punktetabelleeinundzwanzig}
{\vspace{3ex} \begin{center} \renewcommand{\arraystretch}{1.5} \tabcolsep=0.78ex \begin{tabular}{@{}|c|c|c|c|c|c|c|c|c|c|c|c|c|c|c|c|c|c|c|c|c|c|c|c|}
\hline Aufgabe & 1 & 2 & 3 & 4 & 5 & 6 & 7 & 8 & 9 & 10 & 11 & 12 & 13 & 14 & 15 & 16 & 17 & 18 & 19 & 20 & 21 & $\sum$
\\ \hline m\"ogliche Pkt. & \aeins & \azwei & \adrei & \avier & \afuenf & \asechs   & \asieben & \aacht  & \aneun & \azehn & \aelf & \azwoelf & \adreizehn & \avierzehn & \afuenfzehn
& \asechzehn   & \asiebzehn & \aachtzehn & \aneunzehn  & \azwanzig & \aeinundzwanzig  & \azweiundzwanzig
\\ \hline erhaltene Pkt. & \phantom{11} & \phantom{11} & \phantom{11} & \phantom{11} & \phantom{11} & \phantom{11} & \phantom{11}
& \phantom{11} & \phantom{11} & \phantom{11} & \phantom{11} & \phantom{11}& \phantom{11} & \phantom{11} & \phantom{11} & \phantom{11}
& \phantom{11} & \phantom{11} & \phantom{11} & \phantom{11} & & \\ \hline \end{tabular} \end{center} \vspace{2ex} }

\newcommand{\punktetabellezweiundzwanzig}
{\vspace{3ex} \begin{center} \renewcommand{\arraystretch}{1.5} \tabcolsep=0.78ex \begin{tabular}{@{}|c|c|c|c|c|c|c|c|c|c|c|c|c|c|c|c|c|c|c|c|c|c|c|c|c|}
\hline Aufgabe & 1 & 2 & 3 & 4 & 5 & 6 & 7 & 8 & 9 & 10 & 11 & 12 & 13 & 14 & 15 & 16 & 17 & 18 & 19 & 20 & 21 & 22 & $\sum$
\\ \hline  Punkte & \aeins & \azwei & \adrei & \avier & \afuenf & \asechs   & \asieben & \aacht  & \aneun & \azehn & \aelf & \azwoelf & \adreizehn & \avierzehn & \afuenfzehn
& \asechzehn   & \asiebzehn & \aachtzehn & \aneunzehn  & \azwanzig & \aeinundzwanzig  & \azweiundzwanzig  & \adreiundzwanzig
\\ \hline Erhalten & \phantom{11} & \phantom{11} & \phantom{11} & \phantom{11} & \phantom{11} & \phantom{11} & \phantom{11}
& \phantom{11} & \phantom{11} & \phantom{11} & \phantom{11} & \phantom{11}& \phantom{11} & \phantom{11} & \phantom{11}  & \phantom{11} & \phantom{11}
& \phantom{11} & \phantom{11} & \phantom{11} & \phantom{11} & & \\ \hline \end{tabular} \end{center} \vspace{2ex} }

%[[Projekt:Semantische Vorlagen/Wahrheitstabellen/Latex]]

\newcommand{\wahrheitstabelleeins}[4]{
\begin{center}
\begin{tabular}{|c|c|}
\multicolumn{2}{l}{\textbf{#1}}\\
\hline
$p$ & $#2$\\ \hline
w  &	#3 \\ \hline
f  &	#4\\ \hline
\end{tabular}
\end{center} }

\newcommand{\wahrheitstabelleeinszwei}[7]{
\begin{center}
\begin{tabular}{|c|c|c|}
\multicolumn{3}{l}{\textbf{#1}}\\
\hline
$p$ & $#2$ & $#5$ \\ \hline
w  &	#3 &  #4 \\ \hline
f  &	#6 &  #7 \\ \hline
\end{tabular}
\end{center} }

\newcommand{\wahrheitstabelleeinsdrei}[4]{ \begin{center} \begin{tabular}{|c|c|c|c|} \multicolumn{4}{l}{\textbf{#1}} \\  \hline #2  \\  \hline #3  \\  \hline  #4  \\  \hline
\end{tabular} \end{center} }

%Wahrheitstabellen mit zwei Variablen p und q

\newcommand{\wahrheitstabellezweieins}[6]{ \begin{center} \begin{tabular}{|c|c|c|} \multicolumn{3}{l}{\textbf{#1}} \\  \hline  #2  \\  \hline  #3  \\  \hline  #4  \\  \hline  #5  \\  \hline  #6  \\  \hline \end{tabular} \end{center} }

\newcommand{\wahrheitstabellezweizwei}[6]{ \begin{center} \begin{tabular}{|c|c|c|c|} \multicolumn{4}{l}{\textbf{#1}} \\  \hline #2  \\  \hline #3  \\  \hline #4  \\  \hline #5  \\  \hline #6  \\  \hline\end{tabular} \end{center} }

\newcommand{\wahrheitstabellezweidrei}[6]{ \begin{center} \begin{tabular}{|c|c|c|c|c|} \multicolumn{5}{l}{\textbf{#1}} \\  \hline #2  \\  \hline #3  \\  \hline #4  \\  \hline #5  \\  \hline #6  \\  \hline \end{tabular} \end{center} }

\newcommand{\wahrheitstabellezweivier}[6]{ \begin{center} \begin{tabular}{|c|c|c|c|c|c|} \multicolumn{6}{l}{\textbf{#1}} \\  \hline #2  \\  \hline #3  \\  \hline #4  \\  \hline  #5  \\  \hline  #6  \\  \hline \end{tabular} \end{center} }

\newcommand{\wahrheitstabellezweifuenf}[6]{ \begin{center} \begin{tabular}{|c|c|c|c|c|c|c|} \multicolumn{7}{l}{\textbf{#1}} \\  \hline #2  \\  \hline #3  \\  \hline #4  \\  \hline #5  \\  \hline #6  \\  \hline
\end{tabular} \end{center} }

%Wahrheitstabellen mit drei Variablen p,q,r

\newcommand{\wahrheitstabelledreieins}[6]{ \begin{center} \begin{tabular}{|c|c|c|c||} \multicolumn{4}{l}{\textbf{#1}} \\  \hline #2  \\  \hline #3  \\  \hline #4  \\  \hline #5  \\  \hline #6 \\  \hline
\end{tabular} \end{center} }

%[[Projekt:Semantische Vorlagen/Wertetabellen/Latex]]

%Wertetabellen

\newcommand{\wertetabellevorskip}{\par \bigskip}

\newcommand{\wertetabellenachskip}{\par \bigskip}

\newcommand{\wertetabelleeinsausteilzeilen}[4]{\wertetabellevorskip \begin{tabular}{|c|c|} \hline #1 & #2  \\ \hline  #3 & #4  \\ \hline \end{tabular} \wertetabellenachskip }

\newcommand{\wertetabellezweiausteilzeilen}[4]{\wertetabellevorskip \begin{tabular}{|c|c|c|} \hline #1 & #2  \\ \hline  #3 & #4  \\ \hline \end{tabular} \wertetabellenachskip }

\newcommand{\wertetabelledreiausteilzeilen}[4]{\wertetabellevorskip \begin{tabular}{|c|c|c|c|} \hline #1 & #2  \\ \hline  #3 & #4  \\ \hline \end{tabular} \wertetabellenachskip }

\newcommand{\wertetabellevierausteilzeilen}[4]{\wertetabellevorskip \begin{tabular}{|c|c|c|c|c|} \hline #1 & #2  \\ \hline  #3 & #4  \\ \hline \end{tabular} \wertetabellenachskip }

\newcommand{\wertetabellefuenfausteilzeilen}[4]{\wertetabellevorskip \begin{tabular}{|c|c|c|c|c|c|} \hline #1 & #2  \\ \hline  #3 & #4  \\ \hline \end{tabular} \wertetabellenachskip }

\newcommand{\wertetabellesechsausteilzeilen}[6]{\wertetabellevorskip \begin{tabular}{|c|c|c|c|c|c|c|} \hline #1 & #2 & #3 \\ \hline  #4 & #5 & #6 \\ \hline \end{tabular} \wertetabellenachskip }

\newcommand{\wertetabellesiebenausteilzeilen}[6]{\wertetabellevorskip \begin{tabular}{|c|c|c|c|c|c|c|c|} \hline #1 & #2 & #3 \\ \hline  #4 & #5 & #6 \\ \hline \end{tabular} \wertetabellenachskip }

\newcommand{\wertetabelleachtausteilzeilen}[6]{\wertetabellevorskip \begin{tabular}{|c|c|c|c|c|c|c|c|c|} \hline #1 & #2 & #3 \\ \hline  #4 & #5 & #6 \\ \hline \end{tabular} \wertetabellenachskip }

\newcommand{\wertetabelleneunausteilzeilen}[6]{\wertetabellevorskip \begin{tabular}{|c|c|c|c|c|c|c|c|c|c|} \hline #1 & #2 & #3 \\ \hline  #4 & #5 & #6 \\ \hline \end{tabular} \wertetabellenachskip }

\newcommand{\wertetabellezehnausteilzeilen}[6]{\wertetabellevorskip \begin{tabular}{|c|c|c|c|c|c|c|c|c|c|c|} \hline #1 & #2 & #3 \\ \hline  #4 & #5 & #6 \\ \hline \end{tabular} \wertetabellenachskip }

\newcommand{\wertetabelleelfausteilzeilen}[8]{\wertetabellevorskip \begin{tabular}{|c|c|c|c|c|c|c|c|c|c|c|c|} \hline #1 & #2 & #3 & #4 \\ \hline #5 & #6 & #7 & #8 \\ \hline \end{tabular} \wertetabellenachskip }

%[[Projekt:Semantische Vorlagen/Verknüpfungstabellen/Latex]]

%Verknüpfungstabellen

\newcommand{\verknuepfungstabelleeins}[2]{\begin{center}
\begin{tabular}{|c||c|}
\hline #1 \\  \hline \hline #2 \\
\hline
\end{tabular}
\end{center} }

\newcommand{\verknuepfungstabellezwei}[3]{\begin{center}
\begin{tabular}{|c||c|c|}
\hline #1 \\  \hline \hline #2 \\ \hline #3 \\
\hline
\end{tabular}
\end{center} }

\newcommand{\verknuepfungstabelledrei}[4]{\begin{center}
\begin{tabular}{|c||c|c|c|}
\hline #1 \\  \hline \hline #2 \\ \hline #3 \\ \hline #4 \\
\hline
\end{tabular}
\end{center} }

\newcommand{\verknuepfungstabellevier}[5]{\begin{center}
\begin{tabular}{|c||c|c|c|c|}
\hline #1 \\  \hline \hline #2 \\ \hline #3 \\ \hline #4 \\  \hline  #5 \\
\hline
\end{tabular}
\end{center} }

\newcommand{\verknuepfungstabellefuenf}[6]{\begin{center}
\begin{tabular}{|c||c|c|c|c|c|}
\hline #1 \\  \hline \hline #2 \\ \hline #3 \\ \hline #4 \\  \hline  #5 \\ \hline  #6 \\
\hline
\end{tabular}
\end{center} }

\newcommand{\verknuepfungstabellesechs}[7]{\begin{center}
\begin{tabular}{|c||c|c|c|c|c|c|}
\hline #1 \\  \hline \hline #2 \\ \hline #3 \\ \hline #4 \\  \hline  #5 \\ \hline  #6 \\  \hline #7 \\
\hline
\end{tabular}
\end{center} }

\newcommand{\verknuepfungstabellesieben}[8]{\begin{center}
\begin{tabular}{|c||c|c|c|c|c|c|c|c|}
\hline #1 \\  \hline \hline #2 \\ \hline #3 \\ \hline #4 \\  \hline  #5 \\ \hline  #6 \\  \hline #7 \\  \hline  #8 \\
\hline
\end{tabular}
\end{center} }

\newcommand{\verknuepfungstabelleacht}[9]{\begin{center}
\begin{tabular}{|c||c|c|c|c|c|c|c|c|}
\hline #1 \\  \hline \hline #2 \\ \hline #3 \\ \hline #4 \\  \hline  #5 \\ \hline  #6 \\  \hline #7 \\  \hline  #8 \\ \hline  #9 \\
\hline
\end{tabular}
\end{center} }

%[[Projekt:Semantische Vorlagen/Initialisierung für Daten/Latex]]

\newcommand{\aeins}{    }

\newcommand{\azwei}{     }

\newcommand{\adrei}{    }

\newcommand{\avier}{    }

\newcommand{\afuenf}{    }

\newcommand{\asechs}{    }

\newcommand{\asieben}{    }

\newcommand{\aacht}{    }

\newcommand{\aneun}{    }

\newcommand{\azehn}{    }

\newcommand{\aelf}{    }

\newcommand{\azwoelf}{    }

\newcommand{\adreizehn}{     }

\newcommand{\avierzehn}{    }

\newcommand{\afuenfzehn}{    }

\newcommand{\asechzehn}{    }

\newcommand{\asiebzehn}{    }

\newcommand{\aachtzehn}{    }

\newcommand{\aneunzehn}{    }

\newcommand{\azwanzig}{    }

\newcommand{\aeinundzwanzig}{    }

\newcommand{\azweiundzwanzig}{    }

\newcommand{\adreiundzwanzig}{    }

\newcommand{\avierundzwanzig}{    }

\newcommand{\afuenfundzwanzig}{    }

\newcommand{\asechsundzwanzig}{    }

\newcommand{\asiebenundzwanzig}{    }

\newcommand{\aachtundzwanzig}{    }

\newcommand{\aneunundzwanzig}{    }

\newcommand{\leitzeilenull}{}

\newcommand{\leitzeileeins}{}

\newcommand{\leitzeilezwei}{}

\newcommand{\leitzeiledrei}{}

\newcommand{\leitzeilevier}{}

\newcommand{\leitzeilefuenf}{}

\newcommand{\leitzeilesechs}{}

\newcommand{\leitzeilesieben}{}

\newcommand{\leitzeileacht}{}

\newcommand{\leitzeileneun}{}

\newcommand{\leitzeilezehn}{}

\newcommand{\leitzeileelf}{}

\newcommand{\leitzeilezwoelf}{}

\newcommand{\leitspaltenull}{}

\newcommand{\leitspalteeins}{}

\newcommand{\leitspaltezwei}{}

\newcommand{\leitspaltedrei}{}

\newcommand{\leitspaltevier}{}

\newcommand{\leitspaltefuenf}{}

\newcommand{\leitspaltesechs}{}

\newcommand{\leitspaltesieben}{}

\newcommand{\leitspalteacht}{}

\newcommand{\leitspalteneun}{}

\newcommand{\leitspaltezehn}{}

\newcommand{\leitspalteelf}{}

\newcommand{\leitspaltezwoelf}{}

\newcommand{\leitspaltedreizehn}{}

\newcommand{\leitspaltevierzehn}{}

\newcommand{\leitspaltefuenfzehn}{}

\newcommand{\leitspaltesechzehn}{}

\newcommand{\leitspaltesiebzehn}{}

\newcommand{\leitspalteachtzehn}{}

\newcommand{\leitspalteneunzehn}{}

\newcommand{\leitspaltezwanzig}{}

\newcommand{\aeinsxeins}{}

\newcommand{\aeinsxzwei}{}

\newcommand{\aeinsxdrei}{}

\newcommand{\aeinsxvier}{}

\newcommand{\aeinsxfuenf}{}

\newcommand{\aeinsxsechs}{}

\newcommand{\aeinsxsieben}{}

\newcommand{\aeinsxacht}{}

\newcommand{\aeinsxneun}{}

\newcommand{\aeinsxzehn}{}

\newcommand{\aeinsxelf}{}

\newcommand{\aeinsxzwoelf}{}

\newcommand{\azweixeins}{}

\newcommand{\azweixzwei}{}

\newcommand{\azweixdrei}{}

\newcommand{\azweixvier}{}

\newcommand{\azweixfuenf}{}

\newcommand{\azweixsechs}{}

\newcommand{\azweixsieben}{}

\newcommand{\azweixacht}{}

\newcommand{\azweixneun}{}

\newcommand{\azweixzehn}{}

\newcommand{\azweixelf}{}

\newcommand{\azweixzwoelf}{}

\newcommand{\adreixeins}{}

\newcommand{\adreixzwei}{}

\newcommand{\adreixdrei}{}

\newcommand{\adreixvier}{}

\newcommand{\adreixfuenf}{}

\newcommand{\adreixsechs}{}

\newcommand{\adreixsieben}{}

\newcommand{\adreixacht}{}

\newcommand{\adreixneun}{}

\newcommand{\adreixzehn}{}

\newcommand{\adreixelf}{}

\newcommand{\adreixzwoelf}{}

\newcommand{\avierxeins}{}

\newcommand{\avierxzwei}{}

\newcommand{\avierxdrei}{}

\newcommand{\avierxvier}{}

\newcommand{\avierxfuenf}{}

\newcommand{\avierxsechs}{}

\newcommand{\avierxsieben}{}

\newcommand{\avierxacht}{}

\newcommand{\avierxneun}{}

\newcommand{\avierxzehn}{}

\newcommand{\avierxelf}{}

\newcommand{\avierxzwoelf}{}

\newcommand{\afuenfxeins}{}

\newcommand{\afuenfxzwei}{}

\newcommand{\afuenfxdrei}{}

\newcommand{\afuenfxvier}{}

\newcommand{\afuenfxfuenf}{}

\newcommand{\afuenfxsechs}{}

\newcommand{\afuenfxsieben}{}

\newcommand{\afuenfxacht}{}

\newcommand{\afuenfxneun}{}

\newcommand{\afuenfxzehn}{}

\newcommand{\afuenfxelf}{}

\newcommand{\afuenfxzwoelf}{}

\newcommand{\asechsxeins}{}

\newcommand{\asechsxzwei}{}

\newcommand{\asechsxdrei}{}

\newcommand{\asechsxvier}{}

\newcommand{\asechsxfuenf}{}

\newcommand{\asechsxsechs}{}

\newcommand{\asechsxsieben}{}

\newcommand{\asechsxacht}{}

\newcommand{\asechsxneun}{}

\newcommand{\asechsxzehn}{}

\newcommand{\asechsxelf}{}

\newcommand{\asechsxzwoelf}{}

\newcommand{\asiebenxeins}{}

\newcommand{\asiebenxzwei}{}

\newcommand{\asiebenxdrei}{}

\newcommand{\asiebenxvier}{}

\newcommand{\asiebenxfuenf}{}

\newcommand{\asiebenxsechs}{}

\newcommand{\asiebenxsieben}{}

\newcommand{\asiebenxacht}{}

\newcommand{\asiebenxneun}{}

\newcommand{\asiebenxzehn}{}

\newcommand{\asiebenxelf}{}

\newcommand{\asiebenxzwoelf}{}

\newcommand{\aachtxeins}{}

\newcommand{\aachtxzwei}{}

\newcommand{\aachtxdrei}{}

\newcommand{\aachtxvier}{}

\newcommand{\aachtxfuenf}{}

\newcommand{\aachtxsechs}{}

\newcommand{\aachtxsieben}{}

\newcommand{\aachtxacht}{}

\newcommand{\aachtxneun}{}

\newcommand{\aachtxzehn}{}

\newcommand{\aachtxelf}{}

\newcommand{\aachtxzwoelf}{}

\newcommand{\aneunxeins}{}

\newcommand{\aneunxzwei}{}

\newcommand{\aneunxdrei}{}

\newcommand{\aneunxvier}{}

\newcommand{\aneunxfuenf}{}

\newcommand{\aneunxsechs}{}

\newcommand{\aneunxsieben}{}

\newcommand{\aneunxacht}{}

\newcommand{\aneunxneun}{}

\newcommand{\aneunxzehn}{}

\newcommand{\aneunxelf}{}

\newcommand{\aneunxzwoelf}{}

\newcommand{\azehnxeins}{}

\newcommand{\azehnxzwei}{}

\newcommand{\azehnxdrei}{}

\newcommand{\azehnxvier}{}

\newcommand{\azehnxfuenf}{}

\newcommand{\azehnxsechs}{}

\newcommand{\azehnxsieben}{}

\newcommand{\azehnxacht}{}

\newcommand{\azehnxneun}{}

\newcommand{\azehnxzehn}{}

\newcommand{\azehnxelf}{}

\newcommand{\azehnxzwoelf}{}

\newcommand{\aelfxeins}{}

\newcommand{\aelfxzwei}{}

\newcommand{\aelfxdrei}{}

\newcommand{\aelfxvier}{}

\newcommand{\aelfxfuenf}{}

\newcommand{\aelfxsechs}{}

\newcommand{\aelfxsieben}{}

\newcommand{\aelfxacht}{}

\newcommand{\aelfxneun}{}

\newcommand{\aelfxzehn}{}

\newcommand{\aelfxelf}{}

\newcommand{\aelfxzwoelf}{}

\newcommand{\azwoelfxeins}{}

\newcommand{\azwoelfxzwei}{}

\newcommand{\azwoelfxdrei}{}

\newcommand{\azwoelfxvier}{}

\newcommand{\azwoelfxfuenf}{}

\newcommand{\azwoelfxsechs}{}

\newcommand{\azwoelfxsieben}{}

\newcommand{\azwoelfxacht}{}

\newcommand{\azwoelfxneun}{}

\newcommand{\azwoelfxzehn}{}

\newcommand{\azwoelfxelf}{}

\newcommand{\azwoelfxzwoelf}{}

\newcommand{\adreizehnxeins}{}

\newcommand{\adreizehnxzwei}{}

\newcommand{\adreizehnxdrei}{}

\newcommand{\adreizehnxvier}{}

\newcommand{\adreizehnxfuenf}{}

\newcommand{\adreizehnxsechs}{}

\newcommand{\adreizehnxsieben}{}

\newcommand{\adreizehnxacht}{}

\newcommand{\adreizehnxneun}{}

\newcommand{\adreizehnxzehn}{}

\newcommand{\adreizehnxelf}{}

\newcommand{\adreizehnxzwoelf}{}

\newcommand{\avierzehnxeins}{}

\newcommand{\avierzehnxzwei}{}

\newcommand{\avierzehnxdrei}{}

\newcommand{\avierzehnxvier}{}

\newcommand{\avierzehnxfuenf}{}

\newcommand{\avierzehnxsechs}{}

\newcommand{\avierzehnxsieben}{}

\newcommand{\avierzehnxacht}{}

\newcommand{\avierzehnxneun}{}

\newcommand{\avierzehnxzehn}{}

\newcommand{\avierzehnxelf}{}

\newcommand{\avierzehnxzwoelf}{}

\newcommand{\afuenfzehnxeins}{}

\newcommand{\afuenfzehnxzwei}{}

\newcommand{\afuenfzehnxdrei}{}

\newcommand{\afuenfzehnxvier}{}

\newcommand{\afuenfzehnxfuenf}{}

\newcommand{\afuenfzehnxsechs}{}

\newcommand{\afuenfzehnxsieben}{}

\newcommand{\afuenfzehnxacht}{}

\newcommand{\afuenfzehnxneun}{}

\newcommand{\afuenfzehnxzehn}{}

\newcommand{\afuenfzehnxelf}{}

\newcommand{\afuenfzehnxzwoelf}{}

\newcommand{\asechzehnxeins}{}

\newcommand{\asechzehnxzwei}{}

\newcommand{\asechzehnxdrei}{}

\newcommand{\asechzehnxvier}{}

\newcommand{\asechzehnxfuenf}{}

\newcommand{\asechzehnxsechs}{}

\newcommand{\asechzehnxsieben}{}

\newcommand{\asechzehnxacht}{}

\newcommand{\asechzehnxneun}{}

\newcommand{\asechzehnxzehn}{}

\newcommand{\asechzehnxelf}{}

\newcommand{\asechzehnxzwoelf}{}

\newcommand{\asiebzehnxeins}{}

\newcommand{\asiebzehnxzwei}{}

\newcommand{\asiebzehnxdrei}{}

\newcommand{\asiebzehnxvier}{}

\newcommand{\asiebzehnxfuenf}{}

\newcommand{\asiebzehnxsechs}{}

\newcommand{\asiebzehnxsieben}{}

\newcommand{\asiebzehnxacht}{}

\newcommand{\asiebzehnxneun}{}

\newcommand{\asiebzehnxzehn}{}

\newcommand{\asiebzehnxelf}{}

\newcommand{\asiebzehnxzwoelf}{}

\newcommand{\aachtzehnxeins}{}

\newcommand{\aachtzehnxzwei}{}

\newcommand{\aachtzehnxdrei}{}

\newcommand{\aachtzehnxvier}{}

\newcommand{\aachtzehnxfuenf}{}

\newcommand{\aachtzehnxsechs}{}

\newcommand{\aachtzehnxsieben}{}

\newcommand{\aachtzehnxacht}{}

\newcommand{\aachtzehnxneun}{}

\newcommand{\aachtzehnxzehn}{}

\newcommand{\aachtzehnxelf}{}

\newcommand{\aachtzehnxzwoelf}{}

%[[Projekt:Semantische Vorlagen/Tabellen mit benannten Parametern/Latex]]

%Tabellen mit zwei Zeilen

\newcommand{\tabelleleitzweixzwei}{

\begin{center}

\begin{tabular}{|c||c|c|}

\hline  \leitzeilenull  &  \leitzeileeins  &  \leitzeilezwei   \\

\hline \hline  \leitspalteeins  &  $ \aeinsxeins $  &  $ \aeinsxzwei $ \\

\hline  \leitspaltezwei  &  $ \azweixeins $  &  $ \azweixzwei $   \\

\hline

\end{tabular}

\end{center} }

%Tabellen mit drei Zeilen

\newcommand{\tabelleleitdreixdrei}{

\begin{center}

\begin{tabular}{|c||c|c|c||}

\hline  \leitzeilenull  &  \leitzeileeins  &  \leitzeilezwei  &  \leitzeiledrei        \\

\hline \hline  \leitspalteeins  &  $ \aeinsxeins $  &  $ \aeinsxzwei $  &  $ \aeinsxdrei $      \\

\hline  \leitspaltezwei  &  $ \azweixeins $  &  $ \azweixzwei $  &  $ \azweixdrei $       \\

\hline  \leitspaltedrei  &  $ \adreixeins $  &  $ \adreixzwei $  &  $ \adreixdrei $       \\

\hline

\end{tabular}

\end{center} }

\newcommand{\tabelleleitdreixvier}{

\begin{center}

\begin{tabular}{|c||c|c|c|c||}

\hline  \leitzeilenull  &  \leitzeileeins  &  \leitzeilezwei  &  \leitzeiledrei  &  \leitzeilevier      \\

\hline \hline  \leitspalteeins  &  $ \aeinsxeins $  &  $ \aeinsxzwei $  &  $ \aeinsxdrei $  &  $ \aeinsxvier $   \\

\hline  \leitspaltezwei  &  $ \azweixeins $  &  $ \azweixzwei $  &  $ \azweixdrei $  &  $ \azweixvier $     \\

\hline  \leitspaltedrei  &  $ \adreixeins $  &  $ \adreixzwei $  &  $ \adreixdrei $  &  $ \adreixvier $     \\

\hline

\end{tabular}

\end{center} }

\newcommand{\tabelleleitdreixsechs}{

\begin{center}

\begin{tabular}{|c||c|c|c|c|c|c|}

\hline  \leitzeilenull  &  \leitzeileeins  &  \leitzeilezwei  &  \leitzeiledrei  &  \leitzeilevier  &  \leitzeilefuenf   &  \leitzeilesechs    \\

\hline \hline  \leitspalteeins  &  $ \aeinsxeins $  &  $ \aeinsxzwei $  &  $ \aeinsxdrei $  &  $ \aeinsxvier $  &  $ \aeinsxfuenf $  &  $ \aeinsxsechs $  \\

\hline  \leitspaltezwei  &  $ \azweixeins $  &  $ \azweixzwei $  &  $ \azweixdrei $  &  $ \azweixvier $  &  $ \azweixfuenf $   &  $ \azweixsechs $  \\

\hline  \leitspaltedrei  &  $ \adreixeins $  &  $ \adreixzwei $  &  $ \adreixdrei $  &  $ \adreixvier $  &  $ \adreixfuenf $   &  $ \adreixsechs $   \\

\hline

\end{tabular}

\end{center} }

\newcommand{\tabelleleitdreixelf}{

\begin{center}

\begin{tabular}{|c||c|c|c|c|c|c|c|c|c|c|c|}

\hline  \leitzeilenull  &  \leitzeileeins  &  \leitzeilezwei  &  \leitzeiledrei  &  \leitzeilevier  &  \leitzeilefuenf   &  \leitzeilesechs  &  \leitzeilesieben  &  \leitzeileacht  &  \leitzeileneun  &  \leitzeilezehn &  \leitzeileelf    \\

\hline \hline  \leitspalteeins  &  $ \aeinsxeins $  &  $ \aeinsxzwei $  &  $ \aeinsxdrei $  &  $ \aeinsxvier $  &  $ \aeinsxfuenf $  &  $ \aeinsxsechs $  &  $ \aeinsxsieben $  &  $ \aeinsxacht $  &  $ \aeinsxneun $ &  $ \aeinsxzehn $  &  $ \aeinsxelf $  \\

\hline  \leitspaltezwei  &  $ \azweixeins $  &  $ \azweixzwei $  &  $ \azweixdrei $  &  $ \azweixvier $  &  $ \azweixfuenf $   &  $ \azweixsechs $  &  $ \azweixsieben $  &  $ \azweixacht $  &  $ \azweixneun $  &  $ \azweixzehn $ &  $ \azweixelf $     \\

\hline  \leitspaltedrei  &  $ \adreixeins $  &  $ \adreixzwei $  &  $ \adreixdrei $  &  $ \adreixvier $  &  $ \adreixfuenf $   &  $ \adreixsechs $  &  $ \adreixsieben $  &  $ \adreixacht $  &  $ \adreixneun $  &  $ \adreixzehn $     &  $ \adreixelf $   \\

\hline

\end{tabular}

\end{center} }

%Tabellen mit vier Zeilen

\newcommand{\tabelleleitvierxzwei}{

\begin{center}

\begin{tabular}{|c||c|c|}

\hline  \leitzeilenull  &  \leitzeileeins  &  \leitzeilezwei   \\

\hline \hline  \leitspalteeins  &  $ \aeinsxeins $  &  $ \aeinsxzwei $    \\

\hline  \leitspaltezwei  &  $ \azweixeins $  &  $ \azweixzwei $    \\

\hline  \leitspaltedrei  &  $ \adreixeins $  &  $ \adreixzwei $    \\

\hline  \leitspaltevier  &  $ \avierxeins $  &  $ \avierxzwei $    \\

\hline

\end{tabular}

\end{center} }

\newcommand{\tabelleleitvierxdrei}{

\begin{center}

\begin{tabular}{|c||c|c|c|}

\hline  \leitzeilenull  &  \leitzeileeins  &  \leitzeilezwei  &  \leitzeiledrei     \\

\hline \hline  \leitspalteeins  &  $ \aeinsxeins $  &  $ \aeinsxzwei $  &  $ \aeinsxdrei $     \\

\hline  \leitspaltezwei  &  $ \azweixeins $  &  $ \azweixzwei $  &  $ \azweixdrei $      \\

\hline  \leitspaltedrei  &  $ \adreixeins $  &  $ \adreixzwei $  &  $ \adreixdrei $    \\

\hline  \leitspaltevier  &  $ \avierxeins $  &  $ \avierxzwei $  &  $ \avierxdrei $    \\

\hline

\end{tabular}

\end{center} }

\newcommand{\tabelleleitvierxvier}{

\begin{center}

\begin{tabular}{|c||c|c|c|c|}

\hline  \leitzeilenull  &  \leitzeileeins  &  \leitzeilezwei  &  \leitzeiledrei  &  \leitzeilevier    \\

\hline \hline  \leitspalteeins  &  $ \aeinsxeins $  &  $ \aeinsxzwei $  &  $ \aeinsxdrei $  &  $ \aeinsxvier $   \\

\hline  \leitspaltezwei  &  $ \azweixeins $  &  $ \azweixzwei $  &  $ \azweixdrei $  &  $ \azweixvier $     \\

\hline  \leitspaltedrei  &  $ \adreixeins $  &  $ \adreixzwei $  &  $ \adreixdrei $  &  $ \adreixvier $   \\

\hline  \leitspaltevier  &  $ \avierxeins $  &  $ \avierxzwei $  &  $ \avierxdrei $  &  $ \avierxvier $   \\

\hline

\end{tabular}

\end{center} }

\newcommand{\tabelleleittextvierxzwei}{

	\begin{center}

		\begin{tabular}{|c||c|c|}

			\hline  \leitzeilenull  &  \leitzeileeins  &  \leitzeilezwei   \\

			\hline \hline  \leitspalteeins  &   \aeinsxeins   &   \aeinsxzwei    \\

			\hline  \leitspaltezwei  &  \azweixeins   &   \azweixzwei     \\

			\hline  \leitspaltedrei  &  \adreixeins   &   \adreixzwei     \\

			\hline  \leitspaltevier  &  \avierxeins   &   \avierxzwei     \\

			\hline

		\end{tabular}

	\end{center} }

\newcommand{\tabelleleitvierxacht}{

\begin{center}

\begin{tabular}{|c||c|c|c|c|c|c|c|c|}

\hline  \leitzeilenull  &  \leitzeileeins  &  \leitzeilezwei  &  \leitzeiledrei  &  \leitzeilevier  &  \leitzeilefuenf  &  \leitzeilesechs  &  \leitzeilesieben  &  \leitzeileacht  \\

\hline \hline  \leitspalteeins  &  $ \aeinsxeins $  &  $ \aeinsxzwei $  &  $ \aeinsxdrei $  &  $ \aeinsxvier $  &  $ \aeinsxfuenf $  &  $ \aeinsxsechs $  &  $ \aeinsxsieben $  &  $ \aeinsxacht $  \\

\hline  \leitspaltezwei  &  $ \azweixeins $  &  $ \azweixzwei $  &  $ \azweixdrei $  &  $ \azweixvier $   &  $ \azweixfuenf $  &  $ \azweixsechs $  &  $ \azweixsieben $  &  $ \azweixacht $   \\

\hline  \leitspaltedrei  &  $ \adreixeins $  &  $ \adreixzwei $  &  $ \adreixdrei $  &  $ \adreixvier $ &  $ \adreixfuenf $  &  $ \adreixsechs $  &  $ \adreixsieben $  &  $ \adreixacht $   \\

\hline  \leitspaltevier  &  $ \avierxeins $  &  $ \avierxzwei $  &  $ \avierxdrei $  &  $ \avierxvier $  &  $ \avierxfuenf $  &  $ \avierxsechs $  &  $ \avierxsieben $  &  $ \avierxacht $  \\

\hline

\end{tabular}

\end{center} }

%Tabellen mit fünf Zeilen

\newcommand{\tabelleleitfuenfxfuenf}{

\begin{center}

\begin{tabular}{|c||c|c|c|c|c|}

\hline  \leitzeilenull  &  \leitzeileeins  &  \leitzeilezwei  &  \leitzeiledrei  &  \leitzeilevier  &  \leitzeilefuenf   \\

\hline \hline  \leitspalteeins  &  $ \aeinsxeins $  &  $ \aeinsxzwei $  &  $ \aeinsxdrei $  &  $ \aeinsxvier $  &  $ \aeinsxfuenf $   \\

\hline  \leitspaltezwei  &  $ \azweixeins $  &  $ \azweixzwei $  &  $ \azweixdrei $  &  $ \azweixvier $  &  $ \azweixfuenf $     \\

\hline  \leitspaltedrei  &  $ \adreixeins $  &  $ \adreixzwei $  &  $ \adreixdrei $  &  $ \adreixvier $  &  $ \adreixfuenf $      \\

\hline  \leitspaltevier  &  $ \avierxeins $  &  $ \avierxzwei $  &  $ \avierxdrei $  &  $ \avierxvier $  &  $ \avierxfuenf $     \\

\hline  \leitspaltefuenf  &  $ \afuenfxeins $  &  $ \afuenfxzwei $  &  $ \afuenfxdrei $  &  $ \afuenfxvier $  &  $ \afuenfxfuenf $    \\

\hline

\end{tabular}

\end{center} }

%Tabellen mit sechs Zeilen

\newcommand{\tabelleleitsechsxdrei}{

\begin{center}

\begin{tabular}{|c||c|c|c|}

\hline  \leitzeilenull  &  \leitzeileeins  &  \leitzeilezwei  &  \leitzeiledrei    \\

\hline \hline  \leitspalteeins  &  $ \aeinsxeins $  &  $ \aeinsxzwei $  &  $ \aeinsxdrei $    \\

\hline  \leitspaltezwei  &  $ \azweixeins $  &  $ \azweixzwei $  &  $ \azweixdrei $      \\

\hline  \leitspaltedrei  &  $ \adreixeins $  &  $ \adreixzwei $  &  $ \adreixdrei $     \\

\hline  \leitspaltevier  &  $ \avierxeins $  &  $ \avierxzwei $  &  $ \avierxdrei $      \\

\hline  \leitspaltefuenf  &  $ \afuenfxeins $  &  $ \afuenfxzwei $  &  $ \afuenfxdrei $     \\

\hline  \leitspaltesechs  &  $ \asechsxeins $  &  $ \asechsxzwei $  &  $ \asechsxdrei $      \\

\hline

\end{tabular}

\end{center} }

\newcommand{\tabelleleitsechsxfuenf}{

\begin{center}

\begin{tabular}{|c||c|c|c|c|c|}

\hline  \leitzeilenull  &  \leitzeileeins  &  \leitzeilezwei  &  \leitzeiledrei  &  \leitzeilevier  &  \leitzeilefuenf   \\

\hline \hline  \leitspalteeins  &  $ \aeinsxeins $  &  $ \aeinsxzwei $  &  $ \aeinsxdrei $  &  $ \aeinsxvier $  &  $ \aeinsxfuenf $   \\

\hline  \leitspaltezwei  &  $ \azweixeins $  &  $ \azweixzwei $  &  $ \azweixdrei $  &  $ \azweixvier $  &  $ \azweixfuenf $     \\

\hline  \leitspaltedrei  &  $ \adreixeins $  &  $ \adreixzwei $  &  $ \adreixdrei $  &  $ \adreixvier $  &  $ \adreixfuenf $      \\

\hline  \leitspaltevier  &  $ \avierxeins $  &  $ \avierxzwei $  &  $ \avierxdrei $  &  $ \avierxvier $  &  $ \avierxfuenf $     \\

\hline  \leitspaltefuenf  &  $ \afuenfxeins $  &  $ \afuenfxzwei $  &  $ \afuenfxdrei $  &  $ \afuenfxvier $  &  $ \afuenfxfuenf $    \\

\hline  \leitspaltesechs  &  $ \asechsxeins $  &  $ \asechsxzwei $  &  $ \asechsxdrei $  &  $ \asechsxvier $  &  $ \asechsxfuenf $    \\

\hline

\end{tabular}

\end{center} }

\newcommand{\tabelleleitsechsxsechs}{

\begin{center}

\begin{tabular}{|c||c|c|c|c|c|c|}

\hline  \leitzeilenull  &  \leitzeileeins  &  \leitzeilezwei  &  \leitzeiledrei  &  \leitzeilevier  &  \leitzeilefuenf  &  \leitzeilesechs   \\

\hline \hline  \leitspalteeins  &  $ \aeinsxeins $  &  $ \aeinsxzwei $  &  $ \aeinsxdrei $  &  $ \aeinsxvier $  &  $ \aeinsxfuenf $  &  $ \aeinsxsechs $ \\

\hline  \leitspaltezwei  &  $ \azweixeins $  &  $ \azweixzwei $  &  $ \azweixdrei $  &  $ \azweixvier $  &  $ \azweixfuenf $  &  $ \azweixsechs $   \\

\hline  \leitspaltedrei  &  $ \adreixeins $  &  $ \adreixzwei $  &  $ \adreixdrei $  &  $ \adreixvier $  &  $ \adreixfuenf $   &  $ \adreixsechs $   \\

\hline  \leitspaltevier  &  $ \avierxeins $  &  $ \avierxzwei $  &  $ \avierxdrei $  &  $ \avierxvier $  &  $ \avierxfuenf $   &  $ \avierxsechs $  \\

\hline  \leitspaltefuenf  &  $ \afuenfxeins $  &  $ \afuenfxzwei $  &  $ \afuenfxdrei $  &  $ \afuenfxvier $  &  $ \afuenfxfuenf $ &  $ \afuenfxsechs $   \\

\hline  \leitspaltesechs  &  $ \asechsxeins $  &  $ \asechsxzwei $  &  $ \asechsxdrei $  &  $ \asechsxvier $  &  $ \asechsxfuenf $  &  $ \asechsxsechs $  \\

\hline

\end{tabular}

\end{center} }

%Tabellen mit sieben Zeilen

\newcommand{\tabelleleitsiebenxsieben}{

\begin{center}

\begin{tabular}{|c||c|c|c|c|c|c|c|}

\hline  \leitzeilenull  &  \leitzeileeins  &  \leitzeilezwei  &  \leitzeiledrei  &  \leitzeilevier  &  \leitzeilefuenf  &  \leitzeilesechs &  \leitzeilesieben  \\

\hline \hline  \leitspalteeins  &  $ \aeinsxeins $  &  $ \aeinsxzwei $  &  $ \aeinsxdrei $  &  $ \aeinsxvier $  &  $ \aeinsxfuenf $  &  $ \aeinsxsechs $ &  $ \aeinsxsieben $ \\

\hline  \leitspaltezwei  &  $ \azweixeins $  &  $ \azweixzwei $  &  $ \azweixdrei $  &  $ \azweixvier $  &  $ \azweixfuenf $  &  $ \azweixsechs $  &  $ \azweixsieben $ \\

\hline  \leitspaltedrei  &  $ \adreixeins $  &  $ \adreixzwei $  &  $ \adreixdrei $  &  $ \adreixvier $  &  $ \adreixfuenf $   &  $ \adreixsechs $ &  $ \adreixsieben $  \\

\hline  \leitspaltevier  &  $ \avierxeins $  &  $ \avierxzwei $  &  $ \avierxdrei $  &  $ \avierxvier $  &  $ \avierxfuenf $   &  $ \avierxsechs $ &  $ \avierxsieben $ \\

\hline  \leitspaltefuenf  &  $ \afuenfxeins $  &  $ \afuenfxzwei $  &  $ \afuenfxdrei $  &  $ \afuenfxvier $  &  $ \afuenfxfuenf $ &  $ \afuenfxsechs $ &  $ \afuenfxsieben $  \\

\hline  \leitspaltesechs  &  $ \asechsxeins $  &  $ \asechsxzwei $  &  $ \asechsxdrei $  &  $ \asechsxvier $  &  $ \asechsxfuenf $  &  $ \asechsxsechs $ &  $ \asechsxsieben $ \\

\hline  \leitspaltesieben  &  $ \asiebenxeins $  &  $ \asiebenxzwei $  &  $ \asiebenxdrei $  &  $ \asiebenxvier $  &  $ \asiebenxfuenf $  &  $ \asiebenxsechs $ &  $ \asiebenxsieben $ \\

\hline

\end{tabular}

\end{center} }

%Tabellen mit acht Zeilen

\newcommand{\tabelleleitachtxdrei}{

\begin{center}

\begin{tabular}{|c||c|c|c|}

\hline  \leitzeilenull  &  \leitzeileeins  &  \leitzeilezwei  &  \leitzeiledrei     \\

\hline \hline  \leitspalteeins  &  $ \aeinsxeins $  &  $ \aeinsxzwei $  &  $ \aeinsxdrei $    \\

\hline  \leitspaltezwei  &  $ \azweixeins $  &  $ \azweixzwei $  &  $ \azweixdrei $     \\

\hline  \leitspaltedrei  &  $ \adreixeins $  &  $ \adreixzwei $  &  $ \adreixdrei $    \\

\hline  \leitspaltevier  &  $ \avierxeins $  &  $ \avierxzwei $  &  $ \avierxdrei $      \\

\hline  \leitspaltefuenf  &  $ \afuenfxeins $  &  $ \afuenfxzwei $  &  $ \afuenfxdrei $     \\

\hline  \leitspaltesechs  &  $ \asechsxeins $  &  $ \asechsxzwei $  &  $ \asechsxdrei $      \\

\hline  \leitspaltesieben  &  $ \asiebenxeins $  &  $ \asiebenxzwei $  &  $ \asiebenxdrei $     \\

\hline  \leitspalteacht  &  $ \aachtxeins $  &  $ \aachtxzwei $  &  $ \aachtxdrei $      \\

\hline

\end{tabular}

\end{center} }

%Tabellen mit zehn Zeilen

\newcommand{\tabelleleitzehnxdrei}{

\begin{center}

\begin{tabular}{|c||c|c|c|}

\hline  \leitzeilenull  &  \leitzeileeins  &  \leitzeilezwei  &  \leitzeiledrei     \\

\hline \hline  \leitspalteeins  &  $ \aeinsxeins $  &  $ \aeinsxzwei $  &  $ \aeinsxdrei $    \\

\hline  \leitspaltezwei  &  $ \azweixeins $  &  $ \azweixzwei $  &  $ \azweixdrei $     \\

\hline  \leitspaltedrei  &  $ \adreixeins $  &  $ \adreixzwei $  &  $ \adreixdrei $    \\

\hline  \leitspaltevier  &  $ \avierxeins $  &  $ \avierxzwei $  &  $ \avierxdrei $      \\

\hline  \leitspaltefuenf  &  $ \afuenfxeins $  &  $ \afuenfxzwei $  &  $ \afuenfxdrei $     \\

\hline  \leitspaltesechs  &  $ \asechsxeins $  &  $ \asechsxzwei $  &  $ \asechsxdrei $      \\

\hline  \leitspaltesieben  &  $ \asiebenxeins $  &  $ \asiebenxzwei $  &  $ \asiebenxdrei $     \\

\hline  \leitspalteacht  &  $ \aachtxeins $  &  $ \aachtxzwei $  &  $ \aachtxdrei $      \\

\hline  \leitspalteneun  &  $ \aneunxeins $  &  $ \aneunxzwei $  &  $ \aneunxdrei $      \\

\hline  \leitspaltezehn  &  $ \azehnxeins $  &  $ \azehnxzwei $  &  $ \azehnxdrei $      \\

\hline

\end{tabular}

\end{center} }

%Tabellen mit elf Zeilen

\newcommand{\tabelleleitelfxvier}{

\begin{center}

\begin{tabular}{|c||c|c|c|c|}

\hline  \leitzeilenull  &  \leitzeileeins  &  \leitzeilezwei  &  \leitzeiledrei  &  \leitzeilevier  \\

\hline \hline  \leitspalteeins  &  $ \aeinsxeins $  &  $ \aeinsxzwei $  &  $ \aeinsxdrei $  &  $ \aeinsxvier $  \\

\hline  \leitspaltezwei  &  $ \azweixeins $  &  $ \azweixzwei $  &  $ \azweixdrei $  &  $ \azweixvier $   \\

\hline  \leitspaltedrei  &  $ \adreixeins $  &  $ \adreixzwei $  &  $ \adreixdrei $  &  $ \adreixvier $   \\

\hline  \leitspaltevier  &  $ \avierxeins $  &  $ \avierxzwei $  &  $ \avierxdrei $  &  $ \avierxvier $   \\

\hline  \leitspaltefuenf  &  $ \afuenfxeins $  &  $ \afuenfxzwei $  &  $ \afuenfxdrei $  &  $ \afuenfxvier $   \\

\hline  \leitspaltesechs  &  $ \asechsxeins $  &  $ \asechsxzwei $  &  $ \asechsxdrei $  &  $ \asechsxvier $   \\

\hline  \leitspaltesieben  &  $ \asiebenxeins $  &  $ \asiebenxzwei $  &  $ \asiebenxdrei $  &  $ \asiebenxvier $    \\

\hline  \leitspalteacht  &  $ \aachtxeins $  &  $ \aachtxzwei $  &  $ \aachtxdrei $  &  $ \aachtxvier $   \\

\hline  \leitspalteneun  &  $ \aneunxeins $  &  $ \aneunxzwei $  &  $ \aneunxdrei $  &  $ \aneunxvier $   \\

\hline  \leitspaltezehn  &  $ \azehnxeins $  &  $ \azehnxzwei $  &  $ \azehnxdrei $  &  $ \azehnxvier $   \\

\hline  \leitspalteelf  &  $ \aelfxeins $  &  $ \aelfxzwei $  &  $ \aelfxdrei $  &  $ \aelfxvier $   \\

\hline

\end{tabular}

\end{center} }

%Tabellen mit zwölf Zeilen

\newcommand{\tabelleleitzwoelfxsieben}{

\begin{center}

\begin{tabular}{|c||c|c|c|c|c|c|}

\hline  \leitzeilenull  &  \leitzeileeins  &  \leitzeilezwei  &  \leitzeiledrei  &  \leitzeilevier &  \leitzeilefuenf  &  \leitzeilesechs &  \leitzeilesieben \\

\hline \hline  \leitspalteeins  &  $ \aeinsxeins $  &  $ \aeinsxzwei $  &  $ \aeinsxdrei $  &  $ \aeinsxvier $  &  $ \aeinsxfuenf $  &  $ \aeinsxsechs $  &  $ \aeinsxsieben $  \\

\hline  \leitspaltezwei  &  $ \azweixeins $  &  $ \azweixzwei $  &  $ \azweixdrei $  &  $ \azweixvier $  &  $ \azweixfuenf $  &  $ \azweixsechs $  &  $ \azweixsieben $   \\

\hline  \leitspaltedrei  &  $ \adreixeins $  &  $ \adreixzwei $  &  $ \adreixdrei $  &  $ \adreixvier $  &  $ \adreixfuenf $  &  $ \adreixsechs $  &  $ \adreixsieben $   \\

\hline  \leitspaltevier  &  $ \avierxeins $  &  $ \avierxzwei $  &  $ \avierxdrei $  &  $ \avierxvier $   &  $ \avierxfuenf $  &  $ \avierxsechs $  &  $ \avierxsieben $  \\

\hline  \leitspaltefuenf  &  $ \afuenfxeins $  &  $ \afuenfxzwei $  &  $ \afuenfxdrei $  &  $ \afuenfxvier $ &  $ \afuenfxfuenf $  &  $ \afuenfxsechs $  &  $ \afuenfxsieben $  \\

\hline  \leitspaltesechs  &  $ \asechsxeins $  &  $ \asechsxzwei $  &  $ \asechsxdrei $  &  $ \asechsxvier $  &  $ \asechsxfuenf $  &  $ \asechsxsechs $  &  $ \asechsxsieben $  \\

\hline  \leitspaltesieben  &  $ \asiebenxeins $  &  $ \asiebenxzwei $  &  $ \asiebenxdrei $  &  $ \asiebenxvier $  &  $ \asiebenxfuenf $  &  $ \asiebenxsechs $  &  $ \asiebenxsieben $   \\

\hline  \leitspalteacht  &  $ \aachtxeins $  &  $ \aachtxzwei $  &  $ \aachtxdrei $  &  $ \aachtxvier $  &  $ \aachtxfuenf $  &  $ \aachtxsechs $  &  $ \aachtxsieben $ \\

\hline  \leitspalteneun  &  $ \aneunxeins $  &  $ \aneunxzwei $  &  $ \aneunxdrei $  &  $ \aneunxvier $  &  $ \aneunxfuenf $  &  $ \aneunxsechs $  &  $ \aneunxsieben $  \\

\hline  \leitspaltezehn  &  $ \azehnxeins $  &  $ \azehnxzwei $  &  $ \azehnxdrei $  &  $ \azehnxvier $ &  $ \azehnxfuenf $  &  $ \azehnxsechs $  &  $ \azehnxsieben $   \\

\hline  \leitspalteelf  &  $ \aelfxeins $  &  $ \aelfxzwei $  &  $ \aelfxdrei $  &  $ \aelfxvier $  &  $ \aelfxfuenf $  &  $ \aelfxsechs $  &  $ \aelfxsieben $  \\

\hline  \leitspaltezwoelf  &  $ \azwoelfxeins $  &  $ \azwoelfxzwei $  &  $ \azwoelfxdrei $  &  $ \zwoelfxvier $  &  $ \azwoelfxfuenf $  &  $ \azwoelfxsechs $  &  $ \zwoelfxsieben $ \\

\hline

\end{tabular}

\end{center} }

%Tabellen mit sechzehn Zeilen

\newcommand{\tabelleleitsechzehnxzwei}{

\begin{center}

\begin{tabular}{|c||c|c|}

\hline  \leitzeilenull  &  \leitzeileeins  &  \leitzeilezwei   \\

\hline \hline  \leitspalteeins  &  $ \aeinsxeins $  &  $ \aeinsxzwei $ \\

\hline  \leitspaltezwei  &  $ \azweixeins $  &  $ \azweixzwei $  \\

\hline  \leitspaltedrei  &  $ \adreixeins $  &  $ \adreixzwei $  \\

\hline  \leitspaltevier  &  $ \avierxeins $  &  $ \avierxzwei $  \\

\hline  \leitspaltefuenf  &  $ \afuenfxeins $  &  $ \afuenfxzwei $  \\

\hline  \leitspaltesechs  &  $ \asechsxeins $  &  $ \asechsxzwei $  \\

\hline  \leitspaltesieben  &  $ \asiebenxeins $  &  $ \asiebenxzwei $  \\

\hline  \leitspalteacht  &  $ \aachtxeins $  &  $ \aachtxzwei $  \\

\hline  \leitspalteneun  &  $ \aneunxeins $  &  $ \aneunxzwei $   \\

\hline  \leitspaltezehn  &  $ \azehnxeins $  &  $ \azehnxzwei $  \\

\hline  \leitspalteelf  &  $ \aelfxeins $  &  $ \aelfxzwei $   \\

\hline  \leitspaltezwoelf  &  $ \azwoelfxeins $  &  $ \azwoelfxzwei $   \\

\hline  \leitspaltedreizehn  &  $ \adreizehnxeins $  &  $ \adreizehnxzwei $   \\

\hline  \leitspaltevierzehn  &  $ \avierzehnxeins $  &  $ \avierzehnxzwei $  \\

\hline  \leitspaltefuenfzehn  &  $ \afuenfzehnxeins $  &  $ \afuenfzehnxzwei $   \\

\hline  \leitspaltesechzehn  &  $ \asechzehnxeins $  &  $ \asechzehnxzwei $   \\

\hline

\end{tabular}

\end{center} }

\newcommand{\tabelleleitsechzehnxsechs}{

\begin{center}

\begin{tabular}{|c||c|c|c|c|c|c|}

\hline  \leitzeilenull  &  \leitzeileeins  &  \leitzeilezwei  &  \leitzeiledrei  &  \leitzeilevier  &  \leitzeilefuenf  &  \leitzeilesechs  \\

\hline \hline  \leitspalteeins  &  $ \aeinsxeins $  &  $ \aeinsxzwei $  &  $ \aeinsxdrei $  &  $ \aeinsxvier $  &  $ \aeinsxfuenf $  &  $ \aeinsxsechs $  \\

\hline  \leitspaltezwei  &  $ \azweixeins $  &  $ \azweixzwei $  &  $ \azweixdrei $  &  $ \azweixvier $  &  $ \azweixfuenf $  &  $ \azweixsechs $  \\

\hline  \leitspaltedrei  &  $ \adreixeins $  &  $ \adreixzwei $  &  $ \adreixdrei $  &  $ \adreixvier $  &  $ \adreixfuenf $  &  $ \adreixsechs $  \\

\hline  \leitspaltevier  &  $ \avierxeins $  &  $ \avierxzwei $  &  $ \avierxdrei $  &  $ \avierxvier $  &  $ \avierxfuenf $  &  $ \avierxsechs $  \\

\hline  \leitspaltefuenf  &  $ \afuenfxeins $  &  $ \afuenfxzwei $  &  $ \afuenfxdrei $  &  $ \afuenfxvier $  &  $ \afuenfxfuenf $  &  $ \afuenfxsechs $  \\

\hline  \leitspaltesechs  &  $ \asechsxeins $  &  $ \asechsxzwei $  &  $ \asechsxdrei $  &  $ \asechsxvier $  &  $ \asechsxfuenf $  &  $ \asechsxsechs $  \\

\hline  \leitspaltesieben  &  $ \asiebenxeins $  &  $ \asiebenxzwei $  &  $ \asiebenxdrei $  &  $ \asiebenxvier $  &  $ \asiebenxfuenf $  &  $ \asiebenxsechs $  \\

\hline  \leitspalteacht  &  $ \aachtxeins $  &  $ \aachtxzwei $  &  $ \aachtxdrei $  &  $ \aachtxvier $  &  $ \aachtxfuenf $  &  $ \aachtxsechs $  \\

\hline  \leitspalteneun  &  $ \aneunxeins $  &  $ \aneunxzwei $  &  $ \aneunxdrei $  &  $ \aneunxvier $  &  $ \aneunxfuenf $  &  $ \aneunxsechs $  \\

\hline  \leitspaltezehn  &  $ \azehnxeins $  &  $ \azehnxzwei $  &  $ \azehnxdrei $  &  $ \azehnxvier $  &  $ \azehnxfuenf $  &  $ \azehnxsechs $  \\

\hline  \leitspalteelf  &  $ \aelfxeins $  &  $ \aelfxzwei $  &  $ \aelfxdrei $  &  $ \aelfxvier $  &  $ \aelfxfuenf $  &  $ \aelfxsechs $  \\

\hline  \leitspaltezwoelf  &  $ \azwoelfxeins $  &  $ \azwoelfxzwei $  &  $ \azwoelfxdrei $  &  $ \azwoelfxvier $  &  $ \azwoelfxfuenf $  &  $ \azwoelfxsechs $  \\

\hline  \leitspaltedreizehn  &  $ \adreizehnxeins $  &  $ \adreizehnxzwei $  &  $ \adreizehnxdrei $  &  $ \adreizehnxvier $  &  $ \adreizehnxfuenf $  &  $ \adreizehnxsechs $  \\

\hline  \leitspaltevierzehn  &  $ \avierzehnxeins $  &  $ \avierzehnxzwei $  &  $ \avierzehnxdrei $  &  $ \avierzehnxvier $  &  $ \avierzehnxfuenf $  &  $ \avierzehnxsechs $  \\

\hline  \leitspaltefuenfzehn  &  $ \afuenfzehnxeins $  &  $ \afuenfzehnxzwei $  &  $ \afuenfzehnxdrei $  &  $ \afuenfzehnxvier $  &  $ \afuenfzehnxfuenf $  &  $ \afuenfzehnxsechs $  \\

\hline  \leitspaltesechzehn  &  $ \asechzehnxeins $  &  $ \asechzehnxzwei $  &  $ \asechzehnxdrei $  &  $ \asechzehnxvier $  &  $ \asechzehnxfuenf $  &  $ \asechzehnxsechs $  \\

\hline

\end{tabular}

\end{center} }

%Tabellen mit achtzehn Zeilen

\newcommand{\tabelleleitachtzehnxneun}{

\begin{center}

\begin{tabular}{|c||c|c|c|c|c|c|c|c|c|}

\hline  \leitzeilenull  &  \leitzeileeins  &  \leitzeilezwei  &  \leitzeiledrei  &  \leitzeilevier  &  \leitzeilefuenf  &  \leitzeilesechs &  \leitzeilesieben  &  \leitzeileacht  &  \leitzeileneun \\

\hline \hline  \leitspalteeins  &  $ \aeinsxeins $  &  $ \aeinsxzwei $  &  $ \aeinsxdrei $  &  $ \aeinsxvier $  &  $ \aeinsxfuenf $  &  $ \aeinsxsechs $ &  $ \aeinsxsieben $  &  $ \aeinsxacht $  &  $ \aeinsxneun $  \\

\hline  \leitspaltezwei  &  $ \azweixeins $  &  $ \azweixzwei $  &  $ \azweixdrei $  &  $ \azweixvier $  &  $ \azweixfuenf $  &  $ \azweixsechs $ &  $ \azweixsieben $  &  $ \azweixacht $  &  $ \azweixneun $  \\

\hline  \leitspaltedrei  &  $ \adreixeins $  &  $ \adreixzwei $  &  $ \adreixdrei $  &  $ \adreixvier $  &  $ \adreixfuenf $  &  $ \adreixsechs $ &  $ \adreixsieben $  &  $ \adreixacht $  &  $ \adreixneun $  \\

\hline  \leitspaltevier  &  $ \avierxeins $  &  $ \avierxzwei $  &  $ \avierxdrei $  &  $ \avierxvier $  &  $ \avierxfuenf $  &  $ \avierxsechs $  &  $ \avierxsieben $  &  $ \avierxacht $  &  $ \avierxneun $ \\

\hline  \leitspaltefuenf  &  $ \afuenfxeins $  &  $ \afuenfxzwei $  &  $ \afuenfxdrei $  &  $ \afuenfxvier $  &  $ \afuenfxfuenf $  &  $ \afuenfxsechs $ &  $ \afuenfxsieben $  &  $ \afuenfxacht $  &  $ \afuenfxneun $  \\

\hline  \leitspaltesechs  &  $ \asechsxeins $  &  $ \asechsxzwei $  &  $ \asechsxdrei $  &  $ \asechsxvier $  &  $ \asechsxfuenf $  &  $ \asechsxsechs $ &  $ \asechsxsieben $  &  $ \asechsxacht $  &  $ \asechsxneun $  \\

\hline  \leitspaltesieben  &  $ \asiebenxeins $  &  $ \asiebenxzwei $  &  $ \asiebenxdrei $  &  $ \asiebenxvier $  &  $ \asiebenxfuenf $  &  $ \asiebenxsechs $ &  $ \asiebenxsieben $  &  $ \asiebenxacht $  &  $ \asiebenxneun $  \\

\hline  \leitspalteacht  &  $ \aachtxeins $  &  $ \aachtxzwei $  &  $ \aachtxdrei $  &  $ \aachtxvier $  &  $ \aachtxfuenf $  &  $ \aachtxsechs $ &  $ \aachtxsieben $  &  $ \aachtxacht $  &  $ \aachtxneun $  \\

\hline  \leitspalteneun  &  $ \aneunxeins $  &  $ \aneunxzwei $  &  $ \aneunxdrei $  &  $ \aneunxvier $  &  $ \aneunxfuenf $  &  $ \aneunxsechs $ &  $ \aneunxsieben $  &  $ \aneunxacht $  &  $ \aneunxneun $  \\

\hline  \leitspaltezehn  &  $ \azehnxeins $  &  $ \azehnxzwei $  &  $ \azehnxdrei $  &  $ \azehnxvier $  &  $ \azehnxfuenf $  &  $ \azehnxsechs $ &  $ \azehnxsieben $  &  $ \azehnxacht $  &  $ \azehnxneun $  \\

\hline  \leitspalteelf  &  $ \aelfxeins $  &  $ \aelfxzwei $  &  $ \aelfxdrei $  &  $ \aelfxvier $  &  $ \aelfxfuenf $  &  $ \aelfxsechs $ &  $ \aelfxsieben $  &  $ \aelfxacht $  &  $ \aelfxneun $  \\

\hline  \leitspaltezwoelf  &  $ \azwoelfxeins $  &  $ \azwoelfxzwei $  &  $ \azwoelfxdrei $  &  $ \azwoelfxvier $  &  $ \azwoelfxfuenf $  &  $ \azwoelfxsechs $ &  $ \azwoelfxsieben $  &  $ \azwoelfxacht $  &  $ \azwoelfxneun $  \\

\hline  \leitspaltedreizehn  &  $ \adreizehnxeins $  &  $ \adreizehnxzwei $  &  $ \adreizehnxdrei $  &  $ \adreizehnxvier $  &  $ \adreizehnxfuenf $  &  $ \adreizehnxsechs $ &  $ \adreizehnxsieben $  &  $ \adreizehnxacht $  &  $ \adreizehnxneun $  \\

\hline  \leitspaltevierzehn  &  $ \avierzehnxeins $  &  $ \avierzehnxzwei $  &  $ \avierzehnxdrei $  &  $ \avierzehnxvier $  &  $ \avierzehnxfuenf $  &  $ \avierzehnxsechs $  &  $ \avierzehnxsieben $  &  $ \avierzehnxacht $  &  $ \avierzehnxneun $ \\

\hline  \leitspaltefuenfzehn  &  $ \afuenfzehnxeins $  &  $ \afuenfzehnxzwei $  &  $ \afuenfzehnxdrei $  &  $ \afuenfzehnxvier $  &  $ \afuenfzehnxfuenf $  &  $ \afuenfzehnxsechs $ &  $ \afuenfzehnxsieben $  &  $ \afuenfzehnxacht $  &  $ \afuenfzehnxneun $  \\

\hline  \leitspaltesechzehn  &  $ \asechzehnxeins $  &  $ \asechzehnxzwei $  &  $ \asechzehnxdrei $  &  $ \asechzehnxvier $  &  $ \asechzehnxfuenf $  &  $ \asechzehnxsechs $ &  $ \asechzehnxsieben $  &  $ \asechzehnxacht $  &  $ \asechzehnxneun $  \\

\hline  \leitspaltesiebzehn  &  $ \asiebzehnxeins $  &  $ \asiebzehnxzwei $  &  $ \asiebzehnxdrei $  &  $ \asiebzehnxvier $  &  $ \asiebzehnxfuenf $  &  $ \asiebzehnxsechs $ &  $ \asiebzehnxsieben $  &  $ \asiebzehnxacht $  &  $ \asiebzehnxneun $  \\

\hline  \leitspalteachtzehn  &  $ \aachtzehnxeins $  &  $ \aachtzehnxzwei $  &  $ \aachtzehnxdrei $  &  $ \aachtzehnxvier $  &  $ \aachtzehnxfuenf $  &  $ \aachtzehnxsechs $ &  $ \aachtzehnxsieben $  &  $ \aachtzehnxacht $  &  $ \aachtzehnxneun $  \\

\hline

\end{tabular}

\end{center} }

\newcommand{\bruchhilfealign}{\cr & & }

\newcommand{\bruchhilfedisp}{\]  \[ }

\newcommand{\bruchhilfe}{ \- }

%[[Projekt:Semantische Vorlagen/Notlösungen/Latex]]

%
%Bei verschachtelten ifthen-Befehlen
%kann es Probleme geben, daher
%

\newcommand{\faktuebergangpos}[1]{ \faktuebergangskip  {\hspace{-0,55cm} } {#1} }

%Ende des Grundvorspannes

%Globale Strukturierungen

\newcommand{\seitenueberschrift}[1]{\seitenueberschriftvorskip
\begin{center} \large{ #1} \end{center} \seitenueberschriftnachskip}

\newcommand{\zwischenueberschrift}[1]{\section*{#1} }

\newcommand{\zwischenueberschriftaufgaben}[1]{ \zwischenueberschriftvorskip \begin{center} {\bf #1} \end{center} \zwischenueberschriftnachskip}

\newcommand{\unterueberschrift}[1]{ \unterueberschriftvorskip \hspace{0.5cm} {\em #1}  \unterueberschriftnachskip}

\renewcommand{\definitionbenennung}[1]{}

%Grund:Bei Definitionen im Normaltext
%nicht nötig, da Begriff früh kommt

\renewcommand{\beispielbenennung}[1]{}

\renewcommand{\bemerkungbenennung}[1]{}

%Umbrüche bei großer Schrift (entfällt)

%[[Projekt:Semantische Vorlagen/Mathematischer Modus/Umbrüche/Latex]]

\newcommand{\mathdisplaybruch}{}

\newcommand{\mathdisplaybruchinklammer}{}

\newcommand{\mathl}[2]{$#1$#2}

\newcommand{\mathlk}[2]{$#1$#2}

%Nicht durch semantische Vorlage produzierte Befehle

\newcommand{\mathdisplaybruchbruch}{ \]  \[ }

%Gestaltung der vertikalen und horizontalen Abstände

%[[Projekt:Semantische Vorlagen/Latex/Vorlesungsskript/Skips]]

%Abstände bei Überschriften

\newcommand{\seitenueberschriftvorskip}{\par \bigskip \bigskip \bigskip }

\newcommand{\seitenueberschriftnachskip}{\par \bigskip}

\newcommand{\zwischenueberschriftvorskip}{\par \bigskip}

\newcommand{\zwischenueberschriftnachskip}{\par \smallskip}

\newcommand{\unterueberschriftvorskip}{\par \bigskip}

\newcommand{\unterueberschriftnachskip}{\par \smallskip}

%Aufzählungen und Listen

\newcommand{\listskip}{\smallskip}

%Abstände bei mathematischen Strukturen

\newcommand{\aufgabevorskip}{\bigskip}

\newcommand{\aufgabepunktskip}{\par}

\newcommand{\aufgabenachskip}{}

\newcommand{\aufgabeloesungskip}{\par \bigskip}

\newcommand{\beispielvorskip}{}

\newcommand{\beispielnachskip}{}

\newcommand{\beispielbenennungnachskip}{}

\newcommand{\bemerkungvorskip}{}

\newcommand{\bemerkungnachskip}{}

\newcommand{\bemerkungbenennungnachskip}{}

\newcommand{\definitionvorskip}{}

\newcommand{\definitionnachskip}{}

\newcommand{\definitionbenennungnachskip}{}

\newcommand{\faktvorskip}{}

\newcommand{\faktnachskip}{}

\newcommand{\faktbenennungvorskip}{}

\newcommand{\faktbenennungnachskip}{}

\newcommand{\faktsituationskip}{}

\newcommand{\faktvoraussetzungskip}{}

\newcommand{\faktuebergangskip}{}

\newcommand{\faktfolgerungskip}{}

\newcommand{\faktzusatzskip}{}

%Sonstiges

\newcommand{\deckblattabstand}{\vspace{2cm}}

\newcommand{\gesamtskriptnewpage}{ }

\newcommand{\einzeltextnewpage}{ }

%[[Projekt:Semantische Vorlagen/Latex/Vorlesungsskript/Horizontale Abstände]]

\newcommand{\leerzeichen}{ }

%Seitengestaltung

%[[Projekt:Semantische Vorlagen/Latex/Vorlesungsskript/Seitengestaltung]]

%Einzelne Abstände

% page geometry supplied by reader wrapper
% page geometry supplied by reader wrapper

% page geometry supplied by reader wrapper
% page geometry supplied by reader wrapper

% page geometry supplied by reader wrapper
% page geometry supplied by reader wrapper

\setlength{\parindent}{0cm}

\setlength{\parskip}{1ex}

%Bildgestaltung

%[[Projekt:Semantische Vorlagen/Latex/Vorlesungsskript/Bildgestaltung]]

\newcommand{\bild}[1]{\vspace{4mm} #1}

\newcommand{\bildtext}[1]{\begin{center} {\small #1 } \end{center} }

%Lokale Einstellungen

%\input{Lokalername}

\newcommand{\bildeinlesung}[1]{../authority/media/#1}

\renewcommand{\bildeinlesung}[2]{../authority/media/#1.#2}

\newcommand{\bildeinlesungpng}[2]{../authority/media/#1.png}

\newcommand{\bildeinlesungPNG}[2]{../authority/media/#1.png}

\newcommand{\bildeinlesungjpg}[2]{../authority/media/#1.jpg}

\newcommand{\bildeinlesungJPG}[2]{../authority/media/#1.jpg}

\newcommand{\bildeinlesungjpeg}[2]{../authority/media/#1.jpg}

\newcommand{\bildeinlesungsvg}[2]{generated/media/#1.png}

\newcommand{\bildeinlesunggif}[2]{../authority/media/#1.png}

\newcommand{\bildeinlesungGIF}[2]{../authority/media/#1.png}

\newcommand{\bildeinlesungxcf}[2]{../authority/media/#1.png}

%Klausurgestaltung

%[[Projekt:Semantische Vorlagen/Klausurgestaltung/Latex]]

\newcommand{\fachbereich}{Fachbereich}

\newcommand{\dozent}{Dozent}

\newcommand{\klausurdatum}{Datum}

\newcommand{\klausurgebiet}{Mathematik}

\newcommand{\klausurtyp}{Klausur}

\newcommand{\klausurmodus}{Dauer: Zehn Minuten Einlesezeit und zwei volle Stunden Bearbeitungszeit. Es sind keine Hilfsmittel erlaubt. Alle Antworten sind zu begr\"unden.

Es gibt insgesamt 64 Punkte. Es gilt die Sockelregelung, d.h. die Bewertung pro Aufgabe(nteil) beginnt in der Regel bei der halben Punktzahl. Zum Bestehen braucht man $16$ Punkte, ab $32$ Punkten gibt es eine Eins.

Tragen Sie auf dem Deckblatt und jedem weiteren Blatt Ihren Namen und Ihre
Matrikelnummer leserlich ein. Viel Erfolg!}

\newcommand{\klausurname}{
\renewcommand{\arraystretch}{1.5}
\begin{tabular}{@{}p{3.2cm}p{10cm}}
Name, Vorname: &
..................................................................................\\
Matrikelnummer: &
..................................................................................\\
\end{tabular}
}

\newcommand{\klausurvorspann}[5]{\textbf{\fachbereich}\hfill \textbf{\klausurdatum}
\\
\textbf{\dozent} \hfill
\vspace{4ex}

\begin{center}
{\bfseries{\large \klausurgebiet \vspace{2ex}

\klausurtyp} }
\end{center}

\vspace{2ex}

\klausurmodus

\vspace{2ex}

\klausurname
}

\newcommand{\klausurnote}{\begin{tabular}{@{}p{3.2cm}p{10cm}} Note: & \end{tabular} }

\newcommand{\klausurtaeuschungwarnung}{ \vspace{2ex}{Mir ist bewu\ss t, dass bei einem schweren T\"auschungsversuch s\"amtliche Pr\"u\-fungsvorleis\-tungen in diesem Kurs verfallen. Ein schwerer T\"auschungsversuch liegt insbesondere beim Einsatz von Literatur und von elektronischen Hilfsmitteln vor. \vspace{1ex} } }

\newcommand{\klausureinverstaendnis}{
\vspace{2ex}
Ich erkl\"are mich durch meine Unterschrift einverstanden, dass mein Klausurergebnis mit meiner Matrikelnummer im Internet bekanntgegeben wird.
\vspace{1ex}

\begin{tabular}{@{}p{3.2cm}p{10cm}}
Unterschrift: &
.....................................................................................
\end{tabular}
}

%Variante für Testklausur

\newcommand{\testklausurvorspann}[5]{\textbf{\fachbereich}\hfill \textbf{\klausurdatum} \\ \textbf{\dozent} \hfill \vspace{4ex}

\begin{center} {\bfseries{\large \klausurgebiet \vspace{2ex}

\klausurtyp} } \end{center}

\vspace{2ex}

\testklausurmodus

\vspace{3ex}

\klausurname }

\newcommand{\testklausurmodus}{Dauer: Zehn Minuten Einlesezeit und zwei volle Stunden Bearbeitungszeit.

Es sind keine Hilfsmittel erlaubt.

Alle Antworten sind zu begr\"unden.

Es gibt insgesamt 64 Punkte. Es gilt die Sockelregelung, d.h. die Bewertung pro Aufgabe(nteil) beginnt in der Regel bei der halben Punktzahl.

Zur Orientierung: Zum Bestehen braucht man $16$ Punkte, ab $32$ Punkten gibt es eine Eins.

Tragen Sie auf dem Deckblatt und jedem weiteren Blatt Ihren Namen und Ihre Matrikelnummer leserlich ein.

Viel Erfolg!}

%\renewcommand{\inputaufgabegibtloesung}[4]{\aufgabevorskip \begin{aufgabe} \punkte {#1} \aufgabepunktskip #2 #3 #4\end{aufgabe} \aufgabenachskip}

\newcommand{\klausurmitloesungenvorspann}[5]{\textbf{\fachbereich}\hfill \textbf{\klausurdatum} \\ \textbf{\dozent} \hfill \vspace{4ex}

\begin{center} {\bfseries{\large \klausurgebiet \vspace{2ex}

\klausurtyp \hspace{0mm} mit L\"osungen} } \end{center}

}

%Lizenzierung und Bildverzeichnis

%[[Projekt:Semantische Vorlagen/Latex/Vorlesungsskript/Lizenzierung]]

\newcommand{\Lizenztext}{Teks sumber dibuat oleh Holger Brenner di Wikiversity berbahasa Jerman dan digunakan di sini berdasarkan CC BY-SA 4.0. Media mengikuti lisensi per berkas.}

\newcommand{\Abbildungslizenzierung}{ Erl\"auterung: Die in diesem Text verwendeten Bilder stammen aus Commons (also von http://commons.wikimedia.org) und haben eine Lizenz, die die Verwendung hier erlaubt. Die Bilder werden mit ihren Dateinamen auf Commons angef\"uhrt zusammen mit ihrem Autor bzw. Hochlader und der Lizenz. }

%Konfiguration des Bildverzeichnisses

\newcommand{\bildlizenz}[6]{
\ifthenelse {\equal{#2}{}}
{\addcontentsline{lof}{figure}{Sumber = #1, Kreator = Pengguna #3 di #4, Lisensi = #5 \bildlizenzskip  }}
{ \ifthenelse {\equal{#3}{}}
{ \addcontentsline{lof}{figure}{ Sumber = #1, Kreator = #2, Lisensi = #5 \bildlizenzskip }}
{ \addcontentsline{lof}{figure}{ Sumber = #1, Kreator = #2 (diunggah oleh Pengguna #3 di #4), Lisensi = #5 \bildlizenzskip }} } }

\newcommand{\bildlizenzskip}{ }

%Englische Variante

\newcommand{\ignoriereeng}[1]{}

\ignoriereeng{

%[[Projekt:Semantische Vorlagen/Latex/Englische Zusätze]]

\newtheorem{corollary}[fakt]{Corollary}

\newtheorem{Corollary}[fakt]{Corollary}

\theoremstyle{definition}

\newtheorem{example}[fakt]{Example}

\newtheorem{remark}[fakt]{Remark}

\newcommand{\inputremark}[2] {\bemerkungvorskip \begin{remark} \bemerkungbenennung{#1} \bemerkungbenennungnachskip #2 \end{remark} \bemerkungnachskip}

\newcommand{\inputexample}[2] {\beispielvorskip \begin{example} \beispielbenennung{#1} \beispielbenennungnachskip #2 \end{example} \beispielnachskip}

\newcommand{\inputfaktproof}[5] {\faktvorskip \begin{#2}

%\label{#1} (Absetzen, sonst kann das folgende dahinter rutschen)

\faktbenennung{#3} \faktbenennungnachskip #4 \end{#2} \faktnachskip \beweis{#5}}

\renewcommand{\proofname}{\hspace{-0.64cm}{\it Proof}}

}

%Variante für Artikel

\newcommand{\ignoriereart}[1]{}

\ignoriereart{
\newcommand{\titelueberschrift}[1]{\title[#1] { #1} \maketitle }

\renewcommand{\seitenueberschrift}[1]{ \section {#1} }

\renewcommand{\zwischenueberschrift}[1]{ \subsection {#1} }
}

% Indonesian aliases used by the independent translated reader.
\newtheorem{Teorema}[fakt]{Teorema}
\renewcommand{\proofname}{Bukti}

% The frozen MediaWiki TeX uses underscore-normalized PNG names while some
% canonical Commons filenames retain spaces. The bounded cumulative build
% creates exact generated aliases and routes PNG macros only through that
% generated directory; canonical source files remain untouched.
\renewcommand{\bildeinlesungpng}[2]{generated/media/#1.png}
\renewcommand{\bildeinlesungPNG}[2]{generated/media/#1.png}
% The frozen source numbers facts and exercises by lecture/worksheet section.
% Book chapters are reader scaffolding and must not enter source identifiers.
\renewcommand{\thefakt}{\arabic{section}.\arabic{fakt}}

\addto\captionsindonesian{%
  \renewcommand{\contentsname}{Daftar Isi}%
  \renewcommand{\figurename}{Gambar}%
  \renewcommand{\listfigurename}{Daftar Gambar}%
}
\hypersetup{pdflang={id-ID},bookmarksopen=true}
\emergencystretch=1em
\ifdefined\pdfinfoomitdate\pdfinfoomitdate=1\fi
\ifdefined\pdftrailerid\pdftrailerid{}\fi
\ifdefined\pdfsuppressptexinfo\pdfsuppressptexinfo=15\fi

\title{Geometri Diferensial dan Manifold Mulus\\\large Pembaca kumulatif hingga Unit 10}
\author{Holger Brenner\\Terjemahan Bahasa Indonesia independen}
\date{Sumber: \textit{Differentialgeometrie (Osnabrück 2023)}}

\begin{document}
\hypersetup{pageanchor=false}
\maketitle
\hypersetup{pageanchor=true}
\frontmatter
\chapter*{Tentang edisi ini}
Pembaca kumulatif ini menerjemahkan Kuliah 1--10 dan Lembar Kerja 1--10 dari kursus Holger Brenner di Wikiversity berbahasa Jerman. Teks sumber digunakan berdasarkan CC BY-SA 4.0. Terjemahan ini merupakan karya independen dan bukan edisi resmi atau dukungan dari penulis maupun Wikiversity. Setiap gambar mengikuti status hak atau lisensi berkasnya sendiri, yang dicatat pada daftar gambar dan manifest media.

Sumber permanen dan hash revisi dicatat dalam berkas kontrol edisi ini. Rumus, urutan, latihan, penanda solusi, dan atribusi media dipertahankan; ID mesin hanya merupakan lapisan tambahan. Bagian solusi hanya memuat solusi yang benar-benar disediakan oleh sumber.

Proses terjemahan dan produksi edisi dibantu oleh OpenAI Codex gpt-5.6-sol, Ultra, di bawah arahan pengguna. Kredit penulis, sumber, dan kontributor manusia tetap dipertahankan sebagaimana dicatat dalam edisi dan manifestnya.
\tableofcontents

\mainmatter
\part{Unit 1}
\chapter[Kuliah 1]{Kuliah 1: Analisis dan Geometri}
\input{generated/lecture01.id.build.tex}

\chapter{Lembar Kerja 1}
\setcounter{section}{1}

  \zwischenueberschrift{Soal latihan}

\inputaufgabegibtloesung
{}
{

Misalkan

\mavergleichskette
{\vergleichskette
{ G
}
{ \subseteq }{ \R^n
}
{  }{
}
{    }{
}
{   }{
}
}
{}{}{}
\definitionsverweis {terbuka}{}{}
dan
\maabbdisp {\varphi} { G } { \R^m
} {}
sebuah
\definitionsverweis {pemetaan terdiferensialkan secara kontinu}{}{,}
yang di titik

\mavergleichskette
{\vergleichskette
{ P
}
{ \in }{ G
}
{  }{
}
{    }{
}
{   }{
}
}
{}{}{}
mempunyai
\definitionsverweis {diferensial total}{}{}
yang surjektif. Misalkan

\mavergleichskette
{\vergleichskette
{ v
}
{ \in }{ T_PZ
}
{  }{
}
{    }{
}
{   }{
}
}
{}{}{}
sebuah vektor di
\definitionsverweis {ruang tangen pada serat}{}{}
$Z$ dari $\varphi$ yang melalui $P$. Buktikan bahwa terdapat
\definitionsverweis {kurva terdiferensialkan secara kontinu}{}{}
\maabbdisp {\gamma} { ]- \epsilon, \epsilon[ } { Z
} {}
\zusatzklammer {untuk suatu

\mavergleichskettek
{\vergleichskettek
{ \epsilon
}
{ > }{ 0
}
{  }{
}
{    }{
}
{   }{
}
}
{}{}{}
yang sesuai} {} {}
dengan

\mavergleichskette
{\vergleichskette
{ \gamma(0)
}
{ = }{ P
}
{  }{
}
{    }{
}
{   }{
}
}
{}{}{}
dan

\mavergleichskettedisp
{\vergleichskette
{ \gamma'(0)
}
{ =} { v
}
{ } {
}
{ } {
}
{ } {
}
}
{}{}{.}

}
{} {}

\inputaufgabe
{}
{

Misalkan

\mavergleichskette
{\vergleichskette
{ W
}
{ \subseteq }{ \R^n
}
{  }{
}
{    }{
}
{   }{
}
}
{}{}{}
\definitionsverweis {terbuka}{}{,}
\maabb {h} { W } { \R
} {}
sebuah
\definitionsverweis {fungsi terdiferensialkan secara kontinu}{}{}
dan

\mavergleichskette
{\vergleichskette
{ Y
}
{ = }{ h^{-1}(0)
}
{  }{
}
{    }{
}
{   }{
}
}
{}{}{}
merupakan
\definitionsverweis {serat}{}{}
di atas $0$, dengan $h$
\definitionsverweis {reguler}{}{}
di setiap titik pada $Y$. Misalkan
\maabbdisp {g} { W } {\R
} {}
sebuah fungsi terdiferensialkan secara kontinu yang tidak mempunyai titik nol pada $Y$. Buktikan bahwa $Y$ juga dapat diperoleh sebagai serat dari

\mavergleichskettedisp
{\vergleichskette
{ \tilde{h}
}
{ =} { g h
}
{ } {
}
{ } {
}
{ } {
}
}
{}{}{,}
bahwa
\definitionsverweis {gradien}{}{}
dari $h$ dan $\tilde{h}$ merupakan kelipatan skalar satu sama lain, dan bahwa
\definitionsverweis {ruang-ruang tangen}{}{}
hanya bergantung pada $Y$.

}
{} {}

\inputaufgabe
{}
{

Misalkan

\mavergleichskette
{\vergleichskette
{ W
}
{ \subseteq }{ \R^n
}
{  }{
}
{    }{
}
{   }{
}
}
{}{}{}
\definitionsverweis {terbuka}{}{,}
\maabb {h} { W } { \R
} {}
sebuah
\definitionsverweis {fungsi terdiferensialkan secara kontinu}{}{}
dan

\mavergleichskette
{\vergleichskette
{ Y
}
{ = }{ h^{-1}(c)
}
{  }{
}
{    }{
}
{   }{
}
}
{}{}{}
merupakan
\definitionsverweis {serat}{}{}
di atas

\mavergleichskette
{\vergleichskette
{ c
}
{ \in }{ \R
}
{  }{
}
{    }{
}
{   }{
}
}
{}{}{,}
dengan $h$
\definitionsverweis {reguler}{}{}
di setiap titik pada $Y$. Misalkan
\maabbdisp {L} {\R^n } { \R^n
} {}
sebuah
\definitionsverweis {isomorfisme linear}{}{,}

\mavergleichskette
{\vergleichskette
{ W'
}
{ = }{ L^{-1}(W)
}
{  }{
}
{    }{
}
{   }{
}
}
{}{}{}
dan

\mavergleichskettedisp
{\vergleichskette
{ Z
}
{ =} { L^{-1}(Y)
}
{ =} { (h \circ L)^{-1}(c)
}
{ } {
}
{ } {
}
}
{}{}{.}
Buktikan bahwa konsep
\definitionsverweis {ruang tangen}{}{,}
medan vektor tangensial, dan solusi persamaan diferensial tangensial yang terkait pada $Y$ dan $Z$ secara umum saling bersesuaian
\zusatzklammer {yakni dapat dipetakan satu sama lain dengan bantuan $L$} {} {.}
Bagaimana dengan gradiennya?

}
{} {}

\inputaufgabe
{}
{

Misalkan

\mavergleichskette
{\vergleichskette
{ Y
}
{ = }{ h^{-1}(c)
}
{  }{
}
{    }{
}
{   }{
}
}
{}{}{}
sebuah
\definitionsverweis {hipermuka terdiferensialkan}{}{}
dan misalkan
\maabbdisp {f} { U} { \R
} {}
sebuah fungsi terdiferensialkan secara kontinu yang didefinisikan pada suatu lingkungan terbuka

\mavergleichskette
{\vergleichskette
{ Y
}
{ \subseteq }{ U
}
{  }{
}
{    }{
}
{   }{
}
}
{}{}{}
Andaikan $f$ di

\mavergleichskette
{\vergleichskette
{ P
}
{ \in }{ Y
}
{  }{
}
{    }{
}
{   }{
}
}
{}{}{}
mempunyai
\definitionsverweis {ekstremum lokal}{}{}
pada $Y$. Buktikan bahwa
\definitionsverweis {komponen tangensial}{}{}
dari $\operatorname{Grad} \,  f  (  P )$ sama dengan $0$.

}
{} {}

\inputaufgabe
{}
{

Misalkan

\mavergleichskette
{\vergleichskette
{ Y
}
{ \subseteq }{ \R^3
}
{  }{
}
{    }{
}
{   }{
}
}
{}{}{}
adalah permukaan bola satuan. Tentukan
\definitionsverweis {dekomposisi ortogonal}{}{}
vektor

\mavergleichskette
{\vergleichskette
{ \begin{pmatrix}  6  \\ 7 \\  -3   \end{pmatrix}
}
{ \in }{ \R^3
}
{  }{
}
{    }{
}
{   }{
}
}
{}{}{}
menjadi komponen tangensial dan komponen ortogonal di titik

\mavergleichskette
{\vergleichskette
{ \begin{pmatrix}  -  { \frac{   1  }{  2 } }  \\ { \frac{   1  }{  2 } } \\  { \frac{   1  }{  \sqrt{2}  } }   \end{pmatrix}
}
{ \in }{ Y
}
{  }{
}
{    }{
}
{   }{
}
}
{}{}{.}

}
{} {}

\inputaufgabe
{}
{

Misalkan

\mavergleichskette
{\vergleichskette
{ Y
}
{ \subseteq }{ \R^3
}
{  }{
}
{    }{
}
{   }{
}
}
{}{}{}
adalah
\definitionsverweis {permukaan bola satuan}{}{.}
Misalkan
\maabbdisp {\gamma} {[0, 2 \pi]} { Y
} {}
merupakan gerak melingkar beraturan yang mengelilingi permukaan bola pada ketinggian ${ \frac{   1  }{  2 } }$.
\aufzaehlungzwei {Parametrisasikan $\gamma$.
} {Tentukan
\definitionsverweis {percepatan tangensial}{}{}
dari $\gamma$ pada setiap saat

\mavergleichskette
{\vergleichskette
{ t
}
{ \in }{ [0, 2 \pi]
}
{  }{
}
{    }{
}
{   }{
}
}
{}{}{.}
}

}
{} {}

\inputaufgabe
{}
{

Misalkan $Y$ adalah
\definitionsverweis {grafik}{}{}
dari
\maabbeledisp {f} {\R^2} { \R
} {(x,y)} { x^2+y^3
} {,}
dan misalkan
\maabbdisp {\gamma} {\R} { Y
} {}
adalah lintasan pada grafik tersebut yang berada di atas kurva dasar
\mathl{t \mapsto (t,t)}{.} Tentukan, untuk setiap

\mavergleichskette
{\vergleichskette
{ t
}
{ \in }{ \R
}
{  }{
}
{    }{
}
{   }{
}
}
{}{}{,}
\definitionsverweis {percepatan tangensial}{}{}
dari $\gamma$ di $\gamma(t)$.

}
{} {}

\inputaufgabe
{}
{

Misalkan

\mavergleichskette
{\vergleichskette
{ Y
}
{ \subseteq }{ \R^n
}
{  }{
}
{    }{
}
{   }{
}
}
{}{}{}
sebuah hiperbidang, yaitu serat dari fungsi linear

\mavergleichskettedisp
{\vergleichskette
{ h(x_1 , \ldots , x_n)
}
{ =} { a_1x_1 + \cdots + a_nx_n
}
{ } {
}
{ } {
}
{ } {
}
}
{}{}{}
di atas

\mavergleichskette
{\vergleichskette
{ c
}
{ \in }{ \R
}
{  }{
}
{    }{
}
{   }{
}
}
{}{}{}
dengan

\mavergleichskette
{\vergleichskette
{ a_i
}
{ \neq }{ 0
}
{  }{
}
{    }{
}
{   }{
}
}
{}{}{}
untuk setidaknya satu indeks $i$. Tafsirkan pengertian
\definitionsverweis {diferensial total}{}{,}
\definitionsverweis {gradien}{}{,}
\definitionsverweis {ruang tangen}{}{,}
dekomposisi menjadi komponen tangensial dan normal,
\definitionsverweis {percepatan tangensial}{}{,}
serta
\definitionsverweis {kurva geodesik}{}{}
dalam kasus khusus ini.

}
{} {}

\inputaufgabe
{}
{

Misalkan

\mavergleichskette
{\vergleichskette
{ Y
}
{ = }{ h^{-1}(c)
}
{  }{
}
{    }{
}
{   }{
}
}
{}{}{}
sebuah
\definitionsverweis {hipermuka terdiferensialkan}{}{,}
yang memuat garis
\mathl{P+ \R v}{.} Buktikan bahwa garis yang dilalui secara seragam tersebut merupakan
\definitionsverweis {kurva geodesik}{}{}
pada $Y$.

}
{} {}

Dua titik

\mathbed {P,Q \in K} {}
 {P \neq Q} {}
 {} {} {} {,}
disebut \stichwort {antipodal} {,} jika garis yang menghubungkannya melalui pusat bola.

\inputaufgabe
{}
{

Misalkan

\mavergleichskette
{\vergleichskette
{ K
}
{ \subseteq }{ \R^3
}
{  }{
}
{    }{
}
{   }{
}
}
{}{}{}
sebuah permukaan bola. Buktikan bahwa setiap dua
\definitionsverweis {titik yang tidak antipodal}{}{}

\mathbed {P,Q \in K} {}
 {P \neq Q} {}
 {} {} {} {,}
terletak pada tepat satu
\definitionsverweis {lingkaran besar}{}{}
dari $K$.

}
{} {}

\bild{ \begin{center}
\includegraphics[width=5.5cm]{ \bildeinlesungsvg {Planned flight map of the Oiseau Blanc} {svg} }
\end{center}
\bildtext {Mengapa pesawat terbang menempuh lintasan melengkung?} }

\bildlizenz { Planned flight map of the Oiseau Blanc.svg } {} {Pethrus} {Commons} {CC BY-SA 3.0} {}

\inputaufgabe
{}
{

Sebuah pesawat akan terbang dari Osnabrück menuju suatu tempat tujuan di belahan Bumi selatan. Mungkinkah rutenya lebih pendek jika pesawat itu mula-mula terbang ke arah utara?

}
{} {}

\inputaufgabe
{}
{

Lucy Sonnenschein merasa bahwa hari-hari terlalu singkat. Karena itu, setiap hari dan selama masih terang, ia memutuskan terbang ke arah barat melintasi satu zona waktu agar setiap harinya berlangsung $25$, bukan $24$, jam.
\aufzaehlungfuenf{Dapatkah Lucy Sonnenschein meningkatkan proporsi jam terang dalam hidupnya dengan strategi ini?
}{Dapatkah Lucy Sonnenschein memperpanjang hidupnya dengan strategi ini?
}{Sekarang Lucy mempunyai teman di seluruh dunia. Setelah berapa hari ia bertemu dengan mereka lagi?
}{Apa yang menyebabkan Lucy mengalami kehilangan waktu yang mengimbangi tambahan waktu hariannya?
}{Apakah hal ini berkaitan dengan teori relativitas?
}

}
{} {}

\inputaufgabe
{}
{

Berikan contoh kurva dua kali terdiferensialkan yang bergerak pada suatu
\definitionsverweis {lingkaran besar}{}{}
dari permukaan bola satuan, tetapi bukan
\definitionsverweis {kurva geodesik}{}{.}

}
{} {}

\inputaufgabe
{}
{

Sebuah \stichwort {segitiga geodesik} {} pada permukaan bola

\mavergleichskettedisp
{\vergleichskette
{ S^2
}
{ =} { \left\{ (x,y,z) \in \R^3  \mid   x^2+y^2+z^2 = 1        \right\}
}
{ } {
}
{ } {
}
{ } {
}
}
{}{}{}
terdiri atas tiga titik berbeda yang tidak terletak pada satu lingkaran besar,

\mavergleichskette
{\vergleichskette
{ A,B,C
}
{ \in }{ S^2
}
{  }{
}
{    }{
}
{   }{
}
}
{}{}{,}
dengan setiap pasangan titik dihubungkan oleh busur yang lebih pendek dari suatu
\definitionsverweis {lingkaran besar}{}{}
\zusatzklammer {setiap pasangan titik tidak antipodal} {} {.}
Buktikan bahwa jumlah sudut dalam segitiga itu lebih besar daripada $\pi$.

}
{} {}

  \zwischenueberschrift{Soal untuk dikumpulkan}

\inputaufgabe
{3}
{

Misalkan $Y$ adalah hipermuka yang diberikan oleh syarat

\mavergleichskettedisp
{\vergleichskette
{ 2x^2 +3y^2+5z^2
}
{ =} { 10
}
{ } {
}
{ } {
}
{ } {
}
}
{}{}{}
di $\R^3$. Tentukan dekomposisi ortogonal vektor

\mavergleichskette
{\vergleichskette
{ \begin{pmatrix}  4  \\ 2 \\  -5   \end{pmatrix}
}
{ \in }{ \R^3
}
{  }{
}
{    }{
}
{   }{
}
}
{}{}{}
menjadi komponen tangensial dan komponen ortogonal di titik

\mavergleichskette
{\vergleichskette
{ \begin{pmatrix}  -1 \\ -1\\  1   \end{pmatrix}
}
{ \in }{ Y
}
{  }{
}
{    }{
}
{   }{
}
}
{}{}{.}

}
{} {}

\inputaufgabe
{4}
{

Misalkan $Y$ adalah
\definitionsverweis {grafik}{}{}
dari
\maabbeledisp {f} {\R^2} { \R
} {(x,y)} { x^2+y^2
} {,}
dan misalkan
\maabbdisp {\gamma} {\R} { Y
} {}
adalah lintasan pada grafik tersebut yang berada di atas kurva dasar
\mathl{x \mapsto (x,1)}{.} Tentukan, untuk setiap

\mavergleichskette
{\vergleichskette
{ x
}
{ \in }{ \R
}
{  }{
}
{    }{
}
{   }{
}
}
{}{}{,}
\definitionsverweis {percepatan tangensial}{}{}
dari $\gamma$ di $\gamma(x)$.

}
{} {}

\inputaufgabe
{4}
{

Pada 11 April 2023 pukul 17.45, secara tiba-tiba dan mengejutkan semua orang, rotasi Bumi pada porosnya dihentikan. Pada saat yang sama, hambatan udara dan semua kerugian akibat gesekan ditiadakan. Semua manusia, hewan, dan benda yang tidak terpasang tetap atau tidak sempat berpegangan akan meneruskan gerak sesaatnya sesuai hukum inersia. Gravitasi tetap bekerja sehingga, untungnya, mereka tidak terlepas ke angkasa. Kapan penghapus papan tulis dari ruang kuliah di Osnabrück
\zusatzklammer {penghapus itu terbang keluar melalui jendela} {} {}
untuk pertama kalinya kembali ke posisi awalnya?

}
{} {}

\inputaufgabe
{4 (3+1)}
{

Misalkan

\mavergleichskette
{\vergleichskette
{ W
}
{ \subseteq }{ \R^n
}
{  }{
}
{    }{
}
{   }{
}
}
{}{}{}
\definitionsverweis {terbuka}{}{,}
\maabb {h} { W } { \R
} {}
sebuah
\definitionsverweis {fungsi terdiferensialkan secara kontinu}{}{}
dan

\mavergleichskette
{\vergleichskette
{ Y
}
{ = }{ h^{-1}(c)
}
{  }{
}
{    }{
}
{   }{
}
}
{}{}{}
merupakan
\definitionsverweis {serat}{}{}
di atas

\mavergleichskette
{\vergleichskette
{ c
}
{ \in }{ \R
}
{  }{
}
{    }{
}
{   }{
}
}
{}{}{,}
dengan $h$
\definitionsverweis {reguler}{}{}
di setiap titik pada $Y$. Misalkan
\maabbdisp {\gamma} { I } { Y
} {}
sebuah kurva dua kali
\definitionsverweis {terdiferensialkan secara kontinu}{}{.}
\aufzaehlungzwei {Buktikan bahwa jika $\gamma$ merupakan
\definitionsverweis {kurva geodesik}{}{,}
maka
\definitionsverweis {norma}{}{}
dari $\gamma'$ adalah konstan.
} {Berikan contoh kurva dengan norma kecepatan konstan yang bukan kurva geodesik.
}

}
{} {}

\inputaufgabe
{3}
{

Misalkan
\maabb {h,g} {\R^n } { \R
} {}
adalah
\definitionsverweis {fungsi-fungsi terdiferensialkan secara kontinu}{}{}
dan

\mavergleichskette
{\vergleichskette
{ Y
}
{ = }{ h^{-1}(c)
}
{  }{
}
{    }{
}
{   }{
}
}
{}{}{}
serta

\mavergleichskette
{\vergleichskette
{ Z
}
{ = }{ g^{-1}(d)
}
{  }{
}
{    }{
}
{   }{
}
}
{}{}{}
adalah
\definitionsverweis {hipermuka terdiferensialkan}{}{}
yang terkait, dengan $h$ di setiap titik pada $Y$ dan $g$ di setiap titik pada $Z$ bersifat
\definitionsverweis {reguler}{}{.}
Misalkan

\mavergleichskette
{\vergleichskette
{ T
}
{ = }{ Y \cap Z
}
{  }{
}
{    }{
}
{   }{
}
}
{}{}{}
dan misalkan
\maabbdisp {\gamma} { I } { T
} {}
sebuah kurva dua kali terdiferensialkan secara kontinu yang, jika dipandang sebagai kurva pada $Y$, merupakan
\definitionsverweis {kurva geodesik}{}{.}
Apakah $\gamma$ juga merupakan kurva geodesik pada $Z$?

}
{} {}


\section*{Solusi yang disediakan oleh sumber}
\addcontentsline{toc}{section}{Solusi yang disediakan oleh sumber}
% Authority: German Wikiversity pageid 131353, revid 1111802, 2026-08-15T16:41:12Z.
% Exact UTF-8 witness SHA-256: b07aac01b6f803481fab9b5331e19480fb9bf0058597159b6152c65c7bc955ce.

\zwischenueberschrift{Solusi untuk Soal 1}

Berlaku

\mavergleichskettedisp
{\vergleichskette
{ v
}
{ \in }{ T_PZ
}
{ = }{ \ker \left((D\varphi)_P\right)
}
{ } {
}
{ } {
}
}
{}{}{.}
Menurut
\definitionsverweis {teorema fungsi implisit}{}{,}
terdapat himpunan terbuka

\mathbed {P \in W} {}
 {W \subseteq G} {}
 {} {} {} {,}
himpunan terbuka

\mavergleichskette
{\vergleichskette
{ V
}
{ \subseteq }{ \R^{n-m}
}
{ } {
}
{ } {
}
{ } {
}
}
{}{}{,}
dan pemetaan terdiferensialkan secara kontinu
\maabbdisp {\psi} { V } { W } {}
sedemikian sehingga

\mavergleichskette
{\vergleichskette
{ \psi(V)
}
{ \subseteq }{ Z \cap W
}
{ } {
}
{ } {
}
{ } {
}
}
{}{}{}
dan $\psi$ menginduksi
\definitionsverweis {bijeksi}{}{}
\maabbdisp {\psi} { V } { Z \cap W } {}.
Misalkan

\mavergleichskette
{\vergleichskette
{ Q
}
{ \in }{ V
}
{ } {
}
{ } {
}
{ } {
}
}
{}{}{}
adalah titik yang memenuhi

\mavergleichskettedisp
{\vergleichskette
{ \psi(Q)
}
{ =} { P
}
{ } {
}
{ } {
}
{ } {
}
}
{}{}{.}
Pemetaan $\psi$ di setiap titik

\mavergleichskette
{\vergleichskette
{ Q
}
{ \in }{ V
}
{ } {
}
{ } {
}
{ } {
}
}
{}{}{}
bersifat
\definitionsverweis {reguler}{}{}
dan
\definitionsverweis {diferensial total}{}{}
dari $\psi$ memenuhi

\mavergleichskettedisp
{\vergleichskette
{ (D\varphi)_P \circ (D\psi)_Q
}
{ =} { 0
}
{ } {
}
{ } {
}
{ } {
}
}
{}{}{.}
Dengan demikian,

\mavergleichskettedisp
{\vergleichskette
{ \operatorname{im} \left((D\psi)_Q\right)
}
{ =} { \ker \left((D\varphi)_P\right)
}
{ =} { T_PZ
}
{ } {
}
{ } {
}
}
{}{}{.}
Karena $\psi$ reguler di $Q$, pemetaan
\maabbdisp {(D\psi)_Q} { \R^{n-m} } { \R^n } {}
bersifat injektif, sedangkan pemetaan
\maabbdisp {(D\psi)_Q} { \R^{n-m} } { T_PZ } {}
bersifat bijektif. Misalkan

\mavergleichskette
{\vergleichskette
{ u
}
{ \in }{ \R^{n-m}
}
{ } {
}
{ } {
}
{ } {
}
}
{}{}{}
adalah prapeta dari $v$, dan misalkan
\maabbeledisp {\theta} { ]-\epsilon,\epsilon[ } { \R^{n-m} } {t} {Q+tu} {,}
dengan $\epsilon$ dipilih cukup kecil agar seluruh citra terletak di dalam $V$. Maka
\maabbdisp {\gamma=\psi\circ\theta} { ]-\epsilon,\epsilon[ } { \R^n } {}
mempunyai sifat

\mavergleichskettedisp
{\vergleichskette
{ \gamma(0)
}
{ =} { \psi(\theta(0))
}
{ =} { \psi(Q)
}
{ =} { P
}
{ } {
}
}
{}{}{}
dan

\mavergleichskettedisp
{\vergleichskette
{ \gamma'(0)
}
{ =} { (D\psi)_Q(\theta'(0))
}
{ =} { (D\psi)_Q(u)
}
{ =} { v
}
{ } {
}
}
{}{}{.}


\part{Unit 2}
\chapter[Kuliah 2]{Kuliah 2: Permukaan Putar, Medan Normal, dan Pemetaan Gauss}
\input{generated/lecture02.id.build.tex}

\chapter{Lembar Kerja 2}
\setcounter{section}{2}

  \zwischenueberschrift{Soal latihan}

\inputaufgabegibtloesung
{}
{

Misalkan
\maabbdisp {f,g} {\R} { \R^n
} {}
\definitionsverweis {kurva-kurva terdiferensialkan}{}{.}
Hitung
\definitionsverweis {turunan}{}{}
fungsi
\maabbeledisp {} {\R} {\R
} { t } { \left\langle f(t) , g(t)  \right\rangle
} {.}
Nyatakan hasilnya menggunakan hasil kali dalam.

}
{} {}

\inputaufgabegibtloesung
{}
{

Misalkan
\maabb {h} { I } { V
} {}
sebuah kurva yang terdiferensialkan dua kali dalam
\definitionsverweis {ruang vektor Euklides}{}{}
$V$. Buktikan bahwa jika

\mavergleichskette
{\vergleichskette
{h'(t)
}
{ \neq }{ 0
}
{  }{
}
{    }{
}
{   }{
}
}
{}{}{,}
maka berlaku persamaan

\mavergleichskettedisp
{\vergleichskette
{ \Vert { h'(t)} \Vert'
}
{ =} { { \frac{   \left\langle h'(t) , h^{\prime \prime} (t)  \right\rangle  }{  \left\langle h'(t) , h'(t)  \right\rangle  } } \Vert { h'(t)} \Vert
}
{ } {
}
{ } {
}
{ } {
}
}
{}{}{.}

}
{} {}

Sebuah
\definitionsverweis {kurva terdiferensialkan}{}{}
\maabbeledisp {f} { [a,b] } {\R^n
} { t } {f(t) = \left( f_1(t)  , \,   \ldots  , \,  f_n(t)                 \right)
} {,}
disebut
\definitionswort {berparametrisasi panjang busur}{,}
jika

\mavergleichskette
{\vergleichskette
{ f_1'(t)^2 + \cdots +  f_n'(t)^2
}
{ = }{ 1
}
{  }{
}
{    }{
}
{   }{
}
}
{}{}{}
untuk setiap $t$.

\inputaufgabe
{}
{

Misalkan
\maabb {\gamma} { I } {\R^n
} {}
sebuah kurva yang dua kali
\definitionsverweis {terdiferensialkan secara kontinu}{}{}
dan
\definitionsverweis {berparametrisasi panjang busur}{}{.}
Buktikan bahwa
\mathkor {} {\gamma'(t)} {dan} {\gamma^{\prime \prime}(t)} {}
saling tegak lurus.

}
{} {}

\inputaufgabe
{}
{

Misalkan

\mavergleichskette
{\vergleichskette
{ P
}
{ \in }{ \R^n
}
{  }{
}
{    }{
}
{   }{
}
}
{}{}{}
sebuah titik dan misalkan

\mavergleichskette
{\vergleichskette
{ I
}
{ = }{ {]{-1},1[}
}
{  }{
}
{    }{
}
{   }{
}
}
{}{}{.}
Kita tinjau himpunan

\mavergleichskettedisp
{\vergleichskette
{ M
}
{ =} { \left\{ f:I \rightarrow \R^n   \mid  f \text{ terdiferensialkan} , \,  f(0) = P       \right\}
}
{ } {
}
{ } {
}
{ } {
}
}
{}{}{.}
Dua
\definitionsverweis {kurva}{}{}

\mavergleichskette
{\vergleichskette
{ f,g
}
{ \in }{ M
}
{  }{
}
{    }{
}
{   }{
}
}
{}{}{}
disebut \stichwort {ekuivalen secara tangensial} {,} jika

\mavergleichskettedisp
{\vergleichskette
{ f'(0)
}
{ =} { g'(0)
}
{ } {
}
{ } {
}
{ } {
}
}
{}{}{.}

a) Buktikan bahwa ini merupakan suatu
\definitionsverweis {relasi ekuivalensi}{}{.}

b) Temukan wakil paling sederhana bagi setiap
\definitionsverweis {kelas ekuivalensi}{}{.}

c) Berikan satu wakil lain untuk setiap kelas.

d) Deskripsikan himpunan kelas-kelas ekuivalensi tersebut
\zusatzklammer {yakni
\definitionsverweis {himpunan hasil bagi}{}{}} {} {.}

}
{} {}

\inputaufgabe
{}
{

Dengan syarat
\zusatzklammer {reguler di setiap titik} {} {,}
berikan sebuah contoh
\definitionsverweis {hipermuka terdiferensialkan}{}{,}
yang tidak
\definitionsverweis {terhubung}{}{.}

}
{} {}

\inputaufgabe
{}
{

Misalkan

\mavergleichskette
{\vergleichskette
{ V
}
{ \subseteq }{ \R^{n-1}
}
{  }{
}
{    }{
}
{   }{
}
}
{}{}{}
sebuah
\definitionsverweis {himpunan terbuka}{}{}
dan misalkan
\maabbdisp {f} { V } {\R
} {}
sebuah fungsi yang
\definitionsverweis {terdiferensialkan secara kontinu}{}{.}
Kita tinjau
\definitionsverweis {grafik}{}{}
$Y$ dari $f$ sebagai hipermuka terdiferensialkan di $\R^n$ dalam pengertian
Contoh 1.2.
Tentukan
\definitionsverweis {medan-medan normal satuan}{}{}
dalam kasus ini.

}
{} {}

\inputaufgabegibtloesung
{}
{

Misalkan
\mathkor {} {r} {dan} {R} {}
bilangan real positif yang memenuhi

\mavergleichskette
{\vergleichskette
{ 0
}
{ < }{ r
}
{ < }{ R
}
{    }{
}
{   }{
}
}
{}{}{.}
Tentukan suatu
\definitionsverweis {medan normal satuan}{}{}
untuk torus
\zusatzklammer {tertanam} {} {}

\mavergleichskettedisp
{\vergleichskette
{ T
}
{ =} { \left\{ (x,y,z) \in \R^3  \mid   \left(  \sqrt{x^2+y^2} -R  \right)^2 +z^2  = r^2         \right\}
}
{ } {
}
{ } {
}
{ } {
}
}
{}{}{.}

}
{} {}

\inputaufgabe
{}
{

Misalkan
\maabbdisp {f} { I } {\R_+
} {}
sebuah
\definitionsverweis {fungsi terdiferensialkan}{}{}
dan misalkan $Y$ adalah permukaan putar yang terkait dengan grafik $f$, dipandang sebagai hipermuka terdiferensialkan di $\R^3$ dalam pengertian
Lemma 2.1.
Tentukan
\definitionsverweis {medan-medan normal satuan}{}{}
dalam kasus ini.

}
{} {}

\inputaufgabe
{}
{

Buktikan bahwa
\definitionsverweis {pemetaan Gauss}{}{}
pada sebuah hiperbidang bersifat konstan. Apakah pernyataan sebaliknya juga berlaku?

}
{} {}

\inputaufgabe
{}
{

Buktikan bahwa
\definitionsverweis {pemetaan Gauss}{}{}
pada sebuah permukaan bola di $\R^n$ yang berpusat di titik asal, untuk setiap pilihan orientasi, diinduksi oleh homoteti skalar pada $\R^n$ (dengan faktor positif untuk orientasi keluar dan faktor negatif untuk orientasi masuk).

}
{} {}

\inputaufgabe
{}
{

Misalkan

\mavergleichskette
{\vergleichskette
{ Y
}
{ \subseteq }{ \R^n
}
{  }{
}
{    }{
}
{   }{
}
}
{}{}{}
sebuah
\definitionsverweis {hipermuka terdiferensialkan}{}{}
dengan
\definitionsverweis {medan normal satuan}{}{}
$N$ dan medan normal satuan yang berlawanan, $-N$. Buktikan bahwa kedua
\definitionsverweis {pemetaan Gauss}{}{}
yang terkait saling berhubungan melalui
\definitionsverweis {pemetaan antipodal}{}{}
\maabbeledisp {} {S^{n-1}} {S^{n-1}
} { P } { -P
} {.}

}
{} {}

\inputaufgabegibtloesung
{}
{

Misalkan

\mavergleichskette
{\vergleichskette
{ \alpha,\beta
}
{ \in }{ \R_+
}
{  }{
}
{    }{
}
{   }{
}
}
{}{}{}
dan misalkan

\mavergleichskettedisp
{\vergleichskette
{ Y
}
{ =} { \left\{ (x,y)\in\R^2  \mid   \alpha x^2+\beta y^2 = 1        \right\}
}
{ } {
}
{ } {
}
{ } {
}
}
{}{}{.}
Orientasikan $Y$ dengan medan normal satuan yang diperoleh dari gradien keluar fungsi
$ (x,y) \longmapsto \alpha x^2+\beta y^2 $.
\aufzaehlungzwei {Tentukan
\definitionsverweis {pemetaan Gauss}{}{}
untuk $Y$.
} { Berikan secara eksplisit suatu
\definitionsverweis {pemetaan invers}{}{}
untuk pemetaan Gauss pada $Y$.
}

}
{} {}

\inputaufgabegibtloesung
{}
{

Misalkan

\mavergleichskette
{\vergleichskette
{ \alpha,\beta
}
{ \in }{ \R_+
}
{  }{
}
{    }{
}
{   }{
}
}
{}{}{}
dan misalkan

\mavergleichskettedisp
{\vergleichskette
{ Y
}
{ =} { \left\{ (x,y)\in\R^2  \mid   \alpha x^2+\beta y^2 = 1        \right\}
}
{ } {
}
{ } {
}
{ } {
}
}
{}{}{.}
Buktikan bahwa
\definitionsverweis {pemetaan Gauss}{}{}
untuk $Y$ bersifat bijektif.

}
{} {}

\inputaufgabe
{}
{

Berikan sebuah contoh
\definitionsverweis {hipermuka terdiferensialkan}{}{}

\mavergleichskette
{\vergleichskette
{ Y
}
{ \subseteq }{ \R^2
}
{  }{
}
{    }{
}
{   }{
}
}
{}{}{}
sedemikian sehingga
\definitionsverweis {pemetaan Gauss}{}{}
bersifat surjektif dan terdapat titik-titik pada $S^1$ yang dicapai tak berhingga kali.

}
{} {}

\inputaufgabe
{}
{

Misalkan

\mavergleichskette
{\vergleichskette
{ W
}
{ \subseteq }{ \R^n
}
{  }{
}
{    }{
}
{   }{
}
}
{}{}{}
\definitionsverweis {terbuka}{}{,}
\maabb {h} { W } { \R
} {}
sebuah
\definitionsverweis {fungsi yang terdiferensialkan secara kontinu}{}{}
dan

\mavergleichskette
{\vergleichskette
{ Y
}
{ = }{ h^{-1}(c)
}
{  }{
}
{    }{
}
{   }{
}
}
{}{}{}
adalah
\definitionsverweis {serat}{}{}
di atas

\mavergleichskette
{\vergleichskette
{ c
}
{ \in }{ \R
}
{  }{
}
{    }{
}
{   }{
}
}
{}{}{,}
dengan $h$
\definitionsverweis {reguler}{}{}
di setiap titik pada $Y$. Misalkan
\maabbdisp {L} {\R^n } { \R^n
} {}
sebuah
\definitionsverweis {isometri linear}{}{,}

\mavergleichskette
{\vergleichskette
{ W'
}
{ = }{ L^{-1}(W)
}
{  }{
}
{    }{
}
{   }{
}
}
{}{}{}
dan

\mavergleichskettedisp
{\vergleichskette
{ Z
}
{ =} { L^{-1}(Y)
}
{ =} { (h \circ L)^{-1}(c)
}
{ } {
}
{ } {
}
}
{}{}{.}
Buktikan bahwa konsep-konsep kurva geodesik, medan normal, medan normal satuan, dan pemetaan Gauss pada $Y$ dan $Z$ saling bersesuaian
\zusatzklammer {yakni dapat dipetakan satu sama lain dengan bantuan $L$} {} {.}

}
{} {}

  \zwischenueberschrift{Soal untuk dikumpulkan}

\inputaufgabe
{3}
{

Tentukan citra parabola standar berorientasi ke atas di bawah
\definitionsverweis {pemetaan Gauss}{}{.}

}
{} {}

\inputaufgabe
{4}
{

Misalkan

\mavergleichskette
{\vergleichskette
{ W
}
{ = }{ \R^3 \setminus \{ (0,0,0) \}
}
{  }{
}
{    }{
}
{   }{
}
}
{}{}{}
dan misalkan

\mavergleichskette
{\vergleichskette
{ Y
}
{ \subseteq }{ W
}
{  }{
}
{    }{
}
{   }{
}
}
{}{}{}
adalah
\definitionsverweis {himpunan nol}{}{}
dari

\mavergleichskettedisp
{\vergleichskette
{ h(x,y,z)
}
{ =} { x^2+y^2-z^2
}
{ } {
}
{ } {
}
{ } {
}
}
{}{}{.}
Tentukan citra $Y$
\zusatzklammer {dengan orientasi yang ditentukan oleh gradien $h$} {} {}
di bawah
\definitionsverweis {pemetaan Gauss}{}{.}

}
{} {}

\inputaufgabe
{5}
{

Misalkan

\mavergleichskette
{\vergleichskette
{ \alpha,\beta, \gamma
}
{ \in }{ \R_+
}
{  }{
}
{    }{
}
{   }{
}
}
{}{}{}
dan misalkan

\mavergleichskettedisp
{\vergleichskette
{ Y
}
{ =} { \left\{ (x,y,z) \in \R^3  \mid   \alpha x^2+ \beta y^2 + \gamma z^2 = 1        \right\}
}
{ } {
}
{ } {
}
{ } {
}
}
{}{}{.}
Buktikan bahwa
\definitionsverweis {pemetaan Gauss}{}{}
untuk $Y$ bersifat bijektif.

}
{} {}

\inputaufgabe
{3}
{

Misalkan

\mavergleichskette
{\vergleichskette
{ W
}
{ \subseteq }{ \R^n
}
{  }{
}
{    }{
}
{   }{
}
}
{}{}{}
\definitionsverweis {terbuka}{}{,}
\maabb {h} { W } { \R
} {}
sebuah
\definitionsverweis {fungsi yang terdiferensialkan secara kontinu}{}{}
dan

\mavergleichskette
{\vergleichskette
{ Y
}
{ = }{ h^{-1}(c)
}
{  }{
}
{    }{
}
{   }{
}
}
{}{}{}
adalah
\definitionsverweis {serat}{}{}
di atas

\mavergleichskette
{\vergleichskette
{ c
}
{ \in }{ \R
}
{  }{
}
{    }{
}
{   }{
}
}
{}{}{,}
dengan $h$
\definitionsverweis {reguler}{}{}
di setiap titik pada $Y$. Misalkan
\maabbdisp {L} {\R^n } { \R^n
} {}
sebuah
\definitionsverweis {isomorfisme linear}{}{,}

\mavergleichskette
{\vergleichskette
{ W'
}
{ = }{ L^{-1}(W)
}
{  }{
}
{    }{
}
{   }{
}
}
{}{}{}
dan

\mavergleichskettedisp
{\vergleichskette
{ Z
}
{ =} { L^{-1}(Y)
}
{ =} { (h \circ L)^{-1}(c)
}
{ } {
}
{ } {
}
}
{}{}{.}
Buktikan bahwa konsep-konsep kurva geodesik, medan normal, medan normal satuan, dan pemetaan Gauss pada $Y$ dan $Z$ secara umum tidak saling bersesuaian
\zusatzklammer {berbeda dengan
Soal 2.14,
yang melibatkan suatu isometri} {} {.}

}
{} {}


\section*{Solusi yang disediakan oleh sumber}
\addcontentsline{toc}{section}{Solusi yang disediakan oleh sumber}
\subsection*{Solusi Soal 2.1}
\input{generated/worksheet02_exercise01_solution.id.build.tex}
\subsection*{Solusi Soal 2.2}
% Authority: German Wikiversity pageid 120435, revid 1089268, 2026-05-31T10:02:12Z.
% Exact UTF-8 wikitext witness SHA-256: 0bdd1cde6d1d0dc74757aa824634d10b6b4bf413efb4431538b7d9bac2b6f06c.
% Expanded German TeX source SHA-256: a92e556f50d2216192fc61703eea4bae3d233ab632c8eedb19bd35dec2ed89b0.

\textit{Catatan edisi: solusi sumber memakai basis ortogonal seolah-olah
ortonormal dan menuliskan turunan koordinat yang keliru. Argumen berikut
mengganti langkah tersebut.}

Karena $h'(t)\neq 0$, norma kecepatannya tak nol. Turunkan kuadrat normanya:

\[
\frac{d}{dt}\Vert h'(t)\Vert^2
=2\Vert h'(t)\Vert\bigl(\Vert h'(t)\Vert\bigr)'
=2\left\langle h'(t),h''(t)\right\rangle .
\]

Membagi dengan $2\Vert h'(t)\Vert$ memberikan

\[
\bigl(\Vert h'(t)\Vert\bigr)'
=\frac{\left\langle h'(t),h''(t)\right\rangle}{\Vert h'(t)\Vert}
=\frac{\left\langle h'(t),h''(t)\right\rangle}
       {\left\langle h'(t),h'(t)\right\rangle}\Vert h'(t)\Vert .
\]

\subsection*{Solusi Soal 2.7}
\begingroup\scriptsize
\input{generated/worksheet02_exercise07_solution.id.build.tex}
\endgroup
\subsection*{Solusi Soal 2.12}
\begingroup\small
\input{generated/worksheet02_exercise12_solution.id.build.tex}
\endgroup
\subsection*{Solusi Soal 2.13}
\input{generated/worksheet02_exercise13_solution.id.build.tex}

\part{Unit 3}
\chapter[Kuliah 3]{Kuliah 3: Kelengkungan Kurva Berparameter Panjang Busur}
\input{generated/lecture03.id.build.tex}

\chapter{Lembar Kerja 3}
\setcounter{section}{3}

  \zwischenueberschrift{Soal latihan}

\inputaufgabe
{}
{

Misalkan
\maabb {\gamma} { I } {\R^2
} {}
sebuah
\definitionsverweis {kurva berparametrisasi panjang busur}{}{}
yang dua kali terdiferensialkan secara kontinu, dan misalkan
\maabbeledisp {\delta} { -I } {\R^2
} { s } { \delta(s) = \gamma( -s)
} {,}
yaitu kurva yang ditelusuri dalam arah sebaliknya. Buktikan bahwa
\definitionsverweis {kelengkungan}{}{}
$\delta$ di $s$ merupakan negatif dari kelengkungan
\mathl{\gamma}{} di $-s$.

}
{} {}

\inputaufgabe
{}
{

Misalkan
\maabb {\gamma} { I } {\R^2
} {}
sebuah
\definitionsverweis {kurva berparametrisasi panjang busur}{}{}
yang dua kali terdiferensialkan secara kontinu, dan
\maabbdisp {L} {\R^2} {\R^2
} {}
sebuah
\definitionsverweis {rotasi}{}{.}
Buktikan bahwa
\definitionsverweis {kelengkungan}{}{}
$\gamma$ sama dengan kelengkungan
\mathl{L\circ \gamma}{}.

}
{} {}

\inputaufgabe
{}
{

Buktikan
Lema 3.6
dalam kasus dengan orientasi nonstandar.

}
{} {}

\inputaufgabe
{}
{

Misalkan
\maabb {\gamma} { I } {\R^2
} {}
sebuah
\definitionsverweis {kurva berparametrisasi panjang busur}{}{}
yang dua kali terdiferensialkan secara kontinu, dan
\maabbdisp {L} {\R^2} {\R^2
} {}
sebuah
\definitionsverweis {homoteti skalar}{}{}
dengan faktor homoteti

\mavergleichskette
{\vergleichskette
{ s
}
{ \neq }{ 0
}
{  }{
}
{    }{
}
{   }{
}
}
{}{}{.}
Bagaimana hubungan antara
\definitionsverweis {kelengkungan}{}{}
$\gamma$ dan kelengkungan
\mathl{L\circ \gamma}{?}

}
{} {}

\inputaufgabe
{}
{

Misalkan
\maabbeledisp {\gamma} {\R} { \R^2
} { t } { \begin{pmatrix}  \cos \left(  \omega t \right)  \\ \sin  \left(  \omega t \right)    \end{pmatrix}
} {,}
dengan

\mavergleichskette
{\vergleichskette
{ \omega
}
{ > }{ 0
}
{  }{
}
{    }{
}
{   }{
}
}
{}{}{.}
Buktikan bahwa terdapat tak berhingga banyak gerak melingkar dengan norma kecepatan konstan yang memiliki nilai dan turunan pertama yang sama dengan $\gamma$ pada

\mavergleichskette
{\vergleichskette
{ t
}
{ = }{ 0
}
{  }{
}
{    }{
}
{   }{
}
}
{}{}{.}

}
{} {}

Soal berikut penting dalam pembuatan komidi putar. Kesenangan menaiki komidi putar berasal dari percepatan, bukan dari kecepatan.

\inputaufgabe
{}
{

Misalkan
\maabbeledisp {\gamma} {\R} { \R^2
} { t } { \begin{pmatrix}  \cos \left(  \omega t \right)  \\ \sin  \left(  \omega t \right)    \end{pmatrix}
} {,}
dengan

\mavergleichskette
{\vergleichskette
{ \omega
}
{ > }{ 0
}
{  }{
}
{    }{
}
{   }{
}
}
{}{}{.}
Buktikan bahwa terdapat tak berhingga banyak gerak melingkar dengan norma kecepatan konstan yang berimpit dengan $\gamma$ pada

\mavergleichskette
{\vergleichskette
{ t
}
{ = }{ 0
}
{  }{
}
{    }{
}
{   }{
}
}
{}{}{}
dan juga memiliki percepatan yang sama di sana.

}
{} {}

\inputaufgabegibtloesung
{}
{

Misalkan
\maabbeledisp {\gamma} {\R} { \R^2
} { t } { \begin{pmatrix}  \cos \left(  \omega t \right)  \\ \sin  \left(  \omega t \right)    \end{pmatrix}
} {,}
dengan

\mavergleichskette
{\vergleichskette
{ \omega
}
{ > }{ 0
}
{  }{
}
{    }{
}
{   }{
}
}
{}{}{.}
Buktikan bahwa tidak ada gerak melingkar lain dengan norma kecepatan konstan yang memiliki nilai serta turunan pertama dan kedua yang sama dengan $\gamma$ pada

\mavergleichskette
{\vergleichskette
{ t
}
{ = }{ 0
}
{  }{
}
{    }{
}
{   }{
}
}
{}{}{.}

}
{} {}

\inputaufgabe
{}
{

Misalkan

\mavergleichskette
{\vergleichskette
{ W
}
{ \subseteq }{ \R^2
}
{  }{
}
{    }{
}
{   }{
}
}
{}{}{}
\definitionsverweis {terbuka}{}{,}
\maabb {h} { W } { \R
} {}
sebuah fungsi yang dua kali
\definitionsverweis {terdiferensialkan secara kontinu}{}{,}
dan

\mavergleichskette
{\vergleichskette
{ Y
}
{ = }{ h^{-1}(c)
}
{  }{
}
{    }{
}
{   }{
}
}
{}{}{}
adalah
\definitionsverweis {serat}{}{}
di atas

\mavergleichskette
{\vergleichskette
{ c
}
{ \in }{ \R
}
{  }{
}
{    }{
}
{   }{
}
}
{}{}{,}
dengan $h$
\definitionsverweis {reguler}{}{}
di setiap titik pada $Y$. Buktikan bahwa
\maabbdisp {\gamma} { I } { Y
} {}
merupakan
\definitionsverweis {kurva geodesik}{}{}
jika dan hanya jika $\gamma$
\definitionsverweis {berparametrisasi panjang busur}{}{.}

}
{} {}

\inputaufgabe
{}
{

\bild{ \begin{center}
\includegraphics[width=5.5cm]{ \bildeinlesungsvg {Euler spiral} {svg} }
\end{center}
\bildtext {} }

\bildlizenz { Euler spiral.svg } {} {AdiJapan} {Commons} {CC-by-sa 3.0} {}

Kita tinjau kurva terdiferensialkan
\zusatzklammer {yang disebut \stichwort {klotoid} {}} {} {}
\maabbeledisp {} {\R} {\R^2
} { t } { \gamma(t)
} {,}
dengan

\mavergleichskettedisp
{\vergleichskette
{ \gamma(t)
}
{ =} { \sqrt{\pi} \int_0^t  \begin{pmatrix}  \cos  { \frac{   \pi s^2 }{  2 } }  \\ \sin  { \frac{   \pi s^2 }{  2 } }     \end{pmatrix}  ds
}
{ } {
}
{ } {
}
{ } {
}
}
{}{}{.}
Buktikan bahwa
\definitionsverweis {kelengkungan}{}{}
kurva ini memenuhi persamaan

\mavergleichskettedisp
{\vergleichskette
{ \kappa(t)
}
{ =} { \sqrt{\pi} t
}
{ } {
}
{ } {
}
{ } {
}
}
{}{}{.}

}
{} {}

\inputaufgabe
{}
{

Misalkan
\maabb {f} { I } {\R
} {}
sebuah fungsi yang dua kali
\definitionsverweis {terdiferensialkan secara kontinu}{}{.}
Tentukan
\definitionsverweis {kelengkungan}{}{}
dan lingkaran kelengkungan dari grafik yang terkait.

}
{} {}

\inputaufgabe
{}
{

Misalkan
\maabb {f} {\R} {\R
} {}
dua kali
\definitionsverweis {terdiferensialkan secara kontinu}{}{}
dan

\mavergleichskettedisp
{\vergleichskette
{ \alpha(t)
}
{ \defeq} { \begin{pmatrix}  \cos f(t)  \\ \sin f(t)   \end{pmatrix}
}
{ } {
}
{ } {
}
{ } {
}
}
{}{}{.}
Tentukan
\definitionsverweis {kelengkungan}{}{}
$\alpha$ dan
\definitionsverweis {lingkaran kelengkungan}{}{}
untuk

\mavergleichskette
{\vergleichskette
{ t
}
{ \in }{ \R
}
{  }{
}
{    }{
}
{   }{
}
}
{}{}{}
dengan syarat $f'(t) \neq 0$, menggunakan rumus-rumus dalam
Lema 3.10.

}
{} {}

\inputaufgabe
{}
{

Tentukan
\definitionsverweis {kelengkungan}{}{}
\definitionsverweis {grafik}{}{}
\maabbdisp {\sin} {\R} {\R
} {}
untuk

\mavergleichskette
{\vergleichskette
{ t
}
{ = }{ { \frac{   \pi }{  2 } }
}
{  }{
}
{    }{
}
{   }{
}
}
{}{}{.}

}
{} {}

\inputaufgabe
{}
{

Tentukan
\definitionsverweis {kelengkungan}{}{}
spiral logaritmik
\maabbeledisp {\gamma} {\R} { \R^2
} { t } { \gamma(t) = e^t \left(  \cos  t   , \,   \sin  t                   \right)
} {,}
untuk setiap

\mavergleichskette
{\vergleichskette
{ t
}
{ \in }{ \R
}
{  }{
}
{    }{
}
{   }{
}
}
{}{}{.}

}
{} {}

Sebuah fungsi
\maabb {f} { I } {\R
} {}
pada interval terbuka

\mavergleichskette
{\vergleichskette
{ I
}
{ \subseteq }{ \R
}
{  }{
}
{    }{
}
{   }{
}
}
{}{}{}
disebut
\definitionswort {analitik}{,}
jika pada setiap titik

\mavergleichskette
{\vergleichskette
{ a
}
{ \in }{ I
}
{  }{
}
{    }{
}
{   }{
}
}
{}{}{}
fungsi tersebut dapat dinyatakan oleh suatu
\definitionsverweis {deret pangkat}{}{.}

Sebuah kurva disebut analitik jika setiap fungsi komponennya analitik.

\inputaufgabe
{}
{

Misalkan
\maabb {\kappa} { I } {\R
} {}
sebuah
\definitionsverweis {fungsi analitik}{}{.}
Buktikan bahwa terdapat kurva analitik
\maabb {\gamma} { I } {\R^2
} {}
yang
\definitionsverweis {kelengkungan}{}{}
di titik $\gamma(t)$ sama dengan $\kappa(t)$.

}
{} {}

\inputaufgabe
{}
{

\bild{ \begin{center}
\includegraphics[width=5.5cm]{ \bildeinlesungsvg {Evolute-parab} {svg} }
\end{center}
\bildtext {} }

\bildlizenz { Evolute-parab.svg } {} {Ag2gaeh} {Commons} {CC-by-sa 4.0} {}

Tentukan
\definitionsverweis {evolut}{}{}
kurva
\maabbeledisp {} {\R} {\R^2
} { t } { \begin{pmatrix}  t  \\ t^2    \end{pmatrix}
} {.}

}
{} {}

\inputaufgabegibtloesung
{}
{

Tentukan
\definitionsverweis {kelengkungan}{}{}
di setiap titik lingkaran satuan yang diberikan secara implisit oleh

\mavergleichskettedisp
{\vergleichskette
{ h(x,y)
}
{ =} { x^2+y^2
}
{ =} { 1
}
{ } {
}
{ } {
}
}
{}{}{}
menggunakan
Lema 3.11
dan medan gradien $h$.

}
{} {}

  \zwischenueberschrift{Soal untuk dikumpulkan}

\inputaufgabe
{2}
{

Tentukan
\definitionsverweis {kelengkungan}{}{}
\definitionsverweis {grafik}{}{}
fungsi eksponensial
\maabbdisp {\exp} {\R} {\R
} {}
untuk setiap

\mavergleichskette
{\vergleichskette
{ t
}
{ \in }{ \R
}
{  }{
}
{    }{
}
{   }{
}
}
{}{}{.}

}
{} {}

\inputaufgabe
{2}
{

Tentukan
\definitionsverweis {kelengkungan}{}{}
spiral Archimedes
\maabbeledisp {\gamma} {\R} { \R^2
} { t } { \gamma(t) = t \left(  \cos  t   , \,   \sin  t                   \right)
} {,}
untuk setiap

\mavergleichskette
{\vergleichskette
{ t
}
{ \in }{ \R
}
{  }{
}
{    }{
}
{   }{
}
}
{}{}{.}

}
{} {}

\inputaufgabe
{4}
{

Berikan sebuah contoh kurva berparametrisasi panjang busur
\maabbdisp {\gamma} {\R} {\R^2
} {}
yang dua kali terdiferensialkan secara kontinu, sedemikian sehingga untuk dua waktu

\mavergleichskette
{\vergleichskette
{ t_0
}
{ \neq }{ t_1
}
{  }{
}
{    }{
}
{   }{
}
}
{}{}{}
berlaku kesamaan

\mavergleichskette
{\vergleichskette
{ \gamma (t_0)
}
{ = }{ \gamma(t_1)
}
{  }{
}
{    }{
}
{   }{
}
}
{}{}{}
dan
\definitionsverweis {lingkaran-lingkaran kelengkungan}{}{}
pada
\mathkor {} {t_0} {dan} {t_1} {}
tidak sama.

}
{} {}

\inputaufgabe
{4}
{

Tentukan
\definitionsverweis {evolut}{}{}
kurva
\maabbeledisp {} {\R} {\R^2
} { t } { \begin{pmatrix}  2 \cos  t  \\ 3 \sin  t    \end{pmatrix}
} {.}
Di titik-titik manakah turunan evolut mempunyai titik nol?

}
{} {}

\inputaufgabe
{4}
{

Tentukan
\definitionsverweis {kelengkungan}{}{}
lingkaran satuan yang diberikan secara implisit oleh

\mavergleichskettedisp
{\vergleichskette
{ h(x,y)
}
{ =} { x^4+y^4
}
{ =} { 1
}
{ } {
}
{ } {
}
}
{}{}{}
menggunakan
Lema 3.11
dan medan gradien $h$ di titik-titik berikut.
\aufzaehlungzwei {
\mavergleichskette
{\vergleichskette
{ P
}
{ = }{ \begin{pmatrix}  1  \\ 0   \end{pmatrix}
}
{  }{
}
{    }{
}
{   }{
}
}
{}{}{,}
} {
\mavergleichskette
{\vergleichskette
{ Q
}
{ = }{ \begin{pmatrix}  { \frac{   1  }{  \sqrt[4]{2}  } }  \\ { \frac{   1  }{  \sqrt[4]{2}  } }     \end{pmatrix}
}
{  }{
}
{    }{
}
{   }{
}
}
{}{}{.}
}

}
{} {}


\section*{Solusi yang disediakan oleh sumber}
\addcontentsline{toc}{section}{Solusi yang disediakan oleh sumber}
\subsection*{Solusi Soal 3.7}
% Authority: German Wikiversity pageid 151168, revid 1095913, 2026-06-15T07:25:24Z.
% Exact UTF-8 wikitext witness SHA-256: c771b7c1ed173ddc05c8c4219c3c34e7e8099fe42079740b292d991b2ce3fb25.
% Expanded German TeX source SHA-256: 7e3437274bf4f79b6a3fa719876b09fdbcee4e10523a215e9e6f4ecb566cdb85.

Nilai-nilainya adalah

\mavergleichskettedisp
{\vergleichskette
{ \gamma(0)
}
{ =} { \begin{pmatrix}  1  \\ 0   \end{pmatrix}
}
{ } {
}
{ } {
}
{ } {
}
}
{}{}{,}

\mavergleichskettedisp
{\vergleichskette
{ \gamma'(0)
}
{ =} { \begin{pmatrix}  0  \\ \omega   \end{pmatrix}
}
{ } {
}
{ } {
}
{ } {
}
}
{}{}{}
dan

\mavergleichskettedisp
{\vergleichskette
{ \gamma^{\prime \prime} (0)
}
{ =} { \begin{pmatrix}  - \omega^2  \\ 0   \end{pmatrix}
}
{ } {
}
{ } {
}
{ } {
}
}
{}{}{.}
Suatu gerak melingkar yang memiliki nilai serta turunan pertama dan kedua yang sama dengan gerak yang diberikan pada

\mavergleichskette
{\vergleichskette
{ t
}
{ = }{ 0
}
{  }{
}
{    }{
}
{   }{
}
}
{}{}{}
harus, khususnya, mempunyai garis singgung yang sama di titik ini. Karena itu, pusat lingkaran harus terletak pada sumbu $x$. Selain itu, pusat tersebut harus berada di sebelah kiri
\mathl{\begin{pmatrix}  1  \\ 0   \end{pmatrix}}{},
karena percepatan mengarah ke kiri. Oleh sebab itu, kita dapat menuliskan gerak melingkar beraturan yang dicari, dengan pusat
\mathl{\begin{pmatrix}  x  \\ 0   \end{pmatrix}}{}
\zusatzklammer {
\mavergleichskettek
{\vergleichskettek
{ x
}
{ < }{  1
}
{  }{
}
{    }{
}
{   }{
}
}
{}{}{}} {} {}
dan jari-jari
\mathl{1-x}{}, sebagai

\mavergleichskettedisp
{\vergleichskette
{  \delta (t)
}
{ =} {  \begin{pmatrix}  x  \\ 0   \end{pmatrix} + (1-x) \begin{pmatrix}  \cos \left(  u t \right)  \\ \sin  \left(  u t \right)     \end{pmatrix}
}
{ } {
}
{ } {
}
{ } {
}
}
{}{}{.}
Kita memperoleh

\mavergleichskettedisp
{\vergleichskette
{ \delta(0)
}
{ =} { \gamma(0)
}
{ } {
}
{ } {
}
{ } {
}
}
{}{}{.}
Selanjutnya,

\mavergleichskettedisp
{\vergleichskette
{ \delta'( 0)
}
{ =} {  (1-x)  \begin{pmatrix}  0  \\ u     \end{pmatrix}
}
{ } {
}
{ } {
}
{ } {
}
}
{}{}{}
dan

\mavergleichskettedisp
{\vergleichskette
{ \delta^{\prime \prime}( 0)
}
{ =} {  (1-x)  \begin{pmatrix}  -u^2  \\ 0     \end{pmatrix}
}
{ } {
}
{ } {
}
{ } {
}
}
{}{}{.}
Jadi, agar syarat-syarat yang diminta terpenuhi, harus berlaku

\mavergleichskettedisp
{\vergleichskette
{ \omega
}
{ =} { (1-x) u
}
{ } {
}
{ } {
}
{ } {
}
}
{}{}{}
dan

\mavergleichskettedisp
{\vergleichskette
{ \omega^2
}
{ =} { (1-x) u^2
}
{ } {
}
{ } {
}
{ } {
}
}
{}{}{.}
Dari sini diperoleh

\mavergleichskettedisp
{\vergleichskette
{ \omega^2
}
{ =} {  \omega u
}
{ } {
}
{ } {
}
{ } {
}
}
{}{}{}
dan dengan demikian

\mavergleichskettedisp
{\vergleichskette
{ u
}
{ =} { \omega
}
{ } {
}
{ } {
}
{ } {
}
}
{}{}{}
serta

\mavergleichskettedisp
{\vergleichskette
{ x
}
{ =} { 0
}
{ } {
}
{ } {
}
{ } {
}
}
{}{}{.}

\subsection*{Solusi Soal 3.16}
\input{generated/worksheet03_exercise16_solution.id.build.tex}

\part{Unit 4}
\chapter[Kuliah 4]{Kuliah 4: Pemetaan Weingarten}
\setcounter{section}{4}

  \zwischenueberschrift{Pemetaan Weingarten}

Misalkan

\mavergleichskette
{\vergleichskette
{ Y
}
{ \subseteq }{ \R^n
}
{  }{
}
{    }{
}
{   }{
}
}
{}{}{}
adalah suatu
\definitionsverweis {hipermuka terdiferensialkan}{}{.}
Untuk suatu titik

\mavergleichskette
{\vergleichskette
{ P
}
{ \in }{ Y
}
{  }{
}
{    }{
}
{   }{
}
}
{}{}{}
kita mendefinisikan suatu pemetaan linear
\maabbdisp {L_P} { T_PY } { T_PY
} {}
pada ruang tangen $Y$ di P. Misalkan $N$ adalah suatu medan normal satuan terdiferensialkan pada $Y$ yang didefinisikan pada suatu lingkungan terbuka di sekitar $Y$. Gagasan pokoknya adalah merepresentasikan suatu vektor tangen

\mavergleichskette
{\vergleichskette
{ v
}
{ \in }{ T_PY
}
{  }{
}
{    }{
}
{   }{
}
}
{}{}{}
dengan suatu kurva terdiferensialkan
\maabbdisp {\gamma} { I } { Y
} {}
sehingga

\mavergleichskettedisp
{\vergleichskette
{ \gamma'(0)
}
{ =} { v
}
{ } {
}
{ } {
}
{ } {
}
}
{}{}{.}
Medan normal satuan itu kemudian mendefinisikan pemetaan
\maabbdisp {N \circ \gamma} { I } { \R^n
} {}
sepanjang $I$. Perubahan infinitesimal medan normal satuan sepanjang $\gamma$ diukur oleh limit

\mathdisp {\operatorname{lim}_{   t  \rightarrow  0   } \,  { \frac{   N(\gamma(t))-N(P)  }{  t  } }} {  }

jika limit ini ada. Menurut
Lemma 43.4 (Analysis (Osnabrück 2021-2023)),
limit ini sama dengan turunan berarah $\left(  D_{ v }  N   \right) \left(   P   \right)$ dan, pada gilirannya, sama dengan
\mathl{\left(  D N   \right)_{ P } \left(   v   \right)}{,} yaitu
\definitionsverweis {diferensial total}{}{}
yang dievaluasi pada vektor $v$. Akan ditunjukkan bahwa pemetaan
\maabbeledisp {} { T_PY } { \R^n
} {v } { \left(  D_{v}  N   \right) \left(   P   \right)
} {,}
ini mengambil nilai di $T_PY$.

\inputdefinition
{}
{

Misalkan

\mavergleichskette
{\vergleichskette
{ W
}
{ \subseteq }{ \R^n
}
{  }{
}
{    }{
}
{   }{
}
}
{}{}{}
\definitionsverweis {terbuka}{}{,}
\maabb {h} { W } { \R
} {}
suatu
\definitionsverweis {fungsi yang dua kali terdiferensialkan secara kontinu}{}{}
dan

\mavergleichskette
{\vergleichskette
{ Y
}
{ = }{ h^{-1}(c)
}
{  }{
}
{    }{
}
{   }{
}
}
{}{}{}
adalah
\definitionsverweis {serat}{}{}
di atas

\mavergleichskette
{\vergleichskette
{ c
}
{ \in }{ \R
}
{  }{
}
{    }{
}
{   }{
}
}
{}{}{,}
dengan $h$
\definitionsverweis {reguler}{}{}
pada setiap titik di $Y$. Misalkan $N$ adalah suatu
\definitionsverweis {medan normal satuan}{}{}
dan misalkan

\mavergleichskette
{\vergleichskette
{ P
}
{ \in }{ Y
}
{  }{
}
{    }{
}
{   }{
}
}
{}{}{.}
Pemetaan
\maabbeledisp {L_P} {T_P Y} { T_PY
} { v } { - \left(  D_{v}  N   \right) \left(   P  \right)
} {,}
disebut
\definitionswort {pemetaan Weingarten}{}
pada $T_PY$.

}

Perhatikan bahwa pemetaan Weingarten bergantung pada medan normal satuan, meskipun hal ini tidak selalu dinyatakan secara eksplisit. Jika $Y$ diberikan sebagai serat dari $h$, biasanya digunakan medan gradien ternormalisasi yang bersesuaian dengan $h$.


\inputfaktbeweis
{Hyperfläche/Weingartenabbildung/Tangentialraum/Fakt}
{Lemma}
{}
{

\faktsituation {Misalkan

\mavergleichskette
{\vergleichskette
{ W
}
{ \subseteq }{ \R^n
}
{  }{
}
{    }{
}
{   }{
}
}
{}{}{}
\definitionsverweis {terbuka}{}{,}
\maabb {h} { W } { \R
} {}
suatu
\definitionsverweis {fungsi yang dua kali terdiferensialkan secara kontinu}{}{}
dan

\mavergleichskette
{\vergleichskette
{ Y
}
{ = }{ h^{-1}(c)
}
{  }{
}
{    }{
}
{   }{
}
}
{}{}{}
adalah
\definitionsverweis {serat}{}{}
di atas

\mavergleichskette
{\vergleichskette
{ c
}
{ \in }{ \R
}
{  }{
}
{    }{
}
{   }{
}
}
{}{}{,}
dengan $h$
\definitionsverweis {reguler}{}{}
pada setiap titik di $Y$. Misalkan

\mavergleichskette
{\vergleichskette
{ P
}
{ \in }{ Y
}
{  }{
}
{    }{
}
{   }{
}
}
{}{}{.}}
\faktfolgerung {Maka
\definitionsverweis {pemetaan Weingarten}{}{}
$L_P$ adalah suatu endomorfisme linear dari
\definitionsverweis {ruang tangen}{}{}
$T_PY$.}
Catatan edisi: sumber memakai huruf kecil p pada ruang tangen ini, padahal titik yang dikuantifikasi adalah P.
\faktzusatz {}
\faktzusatz {}

}
{

Medan

\mavergleichskette
{\vergleichskette
{ N(P)
}
{ = }{ { \frac{
\operatorname{Grad} \,  h  (  P )  }{  \Vert {
\operatorname{Grad} \,  h  (  P ) } \Vert  } }
}
{  }{
}
{    }{
}
{   }{
}
}
{}{}{}
adalah suatu medan vektor yang terdiferensialkan secara kontinu pada suatu lingkungan terbuka

\mavergleichskette
{\vergleichskette
{ Y
}
{ \subseteq }{ U
}
{  }{
}
{    }{
}
{   }{
}
}
{}{}{}
tempat medan tersebut didefinisikan. Oleh karena itu, menurut
Proposition 46.1 (Analysis (Osnabrück 2021-2023)),

\mavergleichskettedisp
{\vergleichskette
{ \left(  D_{v}  N   \right) \left(   P  \right)
}
{ =} { \left(  D N   \right)_{ P } \left(  v  \right)
}
{ } {
}
{ } {
}
{ } {
}
}
{}{}{}
bersifat linear terhadap arah $v$. Dari syarat normal satuan diperoleh

\mavergleichskette
{\vergleichskette
{ \left\langle  N(P)  ,  N(P)   \right\rangle
}
{ = }{ 1
}
{  }{
}
{    }{
}
{   }{
}
}
{}{}{}
untuk setiap

\mavergleichskette
{\vergleichskette
{ P
}
{ \in }{ U
}
{  }{
}
{    }{
}
{   }{
}
}
{}{}{}
dan, dengan menggunakan
Soal 43.14 (Analysis (Osnabrück 2021-2023)), diperoleh

\mavergleichskettedisp
{\vergleichskette
{ 0
}
{ =} { \left(  D_{v}  \left\langle  N  ,  N  \right\rangle   \right) \left(   P  \right)
}
{ =} { 2 \left\langle  N(P)  ,  \left(  D_{v}  N   \right) \left(   P  \right)   \right\rangle
}
{ } {
}
{ } {
}
}
{}{}{.}
Jadi,
\mathl{\left(  D_{v}  N   \right) \left(   P  \right)}{} tegak lurus terhadap $N(P)$ dan, untuk

\mavergleichskette
{\vergleichskette
{ P
}
{ \in }{ Y
}
{  }{
}
{    }{
}
{   }{
}
}
{}{}{}
vektor tersebut termasuk dalam ruang tangen $T_PY$.

}

Dengan demikian, pemetaan Weingarten $L_P$ adalah negatif dari pembatasan diferensial total medan normal satuan pada ruang tangen $T_PY$, dengan kodomain ruang tangen $T_PY$. Jika diferensial total dihitung melalui
\definitionsverweis {matriks Jacobi}{}{}
dari $-N$, kita harus menentukan cara matriks itu bekerja pada suatu basis ruang tangen untuk memperoleh representasi matriks dari pemetaan Weingarten.



\inputfaktbeweis
{Ebene Kurve/Weingartenabbildung/Krümmung/Fakt}
{Lemma}
{}
{

\faktsituation {Misalkan

\mavergleichskette
{\vergleichskette
{ W
}
{ \subseteq }{ \R^2
}
{  }{
}
{    }{
}
{   }{
}
}
{}{}{}
\definitionsverweis {terbuka}{}{,}
\maabb {h} { W } { \R
} {}
suatu
\definitionsverweis {fungsi yang dua kali terdiferensialkan secara kontinu}{}{}
dan

\mavergleichskette
{\vergleichskette
{ C
}
{ = }{ h^{-1}(c)
}
{  }{
}
{    }{
}
{   }{
}
}
{}{}{}
adalah serat di atas

\mavergleichskette
{\vergleichskette
{ c
}
{ \in }{ \R
}
{  }{
}
{    }{
}
{   }{
}
}
{}{}{,}
dengan $h$
\definitionsverweis {reguler}{}{}
pada setiap titik di $C$, dan misalkan

\mavergleichskette
{\vergleichskette
{ P
}
{ \in }{ C
}
{  }{
}
{    }{
}
{   }{
}
}
{}{}{.}}
\faktfolgerung {Maka
\definitionsverweis {pemetaan Weingarten}{}{}
di $P$ adalah perkalian dengan
\definitionsverweis {kelengkungan}{}{}
$\kappa(P)$ dari $C$ di $P$.}
\faktzusatz {}
\faktzusatz {}

}
{

Pernyataan ini merupakan perumusan ulang dari
Lemma 3.11.

}

Pernyataan di atas tidak secara eksplisit merujuk pada orientasi; orientasi yang dimaksud harus turut diperhitungkan.

\inputbeispiel{}
{

Misalkan

\mavergleichskettedisp
{\vergleichskette
{ h(x,y,z)
}
{ =} { x^2+y^2+z^2
}
{ } {
}
{ } {
}
{ } {
}
}
{}{}{}
dan perhatikan serat $Y$ dari $h$ di atas $r^2$, yakni permukaan bola berjari-jari

\mavergleichskette
{\vergleichskette
{ r
}
{ > }{ 0
}
{  }{
}
{    }{
}
{   }{
}
}
{}{}{}
yang berpusat di titik asal. Misalkan

\mavergleichskette
{\vergleichskette
{ P
}
{ = }{ \left(  x_0   , \,   y_0  , \,  z_0                 \right)
}
{ \in }{ Y
}
{    }{
}
{   }{
}
}
{}{}{.}
Dengan suatu isometri, titik ini dapat dipetakan ke

\mavergleichskette
{\vergleichskette
{ Q
}
{ = }{ \left(  r   , \,   0  , \,   0                 \right)
}
{  }{
}
{    }{
}
{   }{
}
}
{}{}{}
tanpa mengubah
\definitionsverweis {pemetaan Weingarten}{}{}.
Suatu basis ruang tangennya ialah
\mathkor {} {\begin{pmatrix}  0  \\ 1 \\  0   \end{pmatrix}} {dan} {\begin{pmatrix}  0  \\ 0\\  1   \end{pmatrix}} {.}
Medan normal satuan $N$ yang mengarah ke dalam adalah
\mathl{- { \frac{   1  }{  2r } } \begin{pmatrix}  2x \\ 2y\\  2z   \end{pmatrix}}{} dan karena itu, untuk

\mavergleichskettedisp
{\vergleichskette
{ v
}
{ =} { \begin{pmatrix}  0  \\ b \\ c   \end{pmatrix}
}
{ } {
}
{ } {
}
{ } {
}
}
{}{}{}

\mavergleichskettedisp
{\vergleichskette
{ \left(  D_{v}  N   \right) \left(   Q   \right)
}
{ =} { b { \frac{   \partial   N   }{  \partial   y  } } (Q) +  c { \frac{   \partial   N   }{  \partial  z  } }  (Q)
}
{ =} { -b { \frac{   1  }{  2r } } \begin{pmatrix}  0  \\ 2 \\  0   \end{pmatrix} - c { \frac{   1  }{  2r } } \begin{pmatrix}  0  \\ 0\\  2   \end{pmatrix}
}
{ =} { - { \frac{   1  }{ r } } \begin{pmatrix}  0  \\ b \\ c   \end{pmatrix}
}
{ } {
}
}
{}{}{.}
Jadi, pemetaan Weingarten adalah perkalian dengan kebalikan ${ \frac{   1  }{ r } }$ dari jari-jari.

}

\inputbeispiel{}
{

Kita melanjutkan
Contoh 2.5.
Matriks Jacobi dari medan normal satuan

\mavergleichskettedisphandlinks
{\vergleichskettedisphandlinks
{ N( \begin{pmatrix}  x  \\ y \\ z   \end{pmatrix} )
}
{ =} { \left(  x^2+y^2+z^2  \right)^{ - { \frac{   1  }{  2 } } } \begin{pmatrix}  x  \\ y \\  -z   \end{pmatrix}
}
{ } {
}
{ } {
}
{ } {
}
}
{}{}{}
adalah

\mathdisp {\left(  x^2+y^2+z^2  \right)^{- { \frac{   3  }{  2 } } } \begin{pmatrix}  y^2+z^2  &   -xy &  -xz \\  -xy &  x^2+z^2 &  -yz \\ xz &  yz &  -x^2-y^2  \end{pmatrix}} { . }

\definitionsverweis {Pemetaan Weingarten}{}{}
adalah negatif dari pembatasan pemetaan ini pada ruang tangen permukaan tersebut; menurut
Lemma 4.2,
hasilnya kembali berada di ruang tangen. Kita mengasumsikan

\mavergleichskette
{\vergleichskette
{ x
}
{ \neq }{ 0
}
{  }{
}
{    }{
}
{   }{
}
}
{}{}{}
dan menggunakan basis
\mathkor {} {\begin{pmatrix}  y  \\ -x\\  0   \end{pmatrix}} {dan} {\begin{pmatrix}  z  \\ 0 \\  x   \end{pmatrix}} {}
dari ruang tangen. Kita mempunyai

\mavergleichskettealignhandlinks
{\vergleichskettealignhandlinks
{ \begin{pmatrix}  y^2+z^2  &   -xy &  -xz \\  -xy &  x^2+z^2 &  -yz \\ xz &  yz &  -x^2-y^2  \end{pmatrix} \begin{pmatrix}  y  \\ -x\\  0   \end{pmatrix}
}
{ =} { \begin{pmatrix}  y \left(  y^2+z^2  \right) +x^2y \\ -xy^2 -x \left(  x^2+z^2  \right) \\  0   \end{pmatrix}
}
{ =} { \left(  x^2+y^2+z^2  \right) \begin{pmatrix}  y  \\ -x\\  0   \end{pmatrix}
}
{ } {
}
{ } {
}
}
{}
{}{,}
sehingga langsung diperoleh vektor eigen $\begin{pmatrix}  y  \\ -x\\  0   \end{pmatrix}$ dengan nilai eigen
\mathl{- \left(  x^2+y^2+z^2  \right)^{ - { \frac{   1  }{  2 } } }}{}.
Selain itu,


\mavergleichskettealigndrucklinks
{\vergleichskettealigndrucklinks
{ \begin{pmatrix}  y^2+z^2  &   -xy &  -xz \\  -xy &  x^2+z^2 &  -yz \\ xz &  yz &  -x^2-y^2  \end{pmatrix} \begin{pmatrix}  z  \\ 0 \\  x   \end{pmatrix}
}
{ =} { \begin{pmatrix} z \left(  y^2+z^2  \right) -x^2z \\ -2xyz\\  xz^2-x \left(  x^2+y^2  \right)   \end{pmatrix}
}
{ =} { 2yz \begin{pmatrix}  y  \\ -x\\  0   \end{pmatrix} + \left(  z^2-x^2-y^2  \right) \begin{pmatrix}  z  \\ 0 \\  x   \end{pmatrix}
}
{ =} { 2yz \begin{pmatrix}  y  \\ -x\\  0   \end{pmatrix}- \begin{pmatrix}  z  \\ 0 \\  x   \end{pmatrix}
}
{ } {
}
}
{}{}{.}

Dengan demikian, terhadap basis ini, pemetaan Weingarten direpresentasikan oleh matriks

\mathdisp {\begin{pmatrix}  - \left(  x^2+y^2+z^2  \right)^{ - { \frac{   1  }{  2 } } }   &   - 2yz \left(  x^2+y^2+z^2  \right)^{ - { \frac{   3  }{  2 } } }    \\  0  &  \left(  x^2+y^2+z^2  \right)^{ - { \frac{   3  }{  2 } } }   \end{pmatrix}} {  }

Untuk memperoleh vektor eigen kedua dalam ruang tangen,
\zusatzklammer {dengan memperhatikan
Korolari 4.8} {} {}
cara termudah ialah menentukan suatu vektor yang tegak lurus terhadap vektor eigen pertama dan vektor normal. Hasilnya adalah vektor
\mathl{\begin{pmatrix}  xz \\ yz\\  x^2+y^2   \end{pmatrix}}{,} dan memang,
\zusatzklammer {tanpa faktor skalar di depan} {} {}


\mavergleichskettealigndrucklinks
{\vergleichskettealigndrucklinks
{ \begin{pmatrix}  y^2+z^2  &   -xy &  -xz \\  -xy &  x^2+z^2 &  -yz \\ xz &  yz &  -x^2-y^2  \end{pmatrix} \begin{pmatrix}  xz \\ yz\\  x^2+y^2   \end{pmatrix}
}
{ =} { \begin{pmatrix}  xz \left(  y^2+z^2  \right) -xy^2z -xz \left(  x^2+y^2  \right)  \\ - x^2yz +yz \left(  x^2+z^2  \right) -yz \left(  x^2+y^2  \right) \\  x^2z^2 +y^2z^2 - \left(  x^2+y^2  \right)^2   \end{pmatrix}
}
{ =} { \left(  - x^2-y^2+z^2  \right) \begin{pmatrix}  xz \\ yz\\  x^2+y^2   \end{pmatrix}
}
{ =} { - \begin{pmatrix}  xz \\ yz\\  x^2+y^2   \end{pmatrix}
}
{ } {
}
}
{}
{}{.}

Jadi, nilai eigennya adalah
\mathl{\left(  x^2+y^2+z^2  \right)^{- { \frac{   3  }{  2 } } }}{}
\zusatzklammer {hal ini juga dapat langsung dibaca dari matriks segitiga yang merepresentasikannya} {} {.}

}




\inputfaktbeweis
{Hyperfläche/Weingartenabbildung/Zweite Ableitung/Skalarprodukt/Fakt}
{Lemma}
{}
{

\faktsituation {Misalkan

\mavergleichskette
{\vergleichskette
{ W
}
{ \subseteq }{ \R^n
}
{  }{
}
{    }{
}
{   }{
}
}
{}{}{}
\definitionsverweis {terbuka}{}{,}
\maabb {h} { W } { \R
} {}
suatu
\definitionsverweis {fungsi yang dua kali terdiferensialkan secara kontinu}{}{}
dan

\mavergleichskette
{\vergleichskette
{ Y
}
{ = }{ h^{-1}(c)
}
{  }{
}
{    }{
}
{   }{
}
}
{}{}{}
adalah
\definitionsverweis {serat}{}{}
di atas

\mavergleichskette
{\vergleichskette
{ c
}
{ \in }{ \R
}
{  }{
}
{    }{
}
{   }{
}
}
{}{}{,}
dengan $h$
\definitionsverweis {reguler}{}{}
pada setiap titik di $Y$. Misalkan $N$ adalah suatu
\definitionsverweis {medan normal satuan}{}{}
dan misalkan

\mavergleichskette
{\vergleichskette
{ P
}
{ \in }{ Y
}
{  }{
}
{    }{
}
{   }{
}
}
{}{}{.}
Misalkan $\gamma$ adalah suatu realisasi yang terdiferensialkan dua kali pada $Y$ dari vektor tangen

\mavergleichskette
{\vergleichskette
{ v
}
{ \in }{ T_PY
}
{  }{
}
{    }{
}
{   }{
}
}
{}{}{.}}
\faktfolgerung {Maka

\mavergleichskettedisp
{\vergleichskette
{ \left\langle  \gamma^{\prime \prime} (0)  ,  N(P)  \right\rangle
}
{ =} { \left\langle  L_P(v) , v  \right\rangle
}
{ } {
}
{ } {
}
{ } {
}
}
{}{}{.}}
\faktzusatz {}
\faktzusatz {}

}
{

Misalkan
\maabbeledisp {\gamma} { I } { Y
} { t } { \gamma(t)
} {,}
terdiferensialkan dua kali, dengan

\mavergleichskette
{\vergleichskette
{ \gamma(0)
}
{ = }{ P
}
{  }{
}
{    }{
}
{   }{
}
}
{}{}{}
dan

\mavergleichskette
{\vergleichskette
{ \gamma'(0)
}
{ = }{ v
}
{ \in }{ T_PY
}
{    }{
}
{   }{
}
}
{}{}{.}
Kita mempunyai

\mavergleichskettedisp
{\vergleichskette
{ \left\langle  \gamma'(t)  ,  N(\gamma(t))   \right\rangle
}
{ =} { 0
}
{ } {
}
{ } {
}
{ } {
}
}
{}{}{}
untuk semua

\mavergleichskette
{\vergleichskette
{ t
}
{ \in }{ I
}
{  }{
}
{    }{
}
{   }{
}
}
{}{}{}
sehingga, dengan
Soal 43.14 (Analysis (Osnabrück 2021-2023)), diperoleh

\mavergleichskettealign
{\vergleichskettealign
{ 0
}
{ =} { \left\langle  \gamma'(t)  ,  N(\gamma(t))   \right\rangle' (0)
}
{ =} { \left\langle  \gamma^{\prime \prime}(0)  ,  N(\gamma(0))  \right\rangle + \left\langle  \gamma'(0)  ,  \left(  D_{ \gamma'(0) }  N   \right) \left(   \gamma(0)   \right)   \right\rangle
}
{ =} { \left\langle  \gamma^{\prime \prime}(0) , N(P)  \right\rangle - \left\langle  L_P(v)  ,  v   \right\rangle
}
{ } {
}
}
{}
{}{.}

}

Dalam pernyataan di atas tidak diperlukan penetapan tambahan mengenai orientasi, sebab medan normal satuan muncul pada kedua ruas, dan pada ruas kanan muncul melalui pemetaan Weingarten. Perhatikan pula bahwa $\gamma$ tidak harus memiliki kecepatan konstan. Dengan demikian, terdapat banyak realisasi dan percepatan $\gamma^{\prime \prime} (0)$ tidak ditentukan secara tunggal. Namun, nilai
\mathl{\left\langle  \gamma^{\prime \prime}(0)  ,  N(P)  \right\rangle}{,} yang menggambarkan komponen normal percepatan, tidak berubah karena $\gamma$ bergerak pada hipermuka tersebut.
Catatan edisi: sumber kehilangan argumen N(P) dalam hasil kali dalam pada kalimat sebelumnya; identitas yang baru saja dibuktikan menentukannya secara unik.



\inputfaktbeweis
{Hyperfläche/Weingartenabbildung/Selbstadjungiert/Fakt}
{Satz}
{}
{

\faktsituation {Misalkan

\mavergleichskette
{\vergleichskette
{ W
}
{ \subseteq }{ \R^n
}
{  }{
}
{    }{
}
{   }{
}
}
{}{}{}
\definitionsverweis {terbuka}{}{,}
\maabb {h} { W } { \R
} {}
suatu
\definitionsverweis {fungsi yang dua kali terdiferensialkan secara kontinu}{}{}
dan

\mavergleichskette
{\vergleichskette
{ Y
}
{ = }{ h^{-1}(c)
}
{  }{
}
{    }{
}
{   }{
}
}
{}{}{}
adalah serat di atas

\mavergleichskette
{\vergleichskette
{ c
}
{ \in }{ \R
}
{  }{
}
{    }{
}
{   }{
}
}
{}{}{,}
dengan $h$
\definitionsverweis {reguler}{}{}
pada setiap titik di $Y$.}
\faktfolgerung {Maka
\definitionsverweis {pemetaan Weingarten}{}{}
$L_P$ pada setiap titik

\mavergleichskette
{\vergleichskette
{ P
}
{ \in }{ Y
}
{  }{
}
{    }{
}
{   }{
}
}
{}{}{}
bersifat
\definitionsverweis {adjoin-diri (self-adjoint)}{}{.}}
\faktzusatz {}
\faktzusatz {}

}
{

Untuk vektor

\mavergleichskette
{\vergleichskette
{ v,w
}
{ \in }{ T_PY
}
{  }{
}
{    }{
}
{   }{
}
}
{}{}{}
kita harus menunjukkan bahwa

\mavergleichskettedisp
{\vergleichskette
{ \left\langle  v  ,  L_P(w)  \right\rangle
}
{ =} { \left\langle L_P(v) ,  w  \right\rangle
}
{ } {
}
{ } {
}
{ } {
}
}
{}{}{}
tersebut. Dengan

\mavergleichskette
{\vergleichskette
{ N(P)
}
{ = }{ { \frac{
\operatorname{Grad} \,  h  (  P )  }{  \Vert {
\operatorname{Grad} \,  h  (  P ) } \Vert  } }
}
{  }{
}
{    }{
}
{   }{
}
}
{}{}{}
menurut
Lemma 45.10 (Analysis (Osnabrück 2021-2023)), kita mempunyai

\mavergleichskettealignhandlinks
{\vergleichskettealignhandlinks
{ \left\langle  v  ,  L_P(w)  \right\rangle
}
{ =} { -\left\langle  v  ,  \left(  D_{w}  N   \right) \left(   P   \right)   \right\rangle
}
{ =} { -\left\langle  v  ,  \left(  D_{w}  { \frac{
\operatorname{Grad} \,  h  }{  \Vert {
\operatorname{Grad} \,  h } \Vert  } }   \right) \left(   P   \right)     \right\rangle
}
{ =} { -\left\langle  v  ,  \left(  D_{w}  { \frac{   1  }{  \Vert {
\operatorname{Grad} \,  h } \Vert  } }   \right) \left(   P   \right) \cdot
\operatorname{Grad} \,  h  (  P ) + { \frac{   1  }{  \Vert {
\operatorname{Grad} \,  h  (  P ) } \Vert  } } \cdot  \left(  D_{w}
\operatorname{Grad} \,  h   \right) \left(   P   \right)   \right\rangle
}
{ =} { - { \frac{   1  }{  \Vert {
\operatorname{Grad} \,  h  (  P ) } \Vert  } } \left\langle  v  ,  \left(  D_{w}
\operatorname{Grad} \,  h   \right) \left(   P   \right)    \right\rangle
}
}
{}
{}{,}
karena suku pertama tegak lurus terhadap vektor tangen $v$. Dalam fungsi-fungsi koordinat,

\mavergleichskettealignhandlinks
{\vergleichskettealignhandlinks
{ \left(  D_{w}
\operatorname{Grad} \,  h   \right) \left(   P   \right)
}
{ =} { \left(  D_{w}  \left(  { \frac{   \partial   h   }{  \partial   x_1   } }   , \,   \ldots  , \,   { \frac{   \partial   h   }{  \partial   x_n   } }                 \right)    \right) \left(   P   \right)
}
{ =} { \sum_{j = 1}^n  w_j  { \frac{   \partial    }{  \partial   x_j   } }  \left(  { \frac{   \partial   h   }{  \partial   x_1   } }   , \,   \ldots  , \,   { \frac{   \partial   h   }{  \partial   x_n   } }                 \right) (P)
}
{ =} { \left(  \sum_{j = 1}^n  w_j  { \frac{   \partial    }{  \partial   x_j   } }  { \frac{   \partial   h   }{  \partial   x_1   } } (P)   , \,   \ldots  , \,   \sum_{j = 1}^n  w_j  { \frac{   \partial    }{  \partial   x_j   } }  { \frac{   \partial   h   }{  \partial   x_n   } } (P)                 \right)
}
{ } {
}
}
{}
{}{.}
Dengan demikian, ungkapan di atas sama dengan

\mavergleichskettealign
{\vergleichskettealign
{ \left\langle  v  ,  L_P(w)  \right\rangle
}
{ =} { - { \frac{   1  }{  \Vert {
\operatorname{Grad} \,  h  (  P ) } \Vert  } } \left\langle  v  ,  \left(  D_{w}
\operatorname{Grad} \,  h   \right) \left(   P   \right)   \right\rangle
}
{ =} { -  { \frac{   1  }{  \Vert {
\operatorname{Grad} \,  h  (  P ) } \Vert  } }  \sum_{i,j = 1}^n v_i  w_j  { \frac{   \partial    }{  \partial   x_j   } } { \frac{   \partial   h   }{  \partial   x_i   } }  (P)
}
{ } {
}
{ } {
}
}
{}
{}{.}
Menurut
Teorema 44.10 (Analysis (Osnabrück 2021-2023)),
urutan turunan parsial dapat dipertukarkan, sehingga peran
\mathkor {} {v} {dan} {w} {}
juga dapat dipertukarkan.

}



\inputfaktbeweis
{Hyperfläche/Weingartenabbildung/Selbstadjungiert/Eigenwerte/Fakt}
{Korollar}
{}
{

\faktsituation {Misalkan

\mavergleichskette
{\vergleichskette
{ W
}
{ \subseteq }{ \R^n
}
{  }{
}
{    }{
}
{   }{
}
}
{}{}{}
\definitionsverweis {terbuka}{}{,}
\maabb {h} { W } { \R
} {}
suatu
\definitionsverweis {fungsi yang dua kali terdiferensialkan secara kontinu}{}{}
dan

\mavergleichskette
{\vergleichskette
{ Y
}
{ = }{ h^{-1}(c)
}
{  }{
}
{    }{
}
{   }{
}
}
{}{}{}
adalah serat di atas

\mavergleichskette
{\vergleichskette
{ c
}
{ \in }{ \R
}
{  }{
}
{    }{
}
{   }{
}
}
{}{}{,}
dengan $h$
\definitionsverweis {reguler}{}{}
pada setiap titik di $Y$.}
\faktfolgerung {Maka
\definitionsverweis {pemetaan Weingarten}{}{}
$L_P$ pada setiap titik

\mavergleichskette
{\vergleichskette
{ P
}
{ \in }{ Y
}
{  }{
}
{    }{
}
{   }{
}
}
{}{}{}
\definitionsverweis {dapat didiagonalkan}{}{}
dengan
\definitionsverweis {nilai-nilai eigen}{}{}
real dan ruang-ruang eigen yang saling ortogonal.}
\faktzusatz {}
\faktzusatz {}

}
{

Hal ini mengikuti dari
Teorema 4.7
dan
Teorema 41.11 (Lineare Algebra (Osnabrück 2024-2025)).

}



\inputfaktbeweis
{Differenzierbare Funktion/Graph/Hyperfläche/Weingartenabbildung/Fakt}
{Lemma}
{}
{

\faktsituation {Misalkan

\mavergleichskette
{\vergleichskette
{ V
}
{ \subseteq }{ \R^{n-1}
}
{  }{
}
{    }{
}
{   }{
}
}
{}{}{}
suatu
\definitionsverweis {himpunan terbuka}{}{}
dan misalkan

\mavergleichskette
{\vergleichskette
{ Y
}
{ \subseteq }{ V \times \R
}
{  }{
}
{    }{
}
{   }{
}
}
{}{}{}
adalah
\definitionsverweis {grafik}{}{}
dari suatu fungsi yang
\definitionsverweis {dua kali terdiferensialkan secara kontinu}{}{}
\maabbdisp {f} { V } {\R
} {.}}
\faktfolgerung {Maka
\definitionsverweis {pemetaan Weingarten}{}{}
$L_P$ pada suatu titik

\mavergleichskette
{\vergleichskette
{ P
}
{ = }{ (x_1 , \ldots , x_{n-1}, f(x_1 , \ldots , x_{n-1}))
}
{ \in }{ Y
}
{    }{
}
{   }{
}
}
{}{}{}
\zusatzklammer {terhadap medan normal satuan yang mengarah ke atas} {} {}
dapat direpresentasikan sebagai berikut. Tuliskan \mathl{g=\operatorname{Grad} f}{} sebagai vektor kolom dan
\mathl{G=I_{n-1}+gg^{\mathsf T}}{,}
yakni matriks metrik dalam basis grafik. Matriks pemetaan Weingarten dalam basis tersebut adalah
\mathl{G^{-1}{ \frac{   1  }{  \sqrt{ 1+ \left(  { \frac{   \partial   f   }{  \partial   x_1   } }  \right)^2 + \cdots + \left(  { \frac{   \partial   f   }{  \partial   x_{n-1}   } }  \right)^2 }  } } \operatorname{Hess}_{   } \,  f}{}, dengan $\operatorname{Hess}_{   } \,  f$ menyatakan
\definitionsverweis {matriks Hesse}{}{}
dari $f$, apabila vektor-vektor dasar

\mavergleichskette
{\vergleichskette
{ v
}
{ \in }{ \R^{n-1}
}
{  }{
}
{    }{
}
{   }{
}
}
{}{}{}
diidentifikasi dengan vektor-vektor tangen dari $T_PY$ dalam pengertian
Contoh 1.2.}
\faktzusatz {}
\faktzusatz {}

}
{

Kita melanjutkan
Contoh 1.2.
Medan normal satuan yang mengarah ke atas diberikan oleh

\mavergleichskettedisphandlinks
{\vergleichskettedisphandlinks
{ N(x_1  , \ldots , x_{n-1}, f(x_1 , \ldots , x_{n-1}))
}
{ =} { { \frac{   1  }{  \sqrt{ 1+ \left(  { \frac{   \partial   f   }{  \partial   x_1   } }  \right)^2 + \cdots + \left(  { \frac{   \partial   f   }{  \partial   x_{n-1}   } }  \right)^2 }  } } \begin{pmatrix}  -{ \frac{   \partial   f   }{  \partial   x_{1}   } }  \\\vdots\\  -{ \frac{   \partial   f   }{  \partial   x_{n-1}   } }\\ 1   \end{pmatrix}
}
{ } {
}
{ } {
}
{ } {
}
}
{}{}{}
yang akan digunakan. Kita meninjau lintasan

\mavergleichskettedisp
{\vergleichskette
{ \gamma(t)
}
{ =} { \begin{pmatrix}  x_1+tv_1  \\ \vdots \\  x_{n-1} +tv_{n-1}\\f \begin{pmatrix}  x_1+tv_1  \\ \vdots \\  x_{n-1} +tv_{n-1}   \end{pmatrix}   \end{pmatrix}
}
{ } {
}
{ } {
}
{ } {
}
}
{}{}{}
pada grafik yang bersesuaian dengan vektor dasar $v$. Turunan keduanya adalah

\mavergleichskettedisp
{\vergleichskette
{ \gamma^{\prime \prime} (0)
}
{ =} { \begin{pmatrix}  0  \\ \vdots \\  0\\ \sum_{1 \leq i, j \leq n-1}  { \frac{   \partial    }{  \partial   x_i   } }  { \frac{   \partial   f   }{  \partial   x_j   } } v_iv_j   \end{pmatrix}
}
{ } {
}
{ } {
}
{ } {
}
}
{}{}{.}
Menurut
Lemma 4.6,

\mavergleichskettealigndrucklinks
{\vergleichskettealigndrucklinks
{ \left\langle  L_P( \gamma'(0) ) ,  \gamma'(0)  \right\rangle
}
{ =} { \left\langle  \gamma^{\prime \prime} (0) ,  N(x_1 , \ldots , x_{n-1}, f(x_1  , \ldots , x_{n-1}))   \right\rangle
}
{ =} { { \frac{   1  }{  \sqrt{ 1 + \left(  { \frac{   \partial   f   }{  \partial   x_1   } }  \right)^2 + \cdots + \left(  { \frac{   \partial   f  }{  \partial   x_{n-1}   } }  \right)^2 }  } } \left\langle  \begin{pmatrix}  0  \\\vdots\\  0\\ \sum_{ 1 \leq i, j \leq n-1 }  { \frac{   \partial    }{  \partial   x_i   } }  { \frac{   \partial   f   }{  \partial   x_j   } } v_iv_j   \end{pmatrix}  ,  \begin{pmatrix}  -{ \frac{   \partial   f   }{  \partial   x_{1}   } }  \\\vdots\\  -{ \frac{   \partial   f   }{  \partial   x_{n-1}   } }\\ 1   \end{pmatrix}   \right\rangle
}
{ =} { { \frac{   \sum_{ 1 \leq i, j \leq n-1 } { \frac{   \partial    }{  \partial   x_i   } }  { \frac{   \partial   f   }{  \partial   x_j   } } v_i v_j     }{  \sqrt{ 1+  \left(  { \frac{   \partial   f   }{  \partial   x_1   } }  \right)^2 + \cdots + \left(  { \frac{   \partial   f  }{  \partial   x_{n-1}   } }  \right)^2 }  } }
}
{ =} { { \frac{   1  }{  \sqrt{ 1+ \left(  { \frac{   \partial   f   }{  \partial   x_1   } }  \right)^2 + \cdots + \left(  { \frac{   \partial   f   }{  \partial   x_{n-1}   } }  \right)^2 }  } }  \left( v_1   , \,   \ldots  , \,   v_{n-1}                 \right)  \operatorname{Hess}_{   } \,  f  \begin{pmatrix} v_1  \\ \vdots \\  v_{n-1}   \end{pmatrix}
}
}
{}{}{.}
Artinya, menurut
Teorema 4.7,
\definitionsverweis {bentuk bilinear}{}{}
simetris

\mathl{\left\langle  L_P(-) ,  -  \right\rangle}{} pada ruang tangen $T_PY$ bersesuaian dengan bentuk bilinear yang diberikan oleh matriks Hesse yang diskalakan
\zusatzklammer {oleh faktor di depan} {} {}
pada $\R^{n-1}$, apabila vektor yang sama disubstitusikan pada kedua tempat.
Menurut
Lemma 38.10 (Lineare Algebra (Osnabrück 2024-2025)),
kedua bentuk bilinear tersebut berimpit untuk sembarang pasangan argumen. Akan tetapi, ini menentukan bentuk fundamental kedua, bukan langsung matriks operator dalam basis grafik. Jika \mathl{A}{} adalah matriks $L_P$ dalam basis grafik dan \mathl{G=I_{n-1}+gg^{\mathsf T}}{} adalah matriks metriknya, maka matriks bentuk bilinear di atas adalah \mathl{GA}{.} Oleh karena itu,

\mathdisp {GA={ \frac{   1  }{  \sqrt{1+\lVert g\rVert^2}  } }\operatorname{Hess}_{   }\, f} { , }

sehingga

\mathdisp {A=G^{-1}{ \frac{   1  }{  \sqrt{1+\lVert g\rVert^2}  } }\operatorname{Hess}_{   }\, f} { . }

Catatan edisi: sumber menghilangkan faktor invers metrik G dan dengan demikian menyamakan matriks bentuk fundamental kedua dengan matriks pemetaan Weingarten. Rumus dan simpulan di atas memperbaiki pembedaan tersebut.

}


\chapter{Lembar Kerja 4}
\setcounter{section}{4}

  \zwischenueberschrift{Soal latihan}

\inputaufgabe
{}
{

Misalkan

\mavergleichskette
{\vergleichskette
{ h(x_1 , \ldots , x_n)
}
{ = }{ \sum_{i = 1}^n a_ix_i
}
{  }{
}
{    }{
}
{   }{
}
}
{}{}{}
adalah suatu
\definitionsverweis {bentuk linear}{}{}
$\neq 0$ dan

\mavergleichskette
{\vergleichskette
{ Y
}
{ = }{ h^{-1}(c)
}
{  }{
}
{    }{
}
{   }{
}
}
{}{}{}
adalah hipermuka yang bersesuaian. Buktikan bahwa
\definitionsverweis {pemetaan Weingarten}{}{}
merupakan pemetaan nol.

}
{} {}

\inputaufgabe
{}
{

Misalkan

\mavergleichskettedisp
{\vergleichskette
{ Y
}
{ =} { \left\{  \left(  x   , \,   y  , \,  z                 \right)   \mid   x^2+y^2 = z^2 , \,  \left(  x   , \,   y  , \,  z                 \right)  \neq  \left(  0   , \,   0  , \,   0                 \right)       \right\}
}
{ \subseteq} { W
}
{ =} { \R^3 \setminus \{  \left(  0   , \,   0  , \,   0                 \right) \}
}
{ } {
}
}
{}{}{}
adalah kerucut standar, dan misalkan

\mavergleichskette
{\vergleichskette
{ P
}
{ = }{ \left(  x   , \,   y  , \,  z                 \right)
}
{ \in }{ Y
}
{    }{
}
{   }{
}
}
{}{}{.}
Misalkan

\mavergleichskettedisp
{\vergleichskette
{ \gamma(t)
}
{ =} { \left(  x\cos t-y\sin t   , \,   x\sin t+y\cos t  , \,  z                 \right)
}
{ } {
}
{ } {
}
{ } {
}
}
{}{}{.}

Catatan edisi: kurva pada sumber tidak selalu melalui P dan tidak selalu tetap berada pada kerucut. Rumus di atas menggantinya dengan rotasi standar pada bidang xy.
Tentukan

\mathdisp {\operatorname{lim}_{   t  \rightarrow  0   } \,  { \frac{  N( \gamma(t) ) -  N( \gamma(0) )  }{ t } }} { , }

dengan $N$ medan normal satuan yang diperoleh dari gradien $x^2+y^2-z^2$.

}
{} {}

\inputaufgabe
{}
{

Misalkan

\mavergleichskettedisp
{\vergleichskette
{ Y
}
{ =} { \left\{  \left(  x   , \,   y  , \,  z                 \right)   \mid   x^2+y^2 =1         \right\}
}
{ \subseteq} { \R^3
}
{ } {
}
{ } {
}
}
{}{}{}
adalah silinder standar, dan misalkan

\mavergleichskette
{\vergleichskette
{ P
}
{ = }{ \left(  1   , \,   0  , \,  z                 \right)
}
{ \in }{ Y
}
{    }{
}
{   }{
}
}
{}{}{.}
Misalkan

\mavergleichskettedisp
{\vergleichskette
{ \gamma(t)
}
{ =} { \left(  \cos  t   , \,   \sin  t  , \,  z                 \right)
}
{ } {
}
{ } {
}
{ } {
}
}
{}{}{.}
Tentukan

\mathdisp {\operatorname{lim}_{   t  \rightarrow  0   } \,  { \frac{  N( \gamma(t) ) -  N( \gamma(0) )  }{ t } }} { , }

dengan $N$ medan normal satuan yang diperoleh dari gradien $x^2+y^2-1$.

}
{} {}

\inputaufgabe
{}
{

Misalkan

\mavergleichskettedisp
{\vergleichskette
{ Y
}
{ =} { \left\{  \left(  x   , \,   y  , \,  z                 \right)   \mid   x^2+y^2 =1         \right\}
}
{ \subseteq} { \R^3
}
{ } {
}
{ } {
}
}
{}{}{}
dan misalkan

\mavergleichskette
{\vergleichskette
{ P
}
{ \in }{ Y
}
{  }{
}
{    }{
}
{   }{}
}
{}{}{}
adalah suatu titik. Buktikan bahwa
\definitionsverweis {pemetaan Weingarten}{}{}
$L_P$ tidak bijektif, tetapi juga bukan pemetaan nol.

}
{} {}

\inputaufgabe
{}
{

Buktikan bahwa suatu
\definitionsverweis {permukaan terdiferensialkan}{}{}

\mavergleichskette
{\vergleichskette
{ Y
}
{ \subseteq }{ \R^3
}
{  }{
}
{    }{
}
{   }{
}
}
{}{}{}
dapat memuat suatu garis

\mavergleichskette
{\vergleichskette
{ G
}
{ \subseteq }{ Y
}
{  }{
}
{    }{
}
{   }{
}
}
{}{}{}
yang vektor arahnya, jika dipandang sebagai vektor di
\definitionsverweis {ruang singgung}{}{}
$T_PY$ untuk suatu titik

\mavergleichskette
{\vergleichskette
{ P
}
{ \in }{ G
}
{  }{
}
{    }{
}
{   }{
}
}
{}{}{,}
tidak harus termasuk dalam
\definitionsverweis {kernel}{}{}
dari
\definitionsverweis {pemetaan Weingarten}{}{}
$L_P$.

}
{} {}

\inputaufgabe
{}
{

Misalkan

\mavergleichskette
{\vergleichskette
{ W
}
{ \subseteq }{ \R^n
}
{  }{
}
{    }{
}
{   }{
}
}
{}{}{}
\definitionsverweis {terbuka}{}{,}
\maabb {h} { W } { \R
} {}
adalah suatu
\definitionsverweis {fungsi yang terdiferensialkan secara kontinu sebanyak dua kali}{}{}
dan

\mavergleichskette
{\vergleichskette
{ Y
}
{ = }{ h^{-1}(c)
}
{  }{
}
{    }{
}
{   }{
}
}
{}{}{}
adalah
\definitionsverweis {serat}{}{}
di atas

\mavergleichskette
{\vergleichskette
{ c
}
{ \in }{ \R
}
{  }{
}
{    }{
}
{   }{
}
}
{}{}{,}
dengan $h$
\definitionsverweis {reguler}{}{}
di setiap titik pada $Y$.
Catatan edisi: sumber hanya mengasumsikan h kelas C1. Di sini hipotesis diperkuat menjadi h kelas C2 agar medan normal gradien dan pemetaan Weingarten yang ditanyakan terdefinisi sebagaimana pada kuliah.
Kita juga memandang $h$ sebagai pemetaan $\tilde{h}$ pada
\mathl{W \times \R^m}{}{,}
dengan $m$ variabel terakhir tidak muncul secara eksplisit dalam $\tilde{h}$. Misalkan

\mavergleichskettedisp
{\vergleichskette
{ Z
}
{ =} { \tilde{h}^{-1}(c)
}
{ =} { Y \times \R^m
}
{ \subseteq} { W \times \R^{m}
}
{ } {
}
}
{}{}{.}
Misalkan

\mavergleichskette
{\vergleichskette
{ Q
}
{ \in }{ Z
}
{  }{
}
{    }{
}
{   }{
}
}
{}{}{}
adalah suatu titik yang terletak di atas

\mavergleichskette
{\vergleichskette
{ P
}
{ \in }{ Y
}
{  }{
}
{    }{
}
{   }{
}
}
{}{}{.}
Bagaimana hubungan antara
\definitionsverweis {pemetaan-pemetaan Weingarten}{}{}
\mathkor {} {L_Q} {dan} {L_P} {?}

}
{} {}

\inputaufgabegibtloesung
{}
{

Misalkan
\mathkor {} {r} {dan} {R} {}
adalah bilangan real positif dengan

\mavergleichskette
{\vergleichskette
{ 0
}
{ < }{ r
}
{ < }{ R
}
{    }{
}
{   }{
}
}
{}{}{,}
dan kita tinjau torus
\zusatzklammer {tertanam} {} {}

\mavergleichskettedisp
{\vergleichskette
{ T
}
{ =} { \left\{ (x,y,z) \in \R^3  \mid   \left(  \sqrt{x^2+y^2} -R  \right)^2 +z^2  = r^2         \right\}
}
{ } {
}
{ } {
}
{ } {
}
}
{}{}{.}
\aufzaehlungdrei{Tentukan suatu
\definitionsverweis {medan normal satuan}{}{}
$N$ pada $T$.
}{Pada titik
\mathl{\left( R-r  , \,   0  , \,   0                 \right)}{} dan untuk vektor

\mavergleichskette
{\vergleichskette
{ v
}
{ = }{ \begin{pmatrix}  0  \\ 1 \\  0   \end{pmatrix}
}
{  }{
}
{    }{
}
{   }{
}
}
{}{}{}
tentukan
\definitionsverweis {turunan berarah}{}{}

\mathl{D_vN}{.}
}{Pada titik
\mathl{\left( R-r  , \,   0  , \,   0                 \right)}{} dan untuk vektor

\mavergleichskette
{\vergleichskette
{ w
}
{ = }{ \begin{pmatrix}  0  \\ 0\\  1   \end{pmatrix}
}
{  }{
}
{    }{
}
{   }{
}
}
{}{}{}
tentukan
\definitionsverweis {turunan berarah}{}{}

\mathl{D_w N}{.}
}

}
{} {}

\inputaufgabe
{}
{

Untuk grafik $Y$ dari fungsi

\mavergleichskettedisp
{\vergleichskette
{ f(x,y)
}
{ =} { xy
}
{ } {
}
{ } {
}
{ } {
}
}
{}{}{}
tentukan
\definitionsverweis {pemetaan Weingarten}{}{}
$L_P$ pada setiap titik

\mavergleichskettedisp
{\vergleichskette
{ P
}
{ =} { (x,y,f(x,y) )
}
{ } {
}
{ } {
}
{ } {
}
}
{}{}{.}
Tentukan nilai eigen, vektor eigen, dan suatu matriks diagonal untuk $L_P$.

}
{} {}

\inputaufgabe
{}
{

Misalkan

\mavergleichskette
{\vergleichskette
{ W
}
{ \subseteq }{ \R^n
}
{  }{
}
{    }{
}
{   }{
}
}
{}{}{}
\definitionsverweis {terbuka}{}{,}
\maabb {h} { W } { \R
} {}
adalah suatu
\definitionsverweis {fungsi yang terdiferensialkan secara kontinu sebanyak dua kali}{}{}
dan

\mavergleichskette
{\vergleichskette
{ Y
}
{ = }{ h^{-1}(c)
}
{  }{
}
{    }{
}
{   }{
}
}
{}{}{}
adalah
\definitionsverweis {serat}{}{}
di atas

\mavergleichskette
{\vergleichskette
{ c
}
{ \in }{ \R
}
{  }{
}
{    }{
}
{   }{
}
}
{}{}{,}
dengan $h$
\definitionsverweis {reguler}{}{}
di setiap titik pada $Y$.
Catatan edisi: sumber hanya mengasumsikan h kelas C1. Di sini hipotesis diperkuat menjadi h kelas C2 agar pemetaan Weingarten yang dibandingkan terdefinisi sebagaimana pada kuliah.
Misalkan
\maabbdisp {L} {\R^n } { \R^n
} {}
adalah suatu
\definitionsverweis {isometri linear}{}{,}

\mavergleichskette
{\vergleichskette
{ W'
}
{ = }{ L^{-1}(W)
}
{  }{
}
{    }{
}
{   }{
}
}
{}{}{}
dan

\mavergleichskettedisp
{\vergleichskette
{ Z
}
{ =} { L^{-1}(Y)
}
{ =} { (h \circ L)^{-1}(c)
}
{ } {
}
{ } {
}
}
{}{}{.}
Buktikan bahwa
\definitionsverweis {pemetaan Weingarten}{}{}
$L_P$ untuk

\mavergleichskette
{\vergleichskette
{ P
}
{ \in }{ Y
}
{  }{
}
{    }{
}
{   }{
}
}
{}{}{}
bersesuaian dengan pemetaan Weingarten untuk

\mavergleichskette
{\vergleichskette
{ Q
}
{ = }{ L^{-1}(P)
}
{ \in }{ Z
}
{    }{
}
{   }{
}
}
{}{}{,}
jika melalui $L$ ruang singgung $T_PY$ dan $T_QZ$ dihubungkan satu sama lain.

}
{} {}

\inputaufgabegibtloesung
{}
{

Untuk permukaan $Y$ yang diberikan oleh

\mavergleichskettedisp
{\vergleichskette
{ x^2+y^2+2z^2
}
{ =} { 4
}
{ } {
}
{ } {
}
{ } {
}
}
{}{}{}
dan titik

\mavergleichskettedisp
{\vergleichskette
{ P
}
{ =} { \left(  1  , \,   1  , \,   1                 \right)
}
{ \in} { Y
}
{ } {
}
{ } {
}
}
{}{}{}
tentukan suatu matriks diagonal untuk
\definitionsverweis {pemetaan Weingarten}{}{}
$L_P$.

}
{} {}

\inputaufgabe
{}
{

Misalkan

\mavergleichskette
{\vergleichskette
{ W
}
{ \subseteq }{ \R^n
}
{  }{
}
{    }{
}
{   }{
}
}
{}{}{}
\definitionsverweis {terbuka}{}{,}
\maabb {h} { W } { \R
} {}
adalah suatu
\definitionsverweis {fungsi yang terdiferensialkan secara kontinu sebanyak dua kali}{}{}
dan

\mavergleichskette
{\vergleichskette
{ Y
}
{ = }{ h^{-1}(c)
}
{  }{
}
{    }{
}
{   }{
}
}
{}{}{}
adalah
\definitionsverweis {serat}{}{}
di atas

\mavergleichskette
{\vergleichskette
{ c
}
{ \in }{ \R
}
{  }{
}
{    }{
}
{   }{
}
}
{}{}{,}
dengan $h$
\definitionsverweis {reguler}{}{}
di setiap titik pada $Y$.
Bagaimana
\definitionsverweis {pemetaan Weingarten}{}{}
dan
\definitionsverweis {ruang-ruang eigennya}{}{}
berubah jika kita beralih dari suatu
\definitionsverweis {orientasi}{}{}
pada $Y$ ke orientasi yang berlawanan?

}
{} {}

  \zwischenueberschrift{Soal untuk dikumpulkan}

\inputaufgabe
{4}
{

Dalam
Contoh 4.5,
tentukan matriks representasi
\definitionsverweis {pemetaan Weingarten}{}{}
untuk titik-titik berbentuk
\mathl{\begin{pmatrix}  0  \\ y \\ z   \end{pmatrix}}{} terhadap basis
\mathkor {} {\begin{pmatrix}  y  \\ 0 \\  0   \end{pmatrix}} {dan} {\begin{pmatrix}  0  \\ z \\  y   \end{pmatrix}} {}
dari ruang singgung. Tentukan nilai-nilai eigen dan suatu basis yang terdiri atas vektor-vektor eigen.

}
{} {}

\inputaufgabe
{5}
{

Untuk permukaan $Y$ yang diberikan oleh

\mavergleichskettedisp
{\vergleichskette
{ x^2+2y^2+5z^2
}
{ =} { 11
}
{ } {
}
{ } {
}
{ } {
}
}
{}{}{}
dan titik

\mavergleichskettedisp
{\vergleichskette
{ P
}
{ =} { \left(  2  , \,   1  , \,   1                 \right)
}
{ \in} { Y
}
{ } {
}
{ } {
}
}
{}{}{}
tentukan suatu matriks diagonal untuk
\definitionsverweis {pemetaan Weingarten}{}{}
$L_P$.

}
{} {}

\inputaufgabe
{4}
{

Untuk hipermuka yang diberikan oleh

\mavergleichskettedisp
{\vergleichskette
{ x^2+2y^2+3z^2+4w^2
}
{ =} { 10
}
{ } {
}
{ } {
}
{ } {
}
}
{}{}{}

\mavergleichskette
{\vergleichskette
{ Y
}
{ \subseteq }{ \R^4
}
{  }{
}
{    }{
}
{   }{
}
}
{}{}{}
dan titik

\mavergleichskettedisp
{\vergleichskette
{ P
}
{ =} { \left(  1  , \,   1  , \,   1 , \,   1                \right)
}
{ \in} { Y
}
{ } {
}
{ } {
}
}
{}{}{}
tentukan suatu matriks untuk
\definitionsverweis {pemetaan Weingarten}{}{}
$L_P$ terhadap suatu
\definitionsverweis {basis}{}{}
yang sesuai dari
\definitionsverweis {ruang singgung}{}{}
$T_PY$.

}
{} {}

\inputaufgabe
{6}
{

Untuk grafik $Y$ dari fungsi

\mavergleichskettedisp
{\vergleichskette
{ f(x,y)
}
{ =} { x^3y^5
}
{ } {
}
{ } {
}
{ } {
}
}
{}{}{}
tentukan
\definitionsverweis {pemetaan Weingarten}{}{}
$L_P$ pada titik-titik

\mavergleichskette
{\vergleichskette
{ P_1
}
{ = }{ (0,0,0)
}
{  }{
}
{    }{
}
{   }{
}
}
{}{}{,}

\mavergleichskette
{\vergleichskette
{ P_2
}
{ = }{ (0,1,0)
}
{  }{
}
{    }{
}
{   }{
}
}
{}{}{,}

\mavergleichskette
{\vergleichskette
{ P_3
}
{ = }{ (2,1,8)
}
{  }{
}
{    }{
}
{   }{
}
}
{}{}{.}
Untuk masing-masing titik, tentukan nilai eigen, vektor eigen, dan suatu matriks diagonal untuk $L_P$.

}
{} {}


\section*{Solusi yang disediakan oleh sumber}
\addcontentsline{toc}{section}{Solusi yang disediakan oleh sumber}
\subsection*{Solusi Soal 4.7}
\input{generated/worksheet04_exercise07_solution.id.build.tex}
\subsection*{Solusi Soal 4.10}
\input{generated/worksheet04_exercise10_solution.id.build.tex}

\part{Unit 5}
\chapter[Kuliah 5]{Kuliah 5: Kelengkungan Utama}
\input{generated/lecture05.id.build.tex}

\chapter{Lembar Kerja 5}
\setcounter{section}{5}

  \zwischenueberschrift{Soal latihan}

\inputaufgabegibtloesung
{}
{

Misalkan

\mavergleichskettedisp
{\vergleichskette
{ f(x,y)
}
{ =} { ax^2+by^2+cxy
}
{ } {
}
{ } {
}
{ } {
}
}
{}{}{.}
Tentukan
\definitionsverweis {kelengkungan utama}{}{}
dan
\definitionsverweis {arah kelengkungan utama}{}{}
pada
\definitionsverweis {grafik}{}{}
fungsi $f$ di titik nol.

}
{} {}

\inputaufgabe
{}
{

Tentukan
\definitionsverweis {kelengkungan utama}{}{}
dan
\definitionsverweis {arah kelengkungan utama}{}{}
pada setiap titik
\definitionsverweis {grafik}{}{}
fungsi
\maabbeledisp {f} {\R^{n-1}} {\R
} { \left(  x_1   , \,   \ldots  , \,   x_{n-1}                 \right) } { a_1x_1^2  + \cdots + a_{n-1} x_{n-1}^2
} {.}

}
{} {}

\inputaufgabe
{}
{

Misalkan

\mavergleichskettedisp
{\vergleichskette
{ f(x,y,z)
}
{ =} { ax^2 +by^2 +c z^2 +dxy
}
{ } {
}
{ } {
}
{ } {
}
}
{}{}{.}
Hitung
\definitionsverweis {polinom karakteristik}{}{}
dari
\definitionsverweis {pemetaan Weingarten}{}{}
pada
\definitionsverweis {grafik}{}{}
fungsi $f$ di titik nol. Hitung pula
\definitionsverweis {kelengkungan utama}{}{}
dan
\definitionsverweis {arah kelengkungan utama}{}{.}

}
{} {}

\inputaufgabe
{}
{

Misalkan

\mavergleichskettedisp
{\vergleichskette
{ f(x,y,z)
}
{ =} { ax^2 +by^2 +c z^2 +dxy + e yz
}
{ } {
}
{ } {
}
{ } {
}
}
{}{}{.}
Hitung
\definitionsverweis {polinom karakteristik}{}{}
dari
\definitionsverweis {pemetaan Weingarten}{}{}
pada
\definitionsverweis {grafik}{}{}
fungsi $f$ di titik nol. Dalam kasus ini, seberapa sulitkah menentukan
\definitionsverweis {kelengkungan utama}{}{?}

}
{} {}

\inputaufgabe
{}
{

Misalkan

\mavergleichskettedisp
{\vergleichskette
{ f(x,y)
}
{ =} { x^2+y^3
}
{ } {
}
{ } {
}
{ } {
}
}
{}{}{.}
Tentukan
\definitionsverweis {kelengkungan utama}{}{,}
\definitionsverweis {arah kelengkungan utama}{}{}
dan
\definitionsverweis {kelengkungan Gauss}{}{}
pada
\definitionsverweis {grafik}{}{}
fungsi $f$ di setiap titik

\mavergleichskette
{\vergleichskette
{ P
}
{ = }{ (Q,f(Q))
}
{  }{
}
{    }{
}
{   }{
}
}
{}{}{,}

\mavergleichskette
{\vergleichskette
{ Q
}
{ \in }{ \R^2
}
{  }{
}
{    }{
}
{   }{
}
}
{}{}{.}

}
{} {}

\inputaufgabe
{}
{

Misalkan

\mavergleichskette
{\vergleichskette
{ V
}
{ \subseteq }{ \R^{n-1}
}
{  }{
}
{    }{
}
{   }{
}
}
{}{}{}
\definitionsverweis {terbuka}{}{}
dan misalkan
\maabb {f} { V } {\R
} {}
adalah suatu
\definitionsverweis {fungsi yang dua kali terdiferensialkan secara kontinu}{}{.}
Buktikan bahwa
\definitionsverweis {kelengkungan utama}{}{}
dan
\definitionsverweis {arah kelengkungan utama}{}{}
pada
\definitionsverweis {grafik}{}{}
fungsi $f$ di suatu titik

\mavergleichskette
{\vergleichskette
{ P
}
{ = }{ (Q,f(Q))
}
{  }{
}
{    }{
}
{   }{
}
}
{}{}{}
hanya bergantung pada turunan-turunan parsial pertama dan kedua dari $f$ di $Q$.

}
{} {}

\inputaufgabe
{}
{

Misalkan

\mavergleichskette
{\vergleichskette
{ W
}
{ \subseteq }{ \R^n
}
{  }{
}
{    }{
}
{   }{
}
}
{}{}{}
\definitionsverweis {terbuka}{}{,}
\maabb {h} { W } { \R
} {}
adalah suatu
\definitionsverweis {fungsi yang dua kali terdiferensialkan secara kontinu}{}{}
dan

\mavergleichskette
{\vergleichskette
{ Y
}
{ = }{ h^{-1}(c)
}
{  }{
}
{    }{
}
{   }{
}
}
{}{}{}
adalah
\definitionsverweis {serat}{}{}
di atas

\mavergleichskette
{\vergleichskette
{ c
}
{ \in }{ \R
}
{  }{
}
{    }{
}
{   }{
}
}
{}{}{,}
dengan $h$
\definitionsverweis {reguler}{}{}
di setiap titik pada $Y$.
Catatan edisi: sumber hanya mengasumsikan $h$ kelas $C^1$, padahal kelengkungan utama dan arah utamanya memerlukan pemetaan Weingarten. Hipotesis di atas diperkuat menjadi $C^2$ agar semua objek yang ditanyakan terdefinisi sebagaimana pada kuliah.
Kita juga memandang $h$ sebagai suatu pemetaan $\tilde{h}$ pada
\mathl{W \times \R^m}{},
dengan $m$ variabel terakhir tidak muncul secara eksplisit dalam $\tilde{h}$. Misalkan

\mavergleichskettedisp
{\vergleichskette
{ Z
}
{ =} { \tilde{h}^{-1}(c)
}
{ =} { Y \times \R^m
}
{ \subseteq} { W \times \R^{m}
}
{ } {
}
}
{}{}{.}
Misalkan

\mavergleichskette
{\vergleichskette
{ Q
}
{ \in }{ Z
}
{  }{
}
{    }{
}
{   }{
}
}
{}{}{}
adalah suatu titik yang terletak di atas

\mavergleichskette
{\vergleichskette
{ P
}
{ \in }{ Y
}
{  }{
}
{    }{
}
{   }{
}
}
{}{}{.}
Bagaimana hubungan antara
\definitionsverweis {kelengkungan utama}{}{}
dan
\definitionsverweis {arah kelengkungan utama}{}{}
dari $Z$ di $Q$ dengan yang dari $Y$ di $P$?

}
{} {}

Misalkan

\mavergleichskette
{\vergleichskette
{ W
}
{ \subseteq }{ \R^n
}
{  }{
}
{    }{
}
{   }{
}
}
{}{}{}
\definitionsverweis {terbuka}{}{,}
\maabb {h} { W } { \R
} {}
adalah suatu
\definitionsverweis {fungsi yang dua kali terdiferensialkan secara kontinu}{}{}
dan

\mavergleichskette
{\vergleichskette
{ Y
}
{ = }{ h^{-1}(c)
}
{  }{
}
{    }{
}
{   }{
}
}
{}{}{}
adalah
\definitionsverweis {serat}{}{}
di atas

\mavergleichskette
{\vergleichskette
{ c
}
{ \in }{ \R
}
{  }{
}
{    }{
}
{   }{
}
}
{}{}{,}
dengan $h$
\definitionsverweis {reguler}{}{}
di setiap titik pada $Y$. Untuk suatu titik

\mavergleichskette
{\vergleichskette
{ P
}
{ \in }{ Y
}
{  }{
}
{    }{
}
{   }{
}
}
{}{}{}
\definitionsverweis {determinan}{}{}
dari
\definitionsverweis {pemetaan Weingarten}{}{}
$L_P$ disebut
\definitionswort {kelengkungan Gauss--Kronecker}{}
dari $Y$ di $P$.

Kelengkungan Gauss--Kronecker merupakan perumuman langsung dari kelengkungan Gauss.

\inputaufgabe
{}
{

Misalkan

\mavergleichskette
{\vergleichskette
{ W
}
{ \subseteq }{ \R^n
}
{  }{
}
{    }{
}
{   }{
}
}
{}{}{}
\definitionsverweis {terbuka}{}{,}
\maabb {h} { W } { \R
} {}
adalah suatu
\definitionsverweis {fungsi yang dua kali terdiferensialkan secara kontinu}{}{}
dan

\mavergleichskette
{\vergleichskette
{ Y
}
{ = }{ h^{-1}(c)
}
{  }{
}
{    }{
}
{   }{
}
}
{}{}{}
adalah
\definitionsverweis {serat}{}{}
di atas

\mavergleichskette
{\vergleichskette
{ c
}
{ \in }{ \R
}
{  }{
}
{    }{
}
{   }{
}
}
{}{}{,}
dengan $h$
\definitionsverweis {reguler}{}{}
di setiap titik pada $Y$. Misalkan

\mavergleichskette
{\vergleichskette
{ P
}
{ \in }{ Y
}
{  }{
}
{    }{
}
{   }{
}
}
{}{}{.}
Buktikan bahwa hasil kali
\definitionsverweis {kelengkungan utama}{}{}
dari $Y$ di $P$ sama dengan
\definitionsverweis {kelengkungan Gauss--Kronecker}{}{}
dari $Y$ di $P$.

}
{} {}

\inputaufgabe
{}
{

Misalkan

\mavergleichskettedisp
{\vergleichskette
{ h(x,y,z,w)
}
{ =} { x^2+y^3+z^4+w^5
}
{ } {
}
{ } {
}
{ } {
}
}
{}{}{}
dan

\mavergleichskette
{\vergleichskette
{ Y
}
{ = }{ h^{-1}(1)
}
{  }{
}
{    }{
}
{   }{
}
}
{}{}{.}
Untuk titik

\mavergleichskettedisp
{\vergleichskette
{ P
}
{ =} { \left(  1   , \,   -1  , \,   1 , \,   0                \right)
}
{ \in} { Y
}
{ } {
}
{ } {
}
}
{}{}{}
tentukan
\definitionsverweis {pemetaan Weingarten}{}{}
$L_P$,
\definitionsverweis {polinom karakteristik}{}{}
darinya, serta
\definitionsverweis {kelengkungan Gauss--Kronecker}{}{}
dari $Y$ di $P$.

}
{} {}

\inputaufgabe
{}
{

Misalkan

\mavergleichskettedisp
{\vergleichskette
{ h(x,y,z)
}
{ =} { x^2+y^3+z^5
}
{ } {
}
{ } {
}
{ } {
}
}
{}{}{}
pada

\mavergleichskette
{\vergleichskette
{ W
}
{ = }{ \R^3 \setminus \{0\}
}
{  }{
}
{    }{
}
{   }{
}
}
{}{}{}
dan

\mavergleichskette
{\vergleichskette
{ Y
}
{ = }{ h^{-1}(0)
}
{  }{
}
{    }{
}
{   }{
}
}
{}{}{.}
Tentukan
\definitionsverweis {kelengkungan normal}{}{}
di titik

\mavergleichskette
{\vergleichskette
{ P
}
{ = }{ \begin{pmatrix}  1  \\ -1\\  0   \end{pmatrix}
}
{  }{
}
{    }{
}
{   }{
}
}
{}{}{}
dalam arah vektor singgung ternormalisasi
\mathl{{ \frac{   1  }{  \sqrt{29}  } }  \begin{pmatrix}  3  \\ -2\\  4   \end{pmatrix}}{.}

}
{} {}

  \zwischenueberschrift{Soal untuk dikumpulkan}

\inputaufgabe
{4}
{

Misalkan

\mavergleichskettedisp
{\vergleichskette
{ f(x,y)
}
{ =} { x^3 +x^2y+3xy^2-y^4
}
{ } {
}
{ } {
}
{ } {
}
}
{}{}{.}
Tentukan
\definitionsverweis {kelengkungan utama}{}{,}
\definitionsverweis {arah kelengkungan utama}{}{}
dan
\definitionsverweis {kelengkungan Gauss}{}{}
pada
\definitionsverweis {grafik}{}{}
fungsi $f$ di setiap titik

\mavergleichskette
{\vergleichskette
{ P
}
{ = }{ (Q,f(Q))
}
{  }{
}
{    }{
}
{   }{
}
}
{}{}{,}

\mavergleichskette
{\vergleichskette
{ Q
}
{ \in }{ \R^2
}
{  }{
}
{    }{
}
{   }{
}
}
{}{}{.}

}
{} {}

\inputaufgabe
{4}
{

Misalkan

\mavergleichskettedisp
{\vergleichskette
{ f(x,y,z)
}
{ =} { -7x^2 +3y^2 -5z^2+ 6xy
}
{ } {
}
{ } {
}
{ } {
}
}
{}{}{.}
Tentukan
\definitionsverweis {kelengkungan utama}{}{}
dan
\definitionsverweis {arah kelengkungan utama}{}{}
pada
\definitionsverweis {grafik}{}{}
fungsi $f$ di titik nol.

}
{} {}

\inputaufgabe
{6}
{

Misalkan

\mavergleichskettedisp
{\vergleichskette
{ h(x,y,z,w)
}
{ =} { x^2+yz+z^3+xyzw
}
{ } {
}
{ } {
}
{ } {
}
}
{}{}{}
dan

\mavergleichskette
{\vergleichskette
{ Y
}
{ = }{ h^{-1}(2)
}
{  }{
}
{    }{
}
{   }{
}
}
{}{}{.}
Untuk titik

\mavergleichskettedisp
{\vergleichskette
{ P
}
{ =} { \left(  1   , \,   1  , \,   1 , \,   -1                \right)
}
{ \in} { Y
}
{ } {
}
{ } {
}
}
{}{}{}
tentukan
\definitionsverweis {pemetaan Weingarten}{}{}
$L_P$,
\definitionsverweis {polinom karakteristik}{}{}
darinya, serta
\definitionsverweis {kelengkungan Gauss--Kronecker}{}{}
dari $Y$ di $P$.

}
{} {Petunjuk: Selesaikan persamaan terhadap suatu variabel yang sesuai dan perlakukan hipermuka tersebut sebagai grafik dalam variabel-variabel lainnya.}

\inputaufgabe
{6 (2+2+2)}
{

Misalkan

\mavergleichskettedisp
{\vergleichskette
{ h(x,y,z)
}
{ =} { x^3+y^3+z^3
}
{ } {
}
{ } {
}
{ } {
}
}
{}{}{}
pada

\mavergleichskette
{\vergleichskette
{ W
}
{ = }{ \R^3 \setminus \{0\}
}
{  }{
}
{    }{
}
{   }{
}
}
{}{}{}
dan

\mavergleichskette
{\vergleichskette
{ Y
}
{ = }{ h^{-1}(0)
}
{  }{
}
{    }{
}
{   }{
}
}
{}{}{.}
Misalkan

\mavergleichskette
{\vergleichskette
{ P
}
{ = }{ \begin{pmatrix}  1  \\ 1\\  - \sqrt[3]{2}   \end{pmatrix}
}
{ \in }{ Y
}
{    }{
}
{   }{
}
}
{}{}{}
dan

\mavergleichskettedisp
{\vergleichskette
{ v
}
{ =} { \begin{pmatrix}  1  \\ -1\\  0   \end{pmatrix}
}
{ \in} { T_PY
}
{ } {
}
{ } {
}
}
{}{}{.}
\aufzaehlungdrei{Tentukan suatu persamaan bidang untuk
\definitionsverweis {bidang normal}{}{}
yang bersesuaian dengan $v$.
}{Tentukan
\definitionsverweis {kelengkungan normal}{}{}
di $P$ dalam arah vektor singgung ternormalisasi
\mathl{{ \frac{   v  }{  \Vert { v } \Vert  } }}{.}
}{Tentukan suatu kurva berparametrisasi panjang busur
\maabbdisp {\gamma} { I } { Y
} {}
dengan

\mavergleichskette
{\vergleichskette
{ \gamma (0)
}
{ = }{ P
}
{  }{
}
{    }{
}
{   }{
}
}
{}{}{,}

\mavergleichskette
{\vergleichskette
{ \gamma' (0)
}
{ = }{ { \frac{   v  }{  \Vert { v } \Vert  } }
}
{  }{
}
{    }{
}
{   }{
}
}
{}{}{.}
Verifikasikan berlakunya
Teorema 5.14
dalam situasi ini.
}

Catatan edisi: pada butir ketiga, orientasikan bidang normal dengan menetapkan pasangan $(v/\lVert v\rVert,N(P))$ sebagai basis positif, sehingga normal satuan kurva irisan di dalam bidang yang bersesuaian dengan arah singgung $v/\lVert v\rVert$ adalah $N(P)$. Tanpa konvensi ini, verifikasi kelengkungan bertanda hanya menentukan kesamaan nilai mutlak.

}
{} {}

\inputaufgabe
{2}
{

Buatlah sketsa suatu permukaan terdiferensialkan $Y$ di ruang dan suatu bidang normal $E$ yang irisannya dengan permukaan tersebut menghasilkan suatu kurva yang tidak reguler pada setidaknya satu titik.

}
{} {}


\section*{Solusi yang disediakan oleh sumber}
\addcontentsline{toc}{section}{Solusi yang disediakan oleh sumber}
\subsection*{Solusi Soal 5.1}
\input{generated/worksheet05_exercise01_solution.id.build.tex}

\part{Unit 6}
\chapter[Kuliah 6]{Kuliah 6: Turunan Kovarian, Medan Vektor Paralel, dan Transpor Paralel}
\input{generated/lecture06.id.build.tex}

\chapter{Lembar Kerja 6}
\input{generated/worksheet06.id.build.tex}

\section*{Solusi yang disediakan oleh sumber}
\addcontentsline{toc}{section}{Solusi yang disediakan oleh sumber}
\subsection*{Solusi Soal 6.2}
\aufzaehlungdrei{Matriks Jacobi tersebut adalah

\mathdisp {\begin{pmatrix}  0   &   1  &  0  \\  -1 &  0  &  0 \\ 0  &  0 &  0  \end{pmatrix}} { . }

}{Sebagai medan normal, kita ambil $\begin{pmatrix}  x  \\ y \\ z   \end{pmatrix}$. Di

\mavergleichskette
{\vergleichskette
{ P
}
{ = }{ \begin{pmatrix}  0  \\ 1 \\  0   \end{pmatrix}
}
{  }{
}
{    }{
}
{   }{
}
}
{}{}{,}
vektor-vektor
\mathkor {} {\begin{pmatrix}  1  \\ 0 \\  0   \end{pmatrix}} {dan} {\begin{pmatrix}  0  \\ 0\\  1   \end{pmatrix}} {}
membentuk suatu basis ruang singgung. Kita mempunyai

\mavergleichskettedisp
{\vergleichskette
{ \begin{pmatrix}  0   &   1  &  0  \\  -1 &  0  &  0 \\ 0  &  0 &  0  \end{pmatrix} \begin{pmatrix}  1  \\ 0 \\  0   \end{pmatrix}
}
{ =} { \begin{pmatrix}  0  \\ -1\\  0   \end{pmatrix}
}
{ } {
}
{ } {
}
{ } {
}
}
{}{}{}
dan

\mavergleichskettedisp
{\vergleichskette
{  \begin{pmatrix}  0   &   1  &  0  \\  -1 &  0  &  0 \\ 0  &  0 &  0  \end{pmatrix} \begin{pmatrix}  0  \\ 0\\  1   \end{pmatrix}
}
{ =} { \begin{pmatrix}  0  \\ 0\\  0   \end{pmatrix}
}
{ } {
}
{ } {
}
{ } {
}
}
{}{}{.}
Di sini,
\mathl{\begin{pmatrix}  0  \\ -1\\  0   \end{pmatrix}}{} merupakan kelipatan vektor normal, sehingga proyeksi ortogonalnya ke ruang singgung juga merupakan vektor nol. Dengan demikian, pemetaan
\mathl{v \mapsto \nabla_vF}{} adalah pemetaan nol.
}{Di

\mavergleichskette
{\vergleichskette
{ Q
}
{ = }{ \begin{pmatrix}  0  \\ 0\\  1   \end{pmatrix}
}
{  }{
}
{    }{
}
{   }{
}
}
{}{}{,}
vektor-vektor
\mathkor {} {\begin{pmatrix}  1  \\ 0 \\  0   \end{pmatrix}} {dan} {\begin{pmatrix}  0  \\ 1 \\  0   \end{pmatrix}} {}
membentuk suatu basis ruang singgung. Kita mempunyai

\mavergleichskettedisp
{\vergleichskette
{ \begin{pmatrix}  0   &   1  &  0  \\  -1 &  0  &  0 \\ 0  &  0 &  0  \end{pmatrix} \begin{pmatrix}  1  \\ 0 \\  0   \end{pmatrix}
}
{ =} { \begin{pmatrix}  0  \\ -1\\  0   \end{pmatrix}
}
{ } {
}
{ } {
}
{ } {
}
}
{}{}{}
dan

\mavergleichskettedisp
{\vergleichskette
{ \begin{pmatrix}  0   &   1  &  0  \\  -1 &  0  &  0 \\ 0  &  0 &  0  \end{pmatrix} \begin{pmatrix}  0  \\ 1 \\  0   \end{pmatrix}
}
{ =} { \begin{pmatrix}  1  \\ 0 \\  0   \end{pmatrix}
}
{ } {
}
{ } {
}
{ } {
}
}
{}{}{.}
Kedua vektor citra ini sudah tangensial pada $Y$ di $Q$. Suatu matriks yang mendeskripsikan
\mathl{v \mapsto \nabla_vF}{} adalah

\mathdisp {\begin{pmatrix}  0   &   1  \\  -1  &  0 \end{pmatrix}} { . }

\emph{Catatan edisi: Tanda pada kedua kolom diperbaiki agar matriks ini sesuai dengan citra kedua vektor basis yang dihitung tepat di atasnya.}

}

\subsection*{Solusi Soal 6.6}
Kita mempunyai

\mavergleichskettealign
{\vergleichskettealign
{ (N (\gamma(t)) \times F(t))'
}
{ =} { (N (\gamma(t)))' \times F(t) +  N (\gamma(t)) \times F'(t)
}
{ =} { \left(  D N   \right)_{ \gamma(t)} \left(   \gamma'(t)   \right)  \times F(t) +  N (\gamma(t)) \times F'(t)
}
{ } {
}
{ } {
}
}
{}
{}{.}
Kita perlu membuktikan bahwa vektor ini selalu tegak lurus terhadap ruang singgung $T_{\gamma(t)}Y$. Karena $F'(t)$ tegak lurus terhadap ruang singgung, turunan tersebut merupakan kelipatan skalar vektor normal $N(\gamma(t))$; oleh karena itu, suku kedua sama dengan $0$ menurut
Fakta (3).
Pada suku di sebelah kiri, $F(t)$ bersifat tangensial menurut asumsi, sedangkan $\left(  D N   \right)_{ \gamma(t)} \left(   \gamma'(t)   \right)$ bersifat tangensial menurut
Fakta.
Oleh karena itu, menurut
Fakta (6),
suku di sebelah kiri tegak lurus terhadap ruang singgung. Jadi, medan vektor tersebut secara keseluruhan bersifat paralel.

\subsection*{Solusi Soal 6.9}
\input{generated/worksheet06_exercise09_solution.id.build.tex}

\part{Unit 7}
\chapter[Kuliah 7]{Kuliah 7: Konsep sebuah manifold}
\setcounter{section}{7}

  \zwischenueberschrift{Konsep sebuah manifold}

Dalam geometri diferensial, gagasan tentang manifold berada di pusat pembahasan. Sebagai contoh, kita meninjau Bumi
\zusatzklammer {permukaannya} {} {,}
yang dalam sejarah ilmu pengetahuan lama dianggap sebagai sebuah cakram, dan ada alasan yang baik untuk itu. Secara lokal Bumi tampak seperti sebuah bidang. Hal ini juga tercermin dalam peta yang dibuat tentangnya. Sebuah peta adalah sebuah \anfuehrung{lembaran}{,} yang titik-titiknya berada dalam bijeksi dengan suatu bagian permukaan Bumi. Khususnya untuk bagian-bagian kecil hal ini tidak bermasalah, tetapi pada peta yang harus menggambarkan bagian besar atau bahkan seluruh Bumi segera muncul pertanyaan tentang apa yang dipertahankan oleh peta dan apa yang tidak: pertanyaan tentang pelestarian panjang, pelestarian luas, pelestarian sudut, titik-titik yang hilang atau muncul berkali-kali, pertanyaan perpanjangan, pertanyaan kelengkungan, dan seterusnya.

\bild{ \begin{center}
\includegraphics[width=5.5cm]{ \bildeinlesungpng {Stereographic_projection_in_3D} {png} }
\end{center}
\bildtext {Proyeksi stereografis ketika bidang tidak diletakkan melalui khatulistiwa, melainkan melalui Kutub Selatan.} }

\bildlizenz { Stereographic projection in 3D.png } {} {Mark.Howison} {en.Wikipedia} {PD} {}

Mula-mula kita membahas \stichwort {proyeksi stereografis} {} dari permukaan bola.

\inputbeispiel{
}
{

Kita meninjau permukaan bola

\mavergleichskettedisp
{\vergleichskette
{ K
}
{ =} { \left\{ (x,y,z)  \mid   x^2+y^2+z^2 = 1        \right\}
}
{ } {
}
{ } {
}
{ } {
}
}
{}{}{}
dan menyebut titik

\mavergleichskette
{\vergleichskette
{ N
}
{ = }{ (0,0,1)
}
{  }{
}
{    }{
}
{   }{
}
}
{}{}{}
sebagai Kutub Utara dan titik

\mavergleichskette
{\vergleichskette
{ S
}
{ = }{ (0,0,-1)
}
{  }{
}
{    }{
}
{   }{
}
}
{}{}{}
sebagai Kutub Selatan. Sebuah titik

\mathbed {P=(x,y,z) \in K} {}
 {P \neq N} {}
 {} {} {} {,}
bersama Kutub Utara menentukan sebuah garis tunggal di ruang, yang diparameterkan oleh

\mathdisp {(tx,ty ,1+t(z-1))=(0,0,1)  + t ((x,y,z) -(0,0,1)), \, t \in \R} { , }

Vektor \mathl{(x,y,z) -(0,0,1)}{,} yang menentukan garis ini, tidak sejajar dengan bidang $x-y$; artinya garis tersebut mempunyai tepat satu titik potong dengan bidang ini. Titik itu diperoleh untuk parameter

\mavergleichskettedisp
{\vergleichskette
{ t
}
{ =} { { \frac{   1  }{  1-z } }
}
{ } {
}
{ } {
}
{ } {
}
}
{}{}{,}
yaitu titik

\mathdisp {\left(  { \frac{   x  }{  1-z  } }, { \frac{   y  }{  1-z  } } ,0  \right)} { . }

Kita rangkum konstruksi ini sebagai pemetaan
\maabbeledisp {\alpha} {K \setminus \{N\}} { \R^2
} {(x,y,z)} { \left(  { \frac{   x  }{  1-z  } }, { \frac{   y  }{  1-z  } }  \right)
} {}
Pemetaan ini secara intuitif jelas merupakan bijeksi, dan hal tersebut juga mudah dihitung dari rumus-rumusnya. Pemetaan invers diperoleh dengan menghubungkan sebuah titik \mathl{(u,v,0)}{} di bidang dengan Kutub Utara, lalu menghitung titik tembusnya $\neq N$ dengan permukaan bola. Hal ini menghasilkan syarat

\mavergleichskettealignhandlinks
{\vergleichskettealignhandlinks
{ \Vert {(0,0,1) +a((u,v,0) - (0,0,1) ) } \Vert^2
}
{ =} { \Vert {(au,av,1-a) } \Vert^2
}
{ =} { a^2u^2+a^2v^2+(1-a)^2
}
{ =} { 1
}
{ } {
}
}
{}
{}{,}
yang berujung pada

\mavergleichskettedisp
{\vergleichskette
{ a \left(  au^2+av^2 -2 +a  \right)
}
{ =} { 0
}
{ } {
}
{ } {
}
{ } {
}
}
{}{}{}
Solusinya adalah

\mavergleichskette
{\vergleichskette
{ a
}
{ = }{ 0
}
{  }{
}
{    }{
}
{   }{
}
}
{}{}{}
yang bersesuaian dengan Kutub Utara dan tidak kita minati, sehingga kita memperoleh

\mavergleichskette
{\vergleichskette
{ a
}
{ = }{ { \frac{   2  }{ u^2+v^2+1 } }
}
{  }{
}
{    }{
}
{   }{
}
}
{}{}{}
dan dengan demikian pemetaan
\maabbeledisp {} { \R^2 } { K \setminus \{N\}
} { (u,v) } { \left(  { \frac{   2u }{ u^2+v^2+1 } }, { \frac{   2v }{ u^2+v^2+1 } } ,1- { \frac{   2  }{ u^2+v^2+1 } }  \right)
} {.}
Jadi bidang riil adalah \definitionsverweis {homeomorf}{}{} dengan permukaan bola yang \anfuehrung{dilubangi}{} pada satu titik.

Pertimbangan yang sama dapat dilakukan untuk \mathl{K \setminus \{S\}}{}. Hal ini menghasilkan pemetaan
\maabbeledisp {\beta} { K \setminus \{S\} } { \R^2
} { (x,y,z) } { \left(  { \frac{   x  }{ z+1 } }, { \frac{   y  }{ z+1 } }  \right)
} {}
dengan pemetaan invers
\maabbeledisp {} { \R^2 } { K \setminus \{S\}
} { (s,t) } { \left(  { \frac{   2s }{ s^2+t^2+1 } }, { \frac{   2t }{ s^2+t^2+1 } }, -1+ { \frac{   2  }{ s^2+t^2+1 } }  \right)
} {.}
Kutub Selatan dipetakan ke pusat bidang Euklides oleh pemetaan pertama, sedangkan Kutub Utara dipetakan ke pusat bidang Euklides oleh pemetaan kedua. Kita menyebut kedua pemetaan, maupun pemetaan inversnya, sebagai peta. Kedua peta tersebut bersama-sama menutupi seluruh permukaan bola. Karena merupakan homeomorfisme, keduanya mempertahankan sifat-sifat topologis utama bola. Keduanya tidak cocok untuk geografi Bumi karena peta membutuhkan seluruh bidang dan sangat mendistorsi panjang.

Kedua peta itu sama baiknya. Titik-titik \zusatzklammer {dan secara umum gambar-gambar lain} {} {} pada satu peta dapat dengan mudah diubah ke peta yang lain. Namun, perlu diperhatikan bahwa kedua kutub hanya muncul pada salah satu peta. Titik-titik dari himpunan \mathl{K \setminus \{N,S\}}{} terdapat pada kedua peta, dan melalui keduanya berada dalam bijeksi dengan bidang yang dilubangi di pusatnya, \mathl{\R^2 \setminus \{0\}}{.} Pemetaan \stichwort {transisi} {} \zusatzklammer {atau \stichwort {pergantian peta} {}} {} {} diberikan oleh

\mavergleichskette
{\vergleichskette
{ \psi
}
{ = }{ \beta \circ \left(  \alpha^{-1}\!\mid_{\R^2 \setminus \{0\} }  \right)
}
{  }{
}
{    }{
}
{   }{
}
}
{}{}{}
Selanjutnya,

\mavergleichskettealignhandlinks
{\vergleichskettealignhandlinks
{ \psi(u,v)
}
{ =} { \beta \left(  { \frac{   2u }{ u^2+v^2+1 } }, { \frac{   2v }{ u^2+v^2+1 } } ,1- { \frac{   2  }{ u^2+v^2+1 } }  \right)
}
{ =} { \left(  { \frac{   { \frac{   2u }{ u^2+v^2+1 } }  }{  2 - { \frac{   2  }{ u^2+v^2+1 } }  } } , { \frac{   { \frac{   2v }{ u^2+v^2+1 } }  }{  2 - { \frac{   2  }{ u^2+v^2+1 } }  } }  \right)
}
{ =} { \left(  { \frac{   { \frac{   u  }{ u^2+v^2+1 } }  }{  { \frac{  (u^2+v^2+1) -1  }{ u^2+v^2+1 } }  } } , { \frac{   { \frac{   v  }{ u^2+v^2+1 } }  }{  { \frac{  (u^2+v^2+1) -1  }{ u^2+v^2+1 } }  } }  \right)
}
{ =} { \left(  { \frac{   u  }{ u^2+v^2  } } , { \frac{   v  }{ u^2+v^2 } }  \right)
}
}
{}
{}{.}
Pemetaan transisi ini tidak hanya menginduksi homeomorfisme antara \mathl{\R^2 \setminus \{0\}}{} dengan \mathl{\R^2 \setminus \{0\}}{\zusatzfussnote {Sebaiknya di sini tidak dikatakan \anfuehrung{dengan dirinya sendiri}{} karena kedua bidang peta sebaiknya dibayangkan sebagai bidang yang independen. Hubungan di antara keduanya muncul semata-mata karena keduanya menggambarkan permukaan bola yang sama} {.} {,}} sebagaimana langsung mengikuti dari fakta bahwa kedua pemetaan peta \mathkor {} {\alpha} {dan} {\beta} {} adalah homeomorfisme, melainkan bahkan sebuah \definitionsverweis {difeomorfisme}{}{.} Ini langsung terlihat dari rumusnya; tetapi tidak masuk akal mengatakan bahwa pemetaan peta adalah difeomorfisme, karena permukaan bola bukan himpunan terbuka di dalam $\R^3$. Yang masih kurang adalah sebuah \anfuehrung{struktur diferensiabel}{} pada permukaan ini agar kita dapat berbicara tentang difeomorfisme.

}
Sebuah \stichwort {manifold} {} adalah suatu bangun geometri yang secara \anfuehrung{lokal}{} tampak seperti ruang Euklides $\R^n$. Kita menganggap bangun ini sebagai sebuah
\definitionsverweis {ruang topologis}{}{}
dan memperjelas kata lokal dengan menyatakan bahwa terdapat tutupan oleh himpunan-himpunan terbuka yang homeomorfik dengan himpunan terbuka di $\R^n$. Walaupun selanjutnya kita bekerja dengan ruang topologis, kandungan intuitif pembahasan tidak berkurang jika ruang topologis dibayangkan sekadar sebagai ruang metrik.

  \zwischenueberschrift{Manifold diferensiabel}

\inputdefinition
{}
{

Sebuah
\definitionsverweis {ruang topologis}{}{}
\definitionsverweis {Hausdorff}{}{}
$M$ disebut \definitionswort {manifold topologis}{} berdimensi \definitionswort {dimensi}{} $n$, jika terdapat suatu
\definitionsverweis {tutupan terbuka}{}{}

\mavergleichskette
{\vergleichskette
{ M
}
{ = }{  \bigcup_{i \in I} U_i
}
{  }{
}
{    }{
}
{   }{
}
}
{}{}{}
sedemikian sehingga setiap $U_i$
\definitionsverweis {homeomorf}{}{}
dengan suatu
\definitionsverweis {himpunan bagian terbuka}{}{}
dari $\R^n$.

}

Untuk setiap titik

\mavergleichskette
{\vergleichskette
{P
}
{ \in }{ M
}
{  }{
}
{    }{
}
{   }{
}
}
{}{}{}
terdapat lingkungan terbuka

\mavergleichskette
{\vergleichskette
{P
}
{ \in }{ U
}
{ \subseteq }{ M
}
{    }{
}
{   }{
}
}
{}{}{,}
yang homeomorfik dengan suatu himpunan terbuka

\mavergleichskette
{\vergleichskette
{V
}
{ \subseteq }{ \R^n
}
{  }{
}
{    }{
}
{   }{
}
}
{}{}{}.
Misalkan
\maabbdisp {\varphi} {U} {V
} {}
sebuah homeomorfisme dan misalkan

\mavergleichskette
{\vergleichskette
{ Q
}
{ = }{ \varphi(P)
}
{  }{
}
{    }{
}
{   }{
}
}
{}{}{.}
Maka lingkungan bola terbuka

\mavergleichskette
{\vergleichskette
{Q
}
{ \in }{ U \left(  Q, \epsilon  \right)
}
{ \subseteq }{ V
}
{    }{
}
{   }{
}
}
{}{}{}
berpadanan dengan lingkungan terbuka

\mavergleichskette
{\vergleichskette
{ U'
}
{ = }{  \varphi^{-1}(U \left(  Q, \epsilon  \right))
}
{  }{
}
{    }{
}
{   }{
}
}
{}{}{}
dengan

\mavergleichskette
{\vergleichskette
{P
}
{ \in }{ U'
}
{ \subseteq }{ U
}
{    }{
}
{   }{
}
}
{}{}{,}
yang menurut konstruksi homeomorfik dengan bola terbuka. Karena itu manifold topologis juga dapat dicirikan sebagai ruang Hausdorff topologis yang \stichwort {secara lokal Euklides} {}.

\inputdefinition
{}
{

Misalkan $M$ adalah sebuah
\definitionsverweis {manifold topologis}{}{.}
Setiap
\definitionsverweis {homeomorfisme}{}{}
\maabbdisp {\varphi} { U } { V
} {,}
dengan

\mavergleichskette
{\vergleichskette
{ U
}
{ \subseteq }{ M
}
{  }{
}
{    }{
}
{   }{
}
}
{}{}{}
dan

\mavergleichskette
{\vergleichskette
{ V
}
{ \subseteq }{ \R^n
}
{  }{
}
{    }{
}
{   }{
}
}
{}{}{}
\definitionsverweis {terbuka}{}{}
adalah sebuah
\zusatzklammer {topologis} {} {}
\definitionswort {peta}{} untuk $M$.

}

Himpunan terbuka

\mavergleichskette
{\vergleichskette
{ U
}
{ \subseteq }{ M
}
{  }{
}
{    }{
}
{   }{
}
}
{}{}{}
kadang-kadang disebut \stichwort {daerah peta} {}, sedangkan

\mavergleichskette
{\vergleichskette
{ V
}
{ \subseteq }{ \R^n
}
{  }{
}
{    }{
}
{   }{
}
}
{}{}{}
kadang-kadang disebut \stichwort {citra peta} {.} Untuk sebuah peta
\maabbdisp {\varphi} {U} {V
} {}
dan himpunan bagian terbuka

\mavergleichskette
{\vergleichskette
{ U'
}
{ \subseteq }{ U
}
{  }{
}
{    }{
}
{   }{
}
}
{}{}{}
pemetaan terinduksi
\maabbdisp {\varphi \vert_{U'}} {U'} {\varphi(U')
} {}
juga merupakan peta. Kadang-kadang pemetaan invers juga disebut peta. Selain peta, kita juga berbicara tentang \stichwort {sistem koordinat lokal} {.} Melalui peta
\maabb {\varphi} {U} {V
} {}
koordinat pada

\mavergleichskette
{\vergleichskette
{V
}
{ \subseteq }{ \R^n
}
{  }{
}
{    }{
}
{   }{
}
}
{}{}{}
dipindahkan ke $U$. Koordinat ke-$j$
\zusatzklammer {yakni proyeksi ke-$j$} {} {}
\maabb {x_j} {V} {\R
} {}
menginduksi \zusatzklammer {koordinat lokal} {} {-}fungsi
\maabbdisp {x_j \circ \varphi} {U} {\R
} {}
\zusatzklammer {yang sering kembali dilambangkan dengan $x_j$} {} {,}
dan sebuah titik

\mavergleichskette
{\vergleichskette
{ Q
}
{ = }{ \left(  x_1   , \,   , \ldots ,  , \,   x_n                 \right)
}
{ \in }{ V
}
{    }{
}
{   }{
}
}
{}{}{}
berpadanan dengan titik

\mavergleichskette
{\vergleichskette
{P
}
{ = }{\varphi^{-1}(Q)
}
{  }{
}
{    }{
}
{   }{
}
}
{}{}{.}

\bild{ \begin{center}
\includegraphics[width=5.5cm]{ \bildeinlesungpng {Manifold_zahyou3} {png} }
\end{center}
\bildtext {} }

% TeX-safe display form for the moving list-of-figures argument.  The exact
% Unicode creator string remains in source/unit_media.json and the rights
% manifest; this ASCII rendering avoids an engine-dependent Unicode failure.
\bildlizenz { Manifold zahyou3.png } {} {132ninme} {ja. Wikipedia} {CC-by-sa 3.0} {}

\inputdefinition
{}
{

Misalkan $M$ adalah sebuah
\definitionsverweis {manifold topologis}{}{,}
misalkan

\mavergleichskette
{\vergleichskette
{  U_1, U_2
}
{ \subseteq }{ M
}
{  }{
}
{    }{
}
{   }{
}
}
{}{}{}
merupakan
\definitionsverweis {himpunan bagian terbuka}{}{}
dan
\maabb {\alpha_1} { U_1 } { V_1
} {} dan \maabb {\alpha_2} { U_2 } { V_2
} {}
merupakan
\definitionsverweis {peta}{}{}
\zusatzklammer {dengan

\mavergleichskettek
{\vergleichskettek
{  V_1,V_2
}
{ \subseteq }{\R^n
}
{  }{
}
{    }{
}
{   }{
}
}
{}{}{}
terbuka} {} {.}
Maka
\definitionsverweis {pemetaan}{}{}
\maabbdisp {\alpha_2  \circ \alpha_1^{-1}} { \alpha_1(U_1 \cap U_2) } { \alpha_2(U_1 \cap U_2)
} {} disebut \definitionswort {pemetaan transisi}{} untuk peta-peta tersebut.

}

Irisan \mathl{U_1 \cap U_2}{} adalah himpunan bagian terbuka tempat kedua peta didefinisikan dan tempat keduanya dapat dibandingkan. Secara lebih teliti, pada definisi tersebut kita harus berbicara tentang pembatasan $\alpha_1^{-1}$ pada himpunan bagian terbuka

\mavergleichskette
{\vergleichskette
{ \alpha_1(U_1 \cap U_2)
}
{ \subseteq }{ V_1
}
{  }{
}
{    }{
}
{   }{
}
}
{}{}{}.

\inputdefinition
{}
{

Misalkan

\mavergleichskette
{\vergleichskette
{ n
}
{ \in }{ \N
}
{  }{
}
{    }{
}
{   }{
}
}
{}{}{}
dan

\mavergleichskette
{\vergleichskette
{ k
}
{ \in }{ \overline{ \N }_+
}
{  }{
}
{    }{
}
{   }{
}
}
{}{}{.}
Sebuah
\definitionsverweis {ruang topologis}{}{}
\definitionsverweis {Hausdorff}{}{}
$M$ bersama suatu
\definitionsverweis {tutupan terbuka}{}{}

\mavergleichskette
{\vergleichskette
{ M
}
{ = }{ \bigcup_{i \in I} U_i
}
{  }{
}
{    }{
}
{   }{
}
}
{}{}{}
dan
\definitionsverweis {peta}{}{}
\maabbdisp {\alpha_i} {U_i } { V_i
} {}
dengan

\mavergleichskette
{\vergleichskette
{ V_i
}
{ \subseteq }{ \R^n
}
{  }{
}
{    }{
}
{   }{
}
}
{}{}{}
terbuka, sedemikian sehingga
\definitionsverweis {pemetaan transisi}{}{}
\maabbdisp {\alpha_j \circ (\alpha_i)^{-1}} {V_i \cap \alpha_i(U_i \cap U_j) } { V_j \cap \alpha_j(U_i \cap U_j)
} {}
merupakan
$C^k$-\definitionsverweis {difeomorfisme}{}{}
untuk semua

\mavergleichskette
{\vergleichskette
{ i,j
}
{ \in }{ I
}
{  }{
}
{    }{
}
{   }{
}
}
{}{}{}
dinamakan
\definitionswortpraemath {C^k}{ manifold }{}
atau \definitionswort {manifold diferensiabel}{}
\zusatzklammer {berdimensi \definitionswort {dimensi}{} $n$ dan berderajat diferensiabilitas $k$} {} {.} Kumpulan peta

\mathbed {(U_i,\alpha_i)} {}
 {i \in I} {}
 {} {} {} {,}
dinamakan juga
\definitionswortpraemath {C^k}{ atlas }{}
manifold tersebut.

}

Manifold diferensiabel khususnya merupakan sebuah
\definitionsverweis {manifold topologis}{}{.}
Menurut definisi kita, atlas merupakan bagian integral dari konsep manifold. Namun, yang lebih penting daripada atlas adalah \stichwort {struktur diferensiabel} {} yang ditentukan oleh atlas tersebut. Hal ini akan menjadi lebih jelas setelah kita memiliki konsep pemetaan diferensiabel antara dua manifold dan dapat berbicara tentang manifold yang difeomorf.

\inputdefinition
{}
{

Misalkan $M$ sebuah
\definitionsverweis {manifold diferensiabel}{}{.}
Sebuah
\definitionsverweis {himpunan bagian terbuka}{}{}

\mavergleichskette
{\vergleichskette
{ U
}
{ \subseteq }{ M
}
{  }{
}
{    }{
}
{   }{
}
}
{}{}{,}
yang dilengkapi dengan
\definitionsverweis {peta}{}{}
\definitionsverweis {terbatas}{}{}
disebut \definitionswort {submanifold terbuka}{.}

}

\inputbeispiel{}
{

Setiap himpunan bagian terbuka

\mavergleichskette
{\vergleichskette
{ V
}
{ \subseteq }{ \R^n
}
{  }{
}
{    }{
}
{   }{
}
}
{}{}{}
merupakan
$C^\infty$-\definitionsverweis {manifold}{}{,}
jika kita menggunakan
\definitionsverweis {identitas}{}{}
\maabb {\operatorname{Id}} { V } { V
} {}
sebagai
\definitionsverweis {peta}{}{}.
Satu-satunya
\definitionsverweis {pemetaan transisi}{}{}
adalah identitas itu sendiri, yang merupakan sebuah $C^\infty$-\definitionsverweis {difeomorfisme}{}{}.
Ini adalah manifold dengan sebuah
\definitionsverweis {atlas}{}{,}
yang terdiri atas satu peta. Namun kita juga dapat menggunakan atlas yang terdiri atas semua himpunan bagian terbuka

\mavergleichskette
{\vergleichskette
{ U
}
{ \subseteq }{ V
}
{  }{
}
{    }{
}
{   }{
}
}
{}{}{}
beserta peta identitas terkait $\varphi_U$. Pemetaan transisinya adalah identitas pada \mathl{U_1 \cap U_2}{.}

}

Kita telah membahas ruang metrik yang
\definitionsverweis {terhubung}{}{}
dalam konteks teorema nilai antara. Definisi yang sama juga kita gunakan untuk ruang topologis.

\inputdefinition
{}
{

Sebuah
\definitionsverweis {ruang topologis}{}{}
$X$ disebut \definitionswort {terhubung}{,} jika di dalam $X$ tepat ada dua himpunan
\zusatzklammer {yaitu $\emptyset$ dan seluruh ruang \mathlk{X \neq \emptyset}{}} {} {,}
yang sekaligus
\definitionsverweis {terbuka}{}{}
dan
\definitionsverweis {tertutup}{}{}.

}

Sering kali kita hanya tertarik pada manifold yang terhubung, terutama karena pada kasus tak terhubung komponen-komponen \anfuehrung{keterhubungan}{} dapat diteliti secara terpisah. Sekarang kita membahas secara singkat manifold berdimensi rendah.

\inputbeispiel{}
{

Pada sebuah
\definitionsverweis {manifold}{}{}
$M$ berdimensi nol, untuk setiap titik

\mavergleichskette
{\vergleichskette
{ P
}
{ \in }{ M
}
{  }{
}
{    }{
}
{   }{
}
}
{}{}{}
terdapat lingkungan terbuka

\mavergleichskette
{\vergleichskette
{ P
}
{ \in }{ U
}
{  }{
}
{    }{
}
{   }{
}
}
{}{}{,}
yang
\definitionsverweis {homeomorf}{}{}
dengan suatu himpunan terbuka dari

\mavergleichskette
{\vergleichskette
{ \R^0
}
{ = }{ \{0\}
}
{  }{
}
{    }{
}
{   }{
}
}
{}{}{}.
Artinya, himpunan satu titik \mathl{\{P\}}{} harus terbuka, sehingga $M$ harus membawa
\definitionsverweis {topologi diskret}{}{},
yaitu setiap himpunan bagiannya terbuka. Karena itu satu-satunya manifold berdimensi nol yang
\definitionsverweis {terhubung}{}{}
adalah himpunan satu titik.

}

\bild{ \begin{center}
\includegraphics[width=5.5cm]{ \bildeinlesungsvg {Circle_-_black_simple} {svg} }
\end{center}
\bildtext {Sebuah garis lingkaran adalah manifold kompak berdimensi satu} }

\bildlizenz { Circle - black simple.svg } {} {Dakdada} {Commons} {PD} {}

\inputbeispiel{}
{

Pada manifold berdimensi satu
\definitionsverweis {manifold}{}{}
mula-mula terdapat himpunan bagian terbuka dari $\R^1$. Himpunan-himpunan ini merupakan gabungan interval terbuka, dan \definitionsverweis {terhubung}{}{,} tepat ketika merupakan sebuah interval terbuka. Setiap interval terbuka, baik terbatas maupun tak terbatas, adalah
\definitionsverweis {homeomorf}{}{}
dan juga
\definitionsverweis {difeomorf}{}{}
dengan
\definitionsverweis {interval satuan terbuka}{}{}
\mathl{]0,1[}{} dan dengan bilangan-bilangan riil $\R$ sendiri. Interval tertutup

\mathbed {[a,b]} {dengan}
 {a < b} {}
 {} {} {} {}
bukan manifold, karena untuk titik-titik batasnya
\zusatzklammer {ujung-ujung interval} {} {}
tidak ada lingkungan terbuka yang homeomorfik dengan interval terbuka
\zusatzklammer {meskipun interval itu disebut \stichwort {manifold dengan batas} {}} {} {.}

Selain itu ada \stichwort {lingkaran} {}
\zusatzklammer {atau \stichwort {sfera} {}} {} {}
$S^1$ sebagai manifold terhubung berdimensi satu lainnya. Berlaku

\mavergleichskettedisp
{\vergleichskette
{ S^1
}
{ =} { \left\{ (x,y) \in \R^2  \mid   x^2+y^2 = 1         \right\}
}
{ } {
}
{ } {
}
{ } {
}
}
{}{}{.}
Untuk setiap titik

\mavergleichskette
{\vergleichskette
{ P
}
{ \in }{ S^1
}
{  }{
}
{    }{
}
{   }{
}
}
{}{}{}
berlaku
\mathl{S^1 \setminus \{P \}}{} homeomorfik dengan $\R$
\zusatzklammer {melalui proyeksi stereografis} {} {.}
Lingkaran tidak homeomorfik dengan $\R$, sebab lingkaran
\definitionsverweis {kompak}{}{\zusatzfussnote {Sebelumnya kita mendefinisikan kekompakan hanya untuk himpunan bagian di dalam $\R^n$; sebentar lagi akan terlihat bahwa ini merupakan konsep absolut yang tidak bergantung pada penyematan. Jadi, tidak mungkin menyematkan $\R$ secara homeomorfik ke dalam $\R^n$ sedemikian sehingga citranya kompak} {.} {}}
sedangkan bilangan riil tidak kompak. Selain \mathkor {} {S^1} {dan} {\R} {}, tidak ada manifold terhubung berdimensi satu lainnya yang memiliki
\definitionsverweis {basis topologi terhitung}{}{}
\zusatzklammer {hal ini dinyatakan tanpa bukti di sini} {} {.}

}
Mulai dari dimensi dua, tanpa syarat tambahan yang kuat, mustahil membuat gambaran menyeluruh tentang semua manifold.


\chapter{Lembar Kerja 7}
\setcounter{section}{7}

  \zwischenueberschrift{Soal latihan}

\inputaufgabe
{}
{

Tunjukkan bahwa untuk dua titik
\mathl{A,B \neq N,S}{} pada bola satuan $S^2$, selisih jarak yang diambil dalam dua peta standar stereografis, yaitu
\mathkor {} {d( \alpha_1(A),\alpha_1(B))} {dan} {d( \alpha_2(A),\alpha_2(B))} {,}
dapat dibuat sekecil dan sebesar apa pun.

}
{} {}
Soal di atas menunjukkan bahwa secara umum kita tidak dapat memperoleh pengertian jarak yang bermakna pada suatu manifold melalui peta. Metrik alami pada bola satuan diperoleh dari metrik terinduksi pada $\R^3$ atau dari jarak geodesik; lihat
Aufgabe 7.16.

\inputaufgabe
{}
{

Tunjukkan bahwa suatu
\definitionsverweis {interval setengah terbuka}{}{}
\mathl{[a,b[}{} bukan
\definitionsverweis {manifold topologis}{}{}.

}
{} {}

\inputaufgabe
{}
{

Misalkan
\mathl{I=[0,1[}{} adalah
\zusatzklammer {setengah terbuka ke atas} {} {}
interval satuan dan $S^1$ adalah
\definitionsverweis {lingkaran satuan}{}{.}
Tunjukkan bahwa terdapat suatu
\definitionsverweis {pemetaan bijektif}{}{}
\definitionsverweis {kontinu}{}{}
\maabbdisp {f} { I } { S^1
} {}
tetapi
\mathkor {} {I} {dan} {S^1} {}
tidak
\definitionsverweis {homeomorfik}{}{}
.

}
{} {}

\inputaufgabegibtloesung
{}
{

Tunjukkan bahwa ketiga manifold berdimensi satu
\mathlistdisp {S^1 = \left\{ (x,y) \in \R^2  \mid   \betrag { (x,y) } = 1         \right\}} {} {\R} {dan} {\R \setminus \{0\}} {}
tidak homeomorfik satu sama lain secara berpasangan.

}
{} {}

\inputaufgabe
{}
{

Misalkan $M$ adalah
\definitionsverweis {grafik}{}{}
dari fungsi

\mavergleichskettedisp
{\vergleichskette
{ f(x)
}
{ =} {\begin{cases} \sin \frac{1}{x} \text{ untuk } x \neq 0\, , \\ 0 \text{ selain itu }\, , \end{cases}
}
{ } {
}
{ } {
}
{ } {
}
}
{}{}{}
yang diberikan oleh pemetaan
\maabb {f} {\R} {\R
} {.}
Tunjukkan bahwa $M$ bukan
\definitionsverweis {manifold topologis}{}{}.

}
{} {}

\inputaufgabe
{}
{

Tunjukkan bahwa sebuah
\definitionsverweis {bola terbuka}{}{}

\mavergleichskette
{\vergleichskette
{ U \left(   P ,r  \right)
}
{ \subseteq }{ \R^m
}
{  }{
}
{    }{
}
{   }{
}
}
{}{}{}
adalah $C^\infty$-\definitionsverweis {difeomorfik}{}{}
dengan $\R^m$.

}
{} {}

\inputaufgabegibtloesung
{}
{

Misalkan

\mathdisp {S= \left\{ P \in \R^3  \mid   \Vert { P } \Vert = 1         \right\}} {  }

adalah bola satuan. Untuk
\mathl{v=(a,b,c) \neq 0}{} berlaku

\mathdisp {E_v = \left\{ (x,y,z) \in \R^3  \mid  ax+by+cz = 0         \right\}} {  }

sebuah bidang melalui titik nol, yang mendefinisikan sebuah lingkaran besar
\zusatzklammer {sebuah \anfuehrung{khatulistiwa}{}} {} {}
dan dua hemisfer terbuka pada $S$.

a) Jelaskan, untuk
\mathl{v=(a,b,c) \neq 0}{}, lingkaran besar terkait dan kedua hemisfer dengan persamaan atau pertidaksamaan.

b) Tunjukkan bahwa $S$ tidak dapat ditutupi oleh tiga hemisfer terbuka.

}
{} {}

\inputaufgabe
{}
{

Tunjukkan bahwa permukaan silinder terbuka
\mathl{S^1 \times {]0,1[}}{} homeomorfik dengan
\mathl{S^1 \times \R}{,} dengan bidang yang ditusuk
\mathl{\R^2 \setminus \{(0,0)\}}{} dan dengan
\mathl{S^2 \setminus \{N,S\}}{}
\definitionsverweis {homeomorfik}{}{.}

}
{} {}

\inputaufgabe
{}
{

Misalkan
\mathl{]a,b[}{} sebuah
\definitionsverweis {interval terbuka}{}{}
dan
\maabbdisp {f} {]a,b[} {\R_{\geq 0}
} {}
sebuah
\definitionsverweis {fungsi kontinu}{}{.}
Misalkan $M$ adalah permukaan luar benda hasil rotasi terkait. Tunjukkan bahwa ketika

\mavergleichskette
{\vergleichskette
{ f
}
{ > }{ 0
}
{  }{
}
{    }{
}
{   }{
}
}
{}{}{}
ia merupakan manifold yang homeomorfik dengan selubung silinder terbuka. Tunjukkan lebih lanjut bahwa tidak ada manifold jika $f$ memiliki baik titik nol maupun nilai positif.

}
{} {}

\inputaufgabe
{}
{

Tunjukkan bahwa sebuah
\definitionsverweis {manifold topologis}{}{}
\definitionsverweis {terhubung}{}{}
tepat jika dan hanya jika
\definitionsverweis {terhubung lintasan}{}{}
.

}
{} {}

\inputaufgabe
{}
{

Misalkan $M$ sebuah
\definitionsverweis {manifold topologis}{}{}
dan misalkan

\mathl{M= \bigcup_{i \in I} U_i}{} sebuah
\definitionsverweis {tutupan terbuka}{}{}
dengan
\definitionsverweis {peta}{}{}
\maabbdisp {\alpha_i} {U_i } { V_i
} {}
dengan
\mathl{V_i \subseteq \R^n}{.} Untuk
\mathl{i,j \in I}{} definisikan
\maabbdisp {\varphi_{ij} = \alpha_j \circ (\alpha_i)^{-1}} {V_i \cap \alpha_i(U_i \cap U_j)  } { V_j \cap \alpha_j(U_i \cap U_j)
} {}
yang disebut
\definitionsverweis {pemetaan transisi}{}{.}
Tunjukkan bahwa untuk
\mathl{i,j,k \in I}{} apa yang disebut \stichwort {syarat koklus} {}
\zusatzklammer {pada himpunan bagian mana} {?} {}

\mavergleichskettedisp
{\vergleichskette
{ \varphi_{ij}
}
{ =} {\varphi_{kj} \circ \varphi_{ik}
}
{ } {
}
{ } {
}
{ } {
}
}
{}{}{}
berlaku.

}
{} {}

\inputaufgabe
{}
{

Periksa syarat koklus
\zusatzklammer {lihat
Aufgabe 7.11} {} {}
untuk bola satuan $M=S^2$ dan tiga proyeksi stereografis dari Kutub Utara, Kutub Selatan, dan
\mathl{P=(1,0,0)}{.}

}
{} {}

\inputaufgabegibtloesung
{}
{

Buat sebuah animasi yang menampilkan objek-objek geometri dari
Aufgabe 7.17.

}
{} {}

\inputaufgabe
{}
{

Perhatikan kertas kado. Dengan cara apa kertas tersebut dapat dipotong dan/atau direkatkan sehingga menghasilkan manifold berdimensi dua? Apakah tepi kertas harus menjadi bagian darinya atau tidak? Manifold mana yang dihasilkan yang terhubung, mana yang kompak? Bagaimana pita Möbius dibuat? Kemungkinan apa yang ada jika kita mengizinkan sejumlah hingga titik pengecualian, tempat tidak ada struktur manifold?

Terapkan teori ini untuk merancang kemasan yang seorisinal mungkin. Ikatlah kemasan tersebut dengan manifold berdimensi satu yang sesuai.

}
{} {}

  \zwischenueberschrift{Soal untuk dikumpulkan}

\inputaufgabe
{3}
{

Tentukan
\definitionsverweis {citra}{}{} dari
\definitionsverweis {lingkaran-lingkaran besar}{}{}
melalui kedua kutub pada
\definitionsverweis {bola satuan}{}{} di bawah
\definitionsverweis {proyeksi stereografis}{}{}
dari Kutub Utara.

}
{} {}

\inputaufgabe
{5}
{

Tunjukkan bahwa pada
\definitionsverweis {bola satuan}{}{}

\mavergleichskette
{\vergleichskette
{ K
}
{ \subset }{ \R^3
}
{  }{
}
{    }{
}
{   }{
}
}
{}{}{}
melalui korespondensi berikut ditentukan sebuah
\definitionsverweis {metrik}{}{} . Untuk

\mavergleichskette
{\vergleichskette
{ P,Q
}
{ \in }{ K
}
{  }{
}
{    }{
}
{   }{
}
}
{}{}{}
\mathl{d(P,Q)}{} adalah panjang lintasan penghubung
\zusatzklammer {yang lebih pendek} {} {} dari $P$ ke $Q$ pada
\definitionsverweis {lingkaran besar}{}{}
\zusatzklammer {perhatikan juga kasus

\mavergleichskette
{\vergleichskette
{ P
}
{ = }{ Q
}
{  }{
}
{    }{
}
{   }{
}
}
{}{}{}
dan $P,Q$
\definitionsverweis {antipodal}{}{}} {} {.}

}
{} {}

\inputaufgabe
{8}
{

Kita tetapkan dua titik $N=(0,0,1)$ dan $P=(1,0,0)$ pada
\definitionsverweis {bola satuan}{}{} $K$. Misalkan $G$ adalah garis penghubung dan $H$ adalah bidang tegak lurus terhadap $G$ melalui $N$. Pada $H$, buat lingkaran satuan berparameter $E$ dengan $N$ sebagai pusat. Tentukan, untuk $S \in E$, panjang lintasan
\zusatzklammer {yang lebih pendek} {} {} dari $N$ ke $P$ pada lingkaran yang ditentukan oleh irisan $K$ dengan bidang yang melalui
\mathkor {} {N,P} {dan} {S} {},
tersebut.

}
{} {}

\inputaufgabe
{6}
{

Berikan suatu
\definitionsverweis {pemetaan injektif}{}{}
\definitionsverweis {kontinu}{}{}
\maabbdisp {\varphi} {\R} {S^2
} {,}
yang
\zusatzklammer {sebagai pemetaan ke $\R^3$} {} {}
\definitionsverweis {dapat direktifikasi}{}{}
memiliki
\definitionsverweis {panjang}{}{}
tak hingga, dan memenuhi
\mathkor {} {\operatorname{lim}_{   t  \rightarrow  \infty  } \, \varphi (t) = N} {dan} {\operatorname{lim}_{   t  \rightarrow  - \infty  } \, \varphi (t) = S} {}.

}
{} {}

\inputaufgabe
{4}
{

Tunjukkan bahwa sumbu-sumbu koordinat bukan
\definitionsverweis {manifold topologis}{}{}.

}
{} {}


\section*{Solusi yang disediakan oleh sumber}
\addcontentsline{toc}{section}{Solusi yang disediakan oleh sumber}
\subsection*{Solusi Soal 7.4}
Sebagai himpunan bagian tertutup dan terbatas dari $\R^2$, lingkaran satuan kompak
\zusatzklammer {Satz von Heine-Borel} {} {.} Garis riil $\R$ dan $\R \setminus \{0\}$, sebaliknya, tidak kompak karena keduanya merupakan himpunan bagian tak terbatas dari $\R$. Jadi
\mathl{S^1 \not \cong \R,\, \R \setminus \{0\}}{.}

Garis riil terhubung, sebagaimana mengikuti dari
Zwischenwertsatz.
Sebaliknya, $\R \setminus \{0\}$ tidak terhubung, karena kita dapat menuliskannya sebagai

\mavergleichskettedisp
{\vergleichskette
{\R_+
}
{ =} { (\R \setminus \{0\}) \cap [0, \infty]
}
{ =} { (\R \setminus \{0\}) \cap \, ]0, \infty]
}
{ } {
}
{ } {
}
}
{}{}{}
sehingga diperoleh suatu himpunan bagian tak kosong yang sekaligus terbuka dan tertutup. Jadi
\mathl{\R \,\not \cong \R \setminus \{0\}}{.}

\subsection*{Solusi Soal 7.7}
a) Lingkaran besar adalah

\mathdisp {G= \left\{ (x,y,z) \in S  \mid  ax+by+cz = 0        \right\}} {  }

dan kedua hemisfer terbuka adalah

\mathdisp {U_+ = \left\{ (x,y,z) \in S  \mid  ax+by+cz > 0        \right\}} {  }

dan

\mathdisp {U_- = \left\{ (x,y,z) \in S  \mid  ax+by+cz < 0        \right\}} { . }

b) Untuk setiap hemisfer terbuka $U$ dan lingkaran besar $G$ yang terkait berlaku
\mathl{U \cap G = \emptyset}{.} Jarak maksimum antara dua titik
\mathl{P,Q \in S}{} adalah
\mathl{2}{,} dan hal ini terjadi tepat ketika kedua titik tersebut saling berhadapan
\zusatzklammer {antipodal} {} {}, yaitu ketika garis penghubungnya melalui pusat bola
\zusatzklammer {yakni pada
\mathl{Q=-P}{}} {} {}. Pasangan antipodal seperti itu tidak terletak pada satu hemisfer terbuka, sebab untuk
\mathl{P=(x,y,z)}{} dan
\mathl{P \in U_+}{} berlaku
\mathl{ax+by+cz>0}{} dan karena itu
\mathl{a(-x) +b(-y)+c(-z) <0}{} berlaku, sehingga
\mathl{-P \in U_-}{.}

Sekarang anggaplah bahwa terdapat suatu tutupan terbuka

\mathdisp {S \subseteq U_1 \cup U_2 \cup U_3} {  }

dengan tiga hemisfer terbuka
\mathl{U_1,U_2,U_3}{}
\zusatzklammer {bidang dan lingkaran besar yang bersesuaian masing-masing dilambangkan dengan
\mathkor {} {E_i} {atau} {G_i} {}} {} {.}
Karena
\mathl{U_1 \cap G_1 = \emptyset}{} diperoleh

\mathdisp {G_1 \subseteq U_2 \cup U_3} { . }

Irisan
\mathl{G_1 \cap E_2}{} memuat setidaknya dua titik antipodal
\mathkor {} {P} {dan} {-P} {}
dan
\mathl{P,-P \not\in U_2}{.} Karena, berdasarkan pengamatan sebelumnya, $U_3$ tidak memuat pasangan titik antipodal, salah satu dari kedua titik ini juga tidak termasuk dalam $U_3$, sehingga kita memperoleh kontradiksi.

\subsection*{Solusi Soal 7.13}
Animasi pertama memperlihatkan secara bertahap konstruksi objek-objek individual. Animasi lainnya memperlihatkan variasi S dan perubahan panjang lintasan yang bersesuaian.

Dari Lukas Freudenberg.
\href{https://commons.wikimedia.org/wiki/File:Aufgabe75.22.1.gif}{Konstruktion der Objekte}
\href{https://commons.wikimedia.org/wiki/File:Aufgabe75.22.2.gif}{Variation von S}


\part{Unit 8}
\chapter[Kuliah 8]{Kuliah 8: Teorema Pemetaan Implisit dan Manifold}
\input{generated/lecture08.id.build.tex}

\chapter{Lembar Kerja 8}
\setcounter{section}{8}

  \zwischenueberschrift{Soal latihan}

\inputaufgabe
{}
{

Misalkan
\mathkor {} {M} {dan} {N} {}
adalah
\definitionsverweis {manifold diferensiabel}{}{}
dan
\maabbdisp {\varphi} { M } { N
} {}
sebuah pemetaan. Misalkan
\mathl{M= \bigcup_{i \in I} U_i}{} suatu
\definitionsverweis {tutupan terbuka}{}{}
dari $M$. Tunjukkan bahwa $\varphi$
\definitionsverweis {diferensiabel}{}{}
jika dan hanya jika semua pembatasan
\mathl{\varphi_i= \varphi \vert_{U_i }}{} diferensiabel.

}
{} {}

\inputaufgabe
{}
{

Misalkan $M$ sebuah
\definitionsverweis {manifold diferensiabel}{}{}
dan
\maabbdisp {\alpha} { U } { V
} {}
sebuah peta
\zusatzklammer {jadi
\mathl{U \subseteq M}{} dan
\mathl{V \subseteq \R^n}{} terbuka} {} {.}
Tunjukkan bahwa $\alpha$ adalah sebuah
\definitionsverweis {difeomorfisme}{}{}.

}
{} {}

\inputaufgabe
{}
{

Misalkan $M$ sebuah
\definitionsverweis {manifold diferensiabel}{}{} dan
\maabbdisp {f,g} { M } {\R
} {}
\definitionsverweis {fungsi-fungsi diferensiabel}{}{} pada $M$. Buktikan pernyataan-pernyataan berikut. \aufzaehlungvier{Pemetaan
\maabbeledisp {f \times g} { M } {\R^2
} { x } {(f(x),g(x))
} {,} diferensiabel.
}{ $f+g$
\definitionsverweis {diferensiabel}{}{.}
}{ $f \cdot g$
\definitionsverweis {diferensiabel}{}{.}
}{Jika $f$ tidak memiliki titik nol, maka
\mathl{f^{-1}}{} juga diferensiabel.
}

}
{} {}

\inputaufgabe
{}
{

Misalkan
\mathl{S^2 \subseteq \R^3}{} adalah
\definitionsverweis {sfera}{}{.}
Dengan menggunakan
\definitionsverweis {peta-peta stereografis}{}{,}
tunjukkan bahwa pembatasan koordinat ruang
\mathl{x,y,z}{} pada sfera merupakan
\definitionsverweis {fungsi-fungsi diferensiabel}{}{}.

}
{} {}

\inputaufgabe
{}
{

Misalkan
\mathkor {} {M} {dan} {N} {}
adalah
\definitionsverweis {manifold diferensiabel}{}{}
dan
\maabbdisp {\varphi} { M } { N
} {}
sebuah
\definitionsverweis {pemetaan diferensiabel}{}{.}
Tunjukkan bahwa pemetaan ini menginduksi sebuah
\definitionsverweis {homomorfisme gelanggang}{}{}
\maabbeledisp {\varphi^*} {C^1(N,\R)} {C^1(M,\R)
} { f } {f \circ \varphi
} {,}
melalui penarikan balik.

}
{} {}

\inputaufgabe
{}
{

Tunjukkan bahwa untuk
\mathl{m \leq n}{}, penyematan subruang
\mathl{\R^m}{} ke dalam $\R^n$ yang diberikan oleh
\mathl{(x_1 , \ldots , x_m) \mapsto (x_1 , \ldots , x_m,0 , \ldots , 0)}{}
\definitionsverweis {diferensiabel hingga orde berapa pun}{}{}.

}
{} {}

\inputaufgabe
{}
{

Tunjukkan bahwa permukaan silinder terbuka
\mathl{S^1 \times {]0,1[}}{},
\mathl{S^1 \times \R}{,} bidang yang ditusuk
\mathl{\R^2 \setminus \{(0,0)\}}{}, dan
\mathl{S^2 \setminus \{N,S\}}{}
saling \definitionsverweis {difeomorfik}{}{}.

}
{} {}

Soal berikut menggunakan definisi di bawah ini.

Sebuah fungsi
\maabbdisp {f} { {\mathbb K}^n } { {\mathbb K}
} {}
disebut
\definitionswort {homogen berderajat}{}
$d$, jika untuk setiap titik
\mathl{P \in {\mathbb K}^n}{} dan setiap
\mathl{\lambda \in {\mathbb K}}{} berlaku hubungan

\mavergleichskettedisp
{\vergleichskette
{f(\lambda P)
}
{ =} { \lambda^d f(P)
}
{ } {
}
{ } {
}
{ } {
}
}
{}{}{}
untuk semua pilihan tersebut.

\inputaufgabe
{}
{

Misalkan
\maabbdisp {f} {\R^n } {\R
} {}
sebuah
\definitionsverweis {fungsi homogen}{}{} yang
\definitionsverweis {terdiferensialkan secara kontinu}{}{,}
dan yang pada
\definitionsverweis {serat}{}{}
$F$ di atas
\mathl{a \neq 0}{}
bersifat \definitionsverweis {reguler}{}{.}
Tunjukkan bahwa setiap serat di atas
\mathl{b \neq 0,\ b=\lambda^d a}{}
untuk suatu skalar real tak nol,
\definitionsverweis {difeomorfik}{}{} dengan $F$ dan merupakan sebuah
\definitionsverweis {manifold}{}{}.

Catatan edisi: Syarat penskalaan real ditambahkan. Klaim sumber untuk semua nilai tak nol gagal pada derajat genap ketika tanda nilai serat berbeda; misalnya, fungsi kuadrat norma mempunyai serat positif berbentuk bola tetapi serat negatif kosong.

}
{} {}

\inputaufgabe
{}
{

Misalkan
\mathl{]a,b[}{} sebuah
\definitionsverweis {interval terbuka}{}{}
dan
\maabbdisp {f} {]a,b[} {\R_{+}
} {}
sebuah
\definitionsverweis {fungsi yang terdiferensialkan secara kontinu}{}{.}
Misalkan $M$ adalah permukaan benda putar terkait. Tunjukkan bahwa himpunan ini merupakan manifold yang difeomorfik dengan silinder terbuka.

}
{} {}

\inputaufgabe
{}
{

Tunjukkan bahwa permukaan elipsoid dan permukaan bola satuan
$C^\infty$-\definitionsverweis {difeomorfik}{}{}.

}
{} {}

\inputaufgabegibtloesung
{}
{

Berikan sebuah contoh
\definitionsverweis {manifold diferensiabel}{}{}
\definitionsverweis {berdimensi dua}{}{} yang
\definitionsverweis {terhubung}{}{},
sebutlah $M$, dan sebuah titik

\mavergleichskette
{\vergleichskette
{ P
}
{ \in }{ M
}
{  }{
}
{    }{
}
{   }{
}
}
{}{}{}
sedemikian sehingga
\mathkor {} {M} {dan} {M \setminus \{P\}} {}
\definitionsverweis {difeomorfik}{}{}
satu sama lain.

}
{} {}

\inputaufgabe
{}
{

Tunjukkan bahwa
\definitionsverweis {pemetaan antipodal}{}{}
\maabbeledisp {} {S^n } { S^n
} {(x_1 , \ldots , x_{n+1})} {(-x_1 , \ldots , -x_{n+1})
} {,}
adalah sebuah
\definitionsverweis {difeomorfisme}{}{} \definitionsverweis {tanpa titik tetap}{}{}
yang inversnya adalah dirinya sendiri.

}
{} {}

\inputaufgabegibtloesung
{}
{

Tunjukkan bahwa sebuah
\definitionsverweis {manifold diferensiabel}{}{}
$M$ berdimensi

\mavergleichskette
{\vergleichskette
{ n
}
{ \geq }{ 1
}
{  }{
}
{    }{
}
{   }{
}
}
{}{}{}
memiliki tak hingga banyak
\definitionsverweis {difeomorfisme}{}{}
\maabbdisp {\varphi} { M } { M
} {}
yang berbeda.

}
{} {}

Soal-soal berikut dimaksudkan untuk menjelaskan mengapa manifold dirumuskan dengan tutupan terbuka.

Misalkan
\mathl{\left(  x_n  \right)_{n \in  \N  }}{} sebuah
\definitionsverweis {barisan}{}{}
dalam suatu
\definitionsverweis {ruang topologis}{}{}
$X$. Kita mengatakan bahwa, untuk suatu

\mavergleichskette
{\vergleichskette
{ x
}
{ \in }{ X
}
{  }{
}
{    }{
}
{   }{
}
}
{}{}{}
barisan tersebut \definitionswort {konvergen ke titik itu}{,} jika sifat berikut terpenuhi.

Untuk setiap
\definitionsverweis {lingkungan terbuka}{}{}

\mavergleichskette
{\vergleichskette
{ U
}
{ \subseteq }{ X
}
{  }{
}
{    }{
}
{   }{
}
}
{}{}{}
yang memuat $x$, terdapat suatu

\mavergleichskette
{\vergleichskette
{ n_0
}
{ \in }{ \N
}
{  }{
}
{    }{
}
{   }{
}
}
{}{}{}
sedemikian sehingga untuk semua

\mavergleichskette
{\vergleichskette
{ n
}
{ \geq }{ n_0
}
{  }{
}
{    }{
}
{   }{
}
}
{}{}{}
suku barisan $x_n$ termasuk dalam $U$.

Dalam hal ini $x$ disebut \definitionswort {nilai batas}{} atau \definitionswort {limit}{} barisan tersebut. Kita juga menuliskannya sebagai

\mavergleichskettedisp
{\vergleichskette
{  \lim_{n \rightarrow \infty}  x_n
}
{ =} { x
}
{ } {
}
{ } {
}
{ } {
}
}
{}{}{.}
Jika barisan memiliki limit, kita juga mengatakan bahwa barisan itu \definitionswort {konvergen}{}
\zusatzklammer {tanpa menyebut nilai limit tertentu} {} {,}
dan jika tidak, bahwa barisan itu \definitionswort {divergen}{.}

\inputaufgabe
{}
{

Misalkan $M$ sebuah
\definitionsverweis {manifold topologis}{}{}
dan
\mathl{\left(   x_n  \right)_{n \in \N }}{} sebuah
\definitionsverweis {barisan}{}{}
dalam $M$. Tunjukkan bahwa barisan tersebut
\definitionsverweis {konvergen}{}{,}
jika dan hanya jika terdapat sebuah
\definitionsverweis {domain peta}{}{}
dari $M$ yang memuat hampir semua suku barisan dan sedemikian sehingga barisan citra yang bersesuaian konvergen di dalam citra peta.

}
{} {}

\inputaufgabe
{}
{

Misalkan $M$ sebuah
\definitionsverweis {manifold topologis}{}{}
dan
\maabbdisp {\alpha_i} {U_i } { V_i
} {}
sebuah keluarga
\definitionsverweis {peta}{}{}
dengan
\definitionsverweis {pemetaan transisi}{}{}
\maabbdisp {\varphi_{ij} = \alpha_j \circ (\alpha_i)^{-1}} {V_i \cap \alpha_i(U_i \cap U_j) } { V_j \cap \alpha_j(U_i \cap U_j)
} {.}
Tunjukkan bahwa manifold $M$ dapat direkonstruksi dari keluarga

\mathbed {V_i} {}
 {i \in I} {}
 {} {} {} {,}
himpunan-himpunan bagian
\mathl{V_{ij} \subseteq V_i}{} dan pemetaan transisi
\maabbdisp {\varphi_{ij}} {V_{ij}} { V_{ji}
} {}
yang bersesuaian.

a) Pada

\mavergleichskettedisp
{\vergleichskette
{N
}
{ \defeq} { \biguplus_{i \in I} V_i
}
{ } {
}
{ } {
}
{ } {
}
}
{}{}{}
perhatikan
\definitionsverweis {relasi ekuivalensi}{}{,}
yang menyamakan dua titik
\mathl{P\in V_i}{} dan
\mathl{Q \in V_j}{} jika
\mathl{\varphi_{ij}}{} memetakan titik pertama ke titik kedua.

b) Lengkapi himpunan hasil bagi
\mathl{N/ \sim}{} dengan sebuah topologi yang sesuai.

c) Definisikan peta-peta pada
\mathl{N/ \sim}{}.

d) Tunjukkan bahwa
\mathkor {} {M} {dan} {N/ \sim} {}
\definitionsverweis {homeomorfik}{}{}.

}
{} {}

Prosedur konstruksi manifold yang dijelaskan dalam soal sebelumnya berlaku untuk suatu keluarga himpunan bagian terbuka di dalam $\R^n$ dengan pemetaan transisi yang memenuhi syarat koklus dari
Soal 7.11.
Namun, ruang topologis yang dihasilkan tidak serta-merta merupakan ruang Hausdorff. Konstruksi ini disebut \stichwort {perekatan terbuka} {} ruang-ruang.

\inputaufgabe
{}
{

Kita mengambil dua salinan garis riil, yaitu

\mavergleichskette
{\vergleichskette
{G_1
}
{ = }{\R
}
{  }{
}
{    }{
}
{   }{
}
}
{}{}{}
dan

\mavergleichskette
{\vergleichskette
{G_2
}
{ = }{\R
}
{  }{
}
{    }{
}
{   }{
}
}
{}{}{}
beserta pemetaan perekatan
\maabbeledisp {\varphi} {G_1 \setminus \{0\} } { G_2 \setminus \{0\}
} { x } { x^{-1}
} {.}
Misalkan $M$ ruang topologis yang dihasilkan menurut konstruksi yang dijelaskan dalam
Soal 8.15.
Tunjukkan bahwa $M$
\definitionsverweis {homeomorfik}{}{}
dengan
$1$-\definitionsverweis {sfera}{}{}.

}
{} {}

\inputaufgabe
{}
{

Kita mengambil dua salinan garis riil, yaitu

\mavergleichskette
{\vergleichskette
{G_1
}
{ = }{\R
}
{  }{
}
{    }{
}
{   }{
}
}
{}{}{}
dan

\mavergleichskette
{\vergleichskette
{G_2
}
{ = }{\R
}
{  }{
}
{    }{
}
{   }{
}
}
{}{}{}
beserta pemetaan perekatan
\maabbdisp {\operatorname{Id}} {G_1 \setminus \{0\} } { G_2 \setminus \{0\}
} {.}
Misalkan $M$ ruang topologis yang dihasilkan menurut konstruksi yang dijelaskan dalam
Soal 8.15.
Tunjukkan bahwa $M$ bukan
\definitionsverweis {manifold}{}{}.

}
{} {}

  \zwischenueberschrift{Soal untuk dikumpulkan}

Dalam soal berikut, tafsirkan ${\mathbb C}$ sebagai $\R^2$.

\inputaufgabe
{4}
{

Perhatikan pemetaan
\maabbeledisp {} { {\mathbb C}^2} { {\mathbb C}
} {(z,w)} { zw
} {.}
Untuk titik manakah
\mathl{u \in {\mathbb C}}{} serat di atas $u$ merupakan manifold? Dalam setiap kasus, berikan deskripsi kelas difeomorfisme yang sesederhana mungkin.

}
{} {}

\inputaufgabe
{6}
{

Misalkan diberikan dua titik
\mathkor {} {P} {dan} {Q} {}
pada
\definitionsverweis {permukaan bola satuan}{}{.}
Tunjukkan bahwa terdapat sebuah
\definitionsverweis {difeomorfisme}{}{}
dari permukaan bola itu ke dirinya sendiri yang memetakan $P$ ke $Q$.

}
{} {}

\inputaufgabe
{4 (1+1+2)}
{

Misalkan $M$ sebuah
\definitionsverweis {manifold diferensiabel}{}{.}
Untuk setiap himpunan bagian terbuka

\mavergleichskette
{\vergleichskette
{ U
}
{ \subseteq }{ M
}
{  }{
}
{    }{
}
{   }{
}
}
{}{}{}
kita perhatikan himpunan
\mathl{C^1(U,\R)}{} yang terdiri atas
\definitionsverweis {fungsi-fungsi diferensiabel}{}{}
pada $U$. Misalkan

\mavergleichskette
{\vergleichskette
{ M
}
{ = }{ \bigcup_{i \in I} U_i
}
{  }{
}
{    }{
}
{   }{
}
}
{}{}{}
sebuah
\definitionsverweis {tutupan terbuka}{}{.}
\aufzaehlungdreiabc{Tunjukkan bahwa jika

\mavergleichskette
{\vergleichskette
{ V
}
{ \subseteq }{ U
}
{  }{
}
{    }{
}
{   }{
}
}
{}{}{}
terbuka dan

\mavergleichskette
{\vergleichskette
{ f
}
{ \in }{ C^1(U,\R)
}
{  }{
}
{    }{
}
{   }{
}
}
{}{}{}
maka
\definitionsverweis {pembatasan}{}{}
$f \vert_V$ juga termasuk dalam
\mathl{C^1(V,\R)}{}.
}{Misalkan

\mavergleichskette
{\vergleichskette
{ f
}
{ \in }{ C^1(M,\R)
}
{  }{
}
{    }{
}
{   }{
}
}
{}{}{.}
Tunjukkan bahwa

\mavergleichskette
{\vergleichskette
{ f
}
{ = }{ 0
}
{  }{
}
{    }{
}
{   }{
}
}
{}{}{}
jika dan hanya jika semua pembatasan

\mavergleichskette
{\vergleichskette
{ f \vert_{U_i }
}
{ = }{ 0
}
{  }{
}
{    }{
}
{   }{
}
}
{}{}{}
berlaku.
}{Misalkan diberikan suatu keluarga

\mavergleichskette
{\vergleichskette
{ f_i
}
{ \in }{ C^1(U_i,\R)
}
{  }{
}
{    }{
}
{   }{
}
}
{}{}{}
fungsi-fungsi yang memenuhi \anfuehrung{syarat kecocokan}{}

\mavergleichskettedisp
{\vergleichskette
{ f_i \vert_{U_i \cap U_j }
}
{ =} { f_j \vert_{U_i \cap U_j }
}
{ } {
}
{ } {
}
{ } {
}
}
{}{}{}
untuk semua $i,j$. Tunjukkan bahwa terdapat suatu

\mavergleichskette
{\vergleichskette
{ f
}
{ \in }{ C^1(M,\R)
}
{  }{
}
{    }{
}
{   }{
}
}
{}{}{}
yang memenuhi

\mavergleichskette
{\vergleichskette
{ f \vert_{U_i }
}
{ = }{ f_i
}
{  }{
}
{    }{
}
{   }{
}
}
{}{}{}
untuk semua $i$.
}

}
{} {}

\inputaufgabe
{5}
{

Misalkan $M$ sebuah
\definitionsverweis {manifold diferensiabel}{}{,} yang memiliki setidaknya dua titik. Tunjukkan bahwa terdapat
\definitionsverweis {fungsi-fungsi diferensiabel}{}{}
\maabbdisp {f,g} { M } {\R
} {}
dengan
\mathl{f,g \neq 0}{,} tetapi
\mathl{fg=0}{.}

}
{} {}


\section*{Solusi yang disediakan oleh sumber}
\addcontentsline{toc}{section}{Solusi yang disediakan oleh sumber}
\subsection*{Solusi Soal 8.11}
Misalkan

\mavergleichskettedisp
{\vergleichskette
{M
}
{ =} { \R^2 \setminus \N_+
}
{ =} { \R^2 \setminus \{ (1,0) , (2,0), (3,0), \ldots  \}
}
{ } {
}
{ } {
}
}
{}{}{,}
yakni dari $\R^2$ kita mengeluarkan bilangan-bilangan asli positif yang terletak pada sumbu $x$. Himpunan ini terbuka di dalam $\R^2$ dan karena itu merupakan manifold diferensiabel. Manifold ini terhubung lintasan: misalnya, setiap titik di $M$ dengan
\mathl{y \geq 0}{} dapat dihubungkan oleh lintasan lurus dengan
\mathl{(0,1)}{}, dan setiap titik di $M$ dengan
\mathl{y \leq 0}{} dapat dihubungkan oleh lintasan lurus dengan
\mathl{(0,-1)}{}; kedua titik tersebut juga dapat dihubungkan oleh lintasan lurus. Sekarang misalkan

\mavergleichskettedisp
{\vergleichskette
{P
}
{ =} { (0,0)
}
{ } {
}
{ } {
}
{ } {
}
}
{}{}{}
dan

\mavergleichskettedisp
{\vergleichskette
{M'
}
{ =} { M \setminus \{P\}
}
{ } {
}
{ } {
}
{ } {
}
}
{}{}{.}
Pemetaan
\maabbeledisp {} {\R^2} {\R^2
} {(x,y)} {(x-1,y)
} {,}
merupakan sebuah difeomorfisme
\zusatzklammer {sebagai translasi} {} {},
dan citra $M$ tepat sama dengan $M'$. Oleh karena itu,
\mathkor {} {M} {dan} {M'=M \setminus \{P\}} {}
difeomorfik satu sama lain.

\subsection*{Solusi Soal 8.13}
\input{generated/worksheet08_exercise13_solution.id.build.tex}

\part{Unit 9}
\chapter[Kuliah 9]{Kuliah 9: Ruang Singgung Manifold Diferensiabel}
\input{generated/lecture09.id.build.tex}

\chapter{Lembar Kerja 9}
\input{generated/worksheet09.id.build.tex}

\part{Unit 10}
\chapter[Kuliah 10]{Kuliah 10: Submanifold Tertutup}
\input{generated/lecture10.id.build.tex}

\chapter{Lembar Kerja 10}
\input{generated/worksheet10.id.build.tex}

\section*{Solusi yang disediakan oleh sumber}
\addcontentsline{toc}{section}{Solusi yang disediakan oleh sumber}
\subsection*{Solusi Soal 10.9}
\input{generated/worksheet10_exercise09_solution.id.build.tex}
\subsection*{Solusi Soal 10.10}
\input{generated/worksheet10_exercise10_solution.id.build.tex}
\subsection*{Solusi Soal 10.15}
Dengan regularitas $C^2$ yang dinyatakan dalam soal edisi ini, kita klaim

\mavergleichskettealignhandlinks
{\vergleichskettealignhandlinks
{ TY
}
{ =} { \left\{  (P,u)   \mid   h(P) = c , \,  \left(  D h   \right)_{ P } \left(  u  \right) = 0       \right\}
}
{ =} { \left\{  (P,u_1 , \ldots , u_n)   \mid   h(P) = c , \,  \sum_{i = 1}^n u_i { \frac{   \partial   h   }{  \partial   x_i   } } (P) = 0       \right\}
}
{ \subseteq} { W \times \R^n
}
{ } {
}
}
{}
{}{}
dengan proyeksi alami ke $Y$.

Definisikan pemetaan
\maabbeledisp {\Phi} {W\times\R^n} {\R^2
} {(P,u)} {\left(h(P),\left(Dh\right)_P(u)\right)
} {.}
Karena $h$ berkelas $C^2$, pemetaan $\Phi$ berkelas $C^1$, dan

\mathdisp {TY=\Phi^{-1}(c,0)} { . }

Kita tunjukkan bahwa $(c,0)$ merupakan nilai reguler. Untuk
\mathl{(P,u)\in TY}{}, diferensial $\Phi$ diberikan oleh

\mathdisp {(D\Phi)_{(P,u)}(v,w)=\left((Dh)_P(v),\,(D^2h)_P(v,u)+(Dh)_P(w)\right)} { . }

Karena $h$ reguler pada $Y$, bentuk linear $(Dh)_P$ surjektif. Untuk sebarang
\mathl{(r,s)\in\R^2}{}, pilih $v$ dengan $(Dh)_P(v)=r$, lalu pilih $w$ dengan

\mathdisp {(Dh)_P(w)=s-(D^2h)_P(v,u)} { . }

Dengan pilihan ini, $(D\Phi)_{(P,u)}(v,w)=(r,s)$, sehingga diferensial $\Phi$ surjektif di setiap titik serat. Teorema pemetaan implisit kini menunjukkan bahwa $TY$ merupakan submanifold tertutup dari $W\times\R^n$ berdimensi $2n-2$. Ketertutupannya juga langsung mengikuti dari kontinuitas $\Phi$ dan ketertutupannya titik $(c,0)$ di $\R^2$.

Untuk mencocokkan realisasi tertanam ini dengan struktur bundel, bagi setiap
\mathl{P\in Y}{}
pilih suatu parametrisasi lokal
\maabbdisp {\beta} {V} {U\subseteq Y
} {,}
dengan $V\subseteq\R^{n-1}$ terbuka. Pemetaan
\maabbeledisp {} {V \times \R^{n-1} } { U \times \R^n
} {(Q,z)} { ( \beta(Q), \left(  D\beta  \right)_{ Q } \left(  z  \right) )
} {,}
memberikan trivialisasi lokal yang citranya tepat terdiri atas pasangan titik dan vektor singgung pada serat. Jadi topologi serta proyeksi realisasi tertanam ini sama dengan topologi dan proyeksi bundel singgung yang telah didefinisikan.

Catatan edisi: Solusi sumber mengganti syarat $h(P)=c$ dengan $h(P)=0$, menulis turunan parsial dari fungsi $f$ yang tidak didefinisikan, dan memakai $\beta(q)$ alih-alih $\beta(Q)$. Lebih mendasar, asumsi $C^1$ sumber hanya membuat turunan pertama kontinu dan tidak cukup untuk menyimpulkan bahwa $TY$ adalah submanifold diferensiabel dari $W\times\R^n$. Edisi ini memperkuat asumsi soal menjadi $C^2$ dan memberikan pemeriksaan nilai reguler yang hilang.

\subsection*{Solusi Soal 10.25}
Misalkan $N=(0,0,1)$ dan gunakan invers proyeksi stereografis
\maabbeledisp {\sigma} {\R^2} {S^2\setminus\{N\}
} {(x,y)} {\left(\frac{2x}{1+x^2+y^2},\frac{2y}{1+x^2+y^2},\frac{x^2+y^2-1}{1+x^2+y^2}\right)
} {.}
Pada $S^2\setminus\{N\}$, dorong medan vektor konstan $e_1$ ke depan dan definisikan

\mathdisp {X(\sigma(x,y))=(D\sigma)_{(x,y)}(e_1)} { . }

Secara eksplisit,

\mathdisp {X(\sigma(x,y))=\frac{1}{(1+x^2+y^2)^2}\left(2(1-x^2+y^2),-4xy,4x\right)} { . }

Karena $\sigma$ merupakan difeomorfisme, diferensialnya injektif; jadi medan ini tidak mempunyai titik nol pada $S^2\setminus\{N\}$. Selain itu,

\mathdisp {\left\|X(\sigma(x,y))\right\|=\frac{2}{1+x^2+y^2}} { , }

yang menuju $0$ ketika $x^2+y^2\to\infty$, yakni ketika $\sigma(x,y)\to N$. Oleh sebab itu, medan tersebut dapat diperluas secara kontinu dengan menetapkan

\mavergleichskettedisp
{\vergleichskette
{X(N)
}
{ =} {0
}
{ \in }{T_NS^2
}
{ } {
}
{ } {
}
}
{}{}{.}
Dalam realisasi $TS^2\subseteq S^2\times\R^3$, rumus norma di atas langsung membuktikan kontinuitas pada $N$. Dengan demikian, $X$ adalah medan vektor kontinu pada $S^2$ dan satu-satunya titik nolnya ialah $N$.

Catatan edisi: Solusi sumber menyatakan medan $e_1/(x^2+y^2)$ pada seluruh $\R^2$, padahal rumus itu tidak terdefinisi di titik asal, lalu menyimpulkan kontinuitas di kutub hanya dari komponen koordinat. Konstruksi di atas memakai medan konstan yang terdefinisi di seluruh bidang dan memeriksa norma dorong-majunya di dalam realisasi ambien, sehingga perluasan di kutub benar-benar kontinu.


\backmatter
\listoffigures
\chapter*{Atribusi dan Hak Media}
\addcontentsline{toc}{chapter}{Atribusi dan Hak Media}
\section*{Unit 1}
\addcontentsline{toc}{section}{Unit 1}
Empat berkas berikut digunakan dalam Unit 1. Setiap berkas tetap mengikuti lisensinya sendiri; tautan sumber dan lisensi dapat diklik.
\begin{description}
\item[\texttt{3d-function-6.svg}] MartinThoma (karya asli).
\href{https://commons.wikimedia.org/wiki/File:3d-function-6.svg}{Halaman sumber di Wikimedia Commons}; lisensi \href{https://creativecommons.org/licenses/by/3.0}{CC BY 3.0}.
\item[\texttt{Great circle passing through two points.svg}] HaEr48; karya turunan dari Polar angle to spherical side.svg oleh Episcophagus.
\href{https://commons.wikimedia.org/wiki/File:Great_circle_passing_through_two_points.svg}{Halaman sumber di Wikimedia Commons}; lisensi \href{https://creativecommons.org/licenses/by-sa/4.0}{CC BY-SA 4.0}.
\item[\texttt{2019-07-Helix.jpg}] Ag2gaeh (karya asli).
\href{https://commons.wikimedia.org/wiki/File:2019-07-Helix.jpg}{Halaman sumber di Wikimedia Commons}; lisensi \href{https://creativecommons.org/licenses/by-sa/4.0}{CC BY-SA 4.0}.
\item[\texttt{Planned flight map of the Oiseau Blanc.svg}] Pethrus; karya turunan dari BlankMap-World8.svg oleh AMK1211.
\href{https://commons.wikimedia.org/wiki/File:Planned_flight_map_of_the_Oiseau_Blanc.svg}{Halaman sumber di Wikimedia Commons}; lisensi \href{https://creativecommons.org/licenses/by-sa/3.0}{CC BY-SA 3.0}.
\end{description}
Identitas revisi, ukuran, SHA-1 Wikimedia, SHA-256, URL berkas asli, dan hash turunan cetak tercatat dalam manifest media yang disertakan.

\section*{Unit 2}
\addcontentsline{toc}{section}{Unit 2}
Dua berkas berikut digunakan dalam Unit 2. Setiap berkas tetap mengikuti status hak atau lisensinya sendiri; tautan sumber dapat diklik dan tautan lisensi dicantumkan jika tersedia.
\begin{description}
\item[]\texttt{Integral} \texttt{apl} \texttt{rot} \texttt{obsah1.svg}
Pajs; dialihkan dari Wikipedia bahasa Ceko ke Wikimedia Commons.
\href{https://commons.wikimedia.org/wiki/File:Integral_apl_rot_obsah1.svg}{Halaman sumber di Wikimedia Commons}; status hak: domain publik.
\item[]\texttt{Hyperboloid1.png}
Lars H. Rohwedder (RokerHRO; karya asli).
\href{https://commons.wikimedia.org/wiki/File:Hyperboloid1.png}{Halaman sumber di Wikimedia Commons}; status hak: domain publik.
\end{description}
Identitas revisi, ukuran, SHA-1 Wikimedia, SHA-256, URL berkas asli, dan hash turunan cetak tercatat dalam manifest media yang disertakan.

\section*{Unit 3}
\addcontentsline{toc}{section}{Unit 3}
Tiga berkas berikut digunakan dalam Unit 3. Setiap berkas tetap mengikuti status hak atau lisensinya sendiri; tautan sumber dapat diklik dan tautan lisensi dicantumkan jika tersedia.
\begin{description}
\item[]\texttt{Parabola} \texttt{circle.svg}
IkamusumeFan (karya sendiri).
\href{https://commons.wikimedia.org/wiki/File:Parabola_circle.svg}{Halaman sumber di Wikimedia Commons}; lisensi \href{https://creativecommons.org/licenses/by-sa/4.0}{CC BY-SA 4.0}.
\item[]\texttt{Euler} \texttt{spiral.svg}
AdiJapan (karya sendiri).
\href{https://commons.wikimedia.org/wiki/File:Euler_spiral.svg}{Halaman sumber di Wikimedia Commons}; lisensi \href{https://creativecommons.org/licenses/by-sa/3.0}{CC BY-SA 3.0}.
\item[]\texttt{Evolute-parab.svg}
Ag2gaeh (karya sendiri).
\href{https://commons.wikimedia.org/wiki/File:Evolute-parab.svg}{Halaman sumber di Wikimedia Commons}; lisensi \href{https://creativecommons.org/licenses/by-sa/4.0}{CC BY-SA 4.0}.
\end{description}
Identitas revisi, ukuran, SHA-1 Wikimedia, SHA-256, URL berkas asli, dan hash turunan cetak tercatat dalam manifest media yang disertakan.

\section*{Unit 4}
\addcontentsline{toc}{section}{Unit 4}
Unit 4 tidak menggunakan berkas media. Penutupan nol-aset ini diverifikasi terhadap kuliah, lembar kerja, dan semua solusi yang disediakan oleh sumber.

\section*{Unit 5}
\addcontentsline{toc}{section}{Unit 5}
Satu berkas berikut digunakan dalam Unit 5. Setiap berkas tetap mengikuti status hak atau lisensinya sendiri; tautan sumber dapat diklik dan tautan lisensi dicantumkan jika tersedia.
\begin{description}
\item[]\texttt{Minimal} \texttt{surface} \texttt{curvature} \texttt{planes-de.svg}
Eric Gaba (Sting); berdasarkan sebuah gambar dalam buku.
\href{https://commons.wikimedia.org/wiki/File:Minimal_surface_curvature_planes-de.svg}{Halaman sumber di Wikimedia Commons}; lisensi \href{https://creativecommons.org/licenses/by-sa/3.0/}{CC BY-SA 3.0}.
\end{description}
Identitas revisi, ukuran, SHA-1 Wikimedia, SHA-256, URL berkas asli, dan hash turunan cetak tercatat dalam manifest media yang disertakan.

\section*{Unit 6}
\addcontentsline{toc}{section}{Unit 6}
Satu berkas berikut digunakan dalam Unit 6. Setiap berkas tetap mengikuti status hak atau lisensinya sendiri; tautan sumber dapat diklik dan tautan lisensi dicantumkan jika tersedia.
\begin{description}
\item[]\texttt{Parallel} \texttt{transport} \texttt{sphere2.svg}
Silly rabbit di Wikipedia bahasa Inggris (karya sendiri).
\href{https://commons.wikimedia.org/wiki/File:Parallel_transport_sphere2.svg}{Halaman sumber di Wikimedia Commons}; lisensi \href{https://creativecommons.org/licenses/by-sa/3.0/}{CC BY-SA 3.0}.
\end{description}
Identitas revisi, ukuran, SHA-1 Wikimedia, SHA-256, URL berkas asli, dan hash turunan cetak tercatat dalam manifest media yang disertakan.

\section*{Unit 7}
\addcontentsline{toc}{section}{Unit 7}
Tiga berkas berikut digunakan dalam Unit 7. Setiap berkas tetap mengikuti status hak atau lisensinya sendiri; tautan sumber dapat diklik dan tautan lisensi dicantumkan jika tersedia.
\begin{description}
\item[]\texttt{Stereographic} \texttt{projection} \texttt{in} \texttt{3D.png}
Mark.Howison (karya sendiri; domain publik).
\href{https://commons.wikimedia.org/wiki/File:Stereographic_projection_in_3D.png}{Halaman sumber di Wikimedia Commons}; status hak: Public domain.
\item[]\texttt{Manifold} \texttt{zahyou3.png}
132ninme di Wikipedia bahasa Jepang (karya sendiri; nama Unicode persis dicatat dalam manifes).
\href{https://commons.wikimedia.org/wiki/File:Manifold_zahyou3.png}{Halaman sumber di Wikimedia Commons}; lisensi \href{https://creativecommons.org/licenses/by-sa/3.0/}{CC BY-SA 3.0}.
\item[]\texttt{Circle} \texttt{-} \texttt{black} \texttt{simple.svg}
Dakdada (karya sendiri; domain publik).
\href{https://commons.wikimedia.org/wiki/File:Circle_-_black_simple.svg}{Halaman sumber di Wikimedia Commons}; status hak: Public domain.
\end{description}
Identitas revisi, ukuran, SHA-1 Wikimedia, SHA-256, URL berkas asli, dan hash turunan cetak tercatat dalam manifest media yang disertakan.

\section*{Unit 8}
\addcontentsline{toc}{section}{Unit 8}
Unit 8 tidak menggunakan berkas media. Penutupan nol-aset ini diverifikasi terhadap kuliah, lembar kerja, dan semua solusi yang disediakan oleh sumber.

\section*{Unit 9}
\addcontentsline{toc}{section}{Unit 9}
Satu berkas berikut digunakan dalam Unit 9. Setiap berkas tetap mengikuti status hak atau lisensinya sendiri; tautan sumber dapat diklik dan tautan lisensi dicantumkan jika tersedia.
\begin{description}
\item[]\texttt{Tangentialvektor.svg}
McSush (karya turunan dari Tangentialvektor.png oleh TN; domain publik).
\href{https://commons.wikimedia.org/wiki/File:Tangentialvektor.svg}{Halaman sumber di Wikimedia Commons}; status hak: domain publik.
\end{description}
Identitas revisi, ukuran, SHA-1 Wikimedia, SHA-256, URL berkas asli, dan hash turunan cetak tercatat dalam manifest media yang disertakan.

\section*{Unit 10}
\addcontentsline{toc}{section}{Unit 10}
Dua berkas berikut digunakan dalam Unit 10. Setiap berkas tetap mengikuti status hak atau lisensinya sendiri; tautan sumber dapat diklik dan tautan lisensi dicantumkan jika tersedia.
\begin{description}
\item[]\texttt{Tangent} \texttt{bundle.svg}
Oleg Alexandrov (karya sendiri dengan Matlab dan Inkscape; domain publik).
\href{https://commons.wikimedia.org/wiki/File:Tangent_bundle.svg}{Halaman sumber di Wikimedia Commons}; status hak: domain publik.
\item[]\texttt{Torus} \texttt{vectors} \texttt{oblique.jpg}
RokerHRO (karya sendiri).
\href{https://commons.wikimedia.org/wiki/File:Torus_vectors_oblique.jpg}{Halaman sumber di Wikimedia Commons}; lisensi \href{https://creativecommons.org/licenses/by-sa/3.0/}{CC BY-SA 3.0}.
\end{description}
Identitas revisi, ukuran, SHA-1 Wikimedia, SHA-256, URL berkas asli, dan hash turunan cetak tercatat dalam manifest media yang disertakan.

\chapter*{Lisensi}
\Lizenztext
\end{document}
